\chapter{Quasi-lisse vertex algebras and the Arakawa--Kawasetsu theorem}\label{ch:L16}
\begingroup\setlength{\emergencystretch}{3em}\sloppy

\providecommand{\lsixteenvac}{\mathbf{1}}
\providecommand{\lsixteenHP}{\mathrm{HP}_0}
\providecommand{\lsixteensl}{\mathfrak{sl}_2}
\providecommand{\lsixteenCas}{\mathrm{Cas}}
\providecommand{\lsixteenNil}{\mathcal{N}}
\providecommand{\lsixteenR}{R}
\providecommand{\lsixteenVA}{V_{\mathcal A}}
\providecommand{\lsixteenOq}{O_q(V)}

This chapter documents the release of the LeanVOA library named in the front matter; it was first written against an earlier state of the library (``the snapshot'' below), and checks that were
run at an earlier state of the library and not re-run for this edition are marked as such. Every Lean quotation below is taken from the source files of the release
and is machine-checked by \texttt{scripts/textbook/sync.py} to occur verbatim (up to whitespace and
comments) in the source. The chapter covers the thirteen files assigned to it in
\texttt{Notes/Textbook/assignment.txt}, about seven thousand three hundred lines of Lean:
\begin{itemize}
\item \leanpath{VOA/\allowbreak Modules/\allowbreak PoissonTrace.lean},
  \leanpath{VOA/\allowbreak Modules/\allowbreak PoissonQuotient.lean},
  \leanpath{VOA/\allowbreak Modules/\allowbreak QuasiLisseFinite.lean} --- the Poisson-trace side of the
  theory and the finiteness half of Arakawa--Kawasetsu;
\item \leanpath{VOA/\allowbreak Modularity/\allowbreak AKGrading.lean},
  \leanpath{AKModularODE.lean}, \leanpath{AKUniformODE.lean}, \leanpath{AKMLDE.lean},
  \leanpath{AKZhuModularity.lean}, \leanpath{AKSl2Payoff.lean}, \leanpath{AKSl2Instance.lean},
  \leanpath{AKGradingInstance.lean}, \leanpath{AKGradingTotal.lean}, \leanpath{AKMLDEUnconditional.lean}
  (all in \leanpath{VOA/\allowbreak Modularity/}) --- the differential-equation and modularity half;
\item \leanpath{VOA/\allowbreak Examples/\allowbreak AffineSl2QuasiLisse.lean},
  \leanpath{AffineSl2Presentation.lean}, \leanpath{AffineSl2NonLisse.lean} (in
  \leanpath{VOA/\allowbreak Examples/}) --- the example $L_{-4/3}(\lsixteensl)$.
\end{itemize}

We rely on earlier chapters for the following, and do not repeat them: Lean idioms
(\lean{Submodule}, quotients \lean{⧸}, \lean{FiniteDimensional}, \lean{∃}/\lean{∀}, coercions) in
Chapter~\ref{ch:01}; the $C_2$-algebra $R_V = V/C_2(V)$ (the library's \lean{A₂ W}), its Poisson
structure, $C_2$-cofiniteness and finite strong generation in Chapter~\ref{ch:L12}
(Section~\ref{l12:sec:c2}); filtrations, symbols and associated graded objects in
Chapter~\ref{ch:L11}; modules, characters and conformal weights in Chapter~\ref{ch:L12}; the
universal affine vertex algebra $V^k(\lsixteensl)$, its singular vectors and symbols in
Chapter~\ref{ch:L08}; the simple quotient $L_{-4/3}(\lsixteensl)$ and Li's invariant-form radical in
Chapter~\ref{ch:L09}; the modular-forms toolkit (formal $q$-series, Eisenstein series, the Serre
derivative, the ring $\C[E_4,E_6]$) in Chapter~\ref{ch:L14}; and Zhu's modularity theorem, square
bracket modes, genus-one conformal blocks and the growth data in Chapter~\ref{ch:L15}.

\paragraph{What this chapter is about, in one paragraph.} Zhu's theorem says that if a vertex
operator algebra $V$ is $C_2$-cofinite (``lisse''), the characters of its modules are modular.
Arakawa and Kawasetsu observed that Zhu's argument uses $C_2$-cofiniteness only to make a certain
module over the ring of modular forms finitely generated, and that a much weaker hypothesis does
the same job: \emph{finiteness of the zeroth Poisson homology} $\lsixteenHP(R_V) = R_V/\{R_V,R_V\}$.
By a theorem of Etingof and Schedler this holds whenever the \emph{associated variety}
$X_V = \operatorname{Specm}R_V$ has finitely many symplectic leaves (``$V$ is quasi-lisse''). Their
conclusion is that the character of every ordinary module of a quasi-lisse $V$ satisfies a
\emph{modular linear differential equation} (MLDE). The motivating examples are admissible affine
vertex algebras such as $L_{-4/3}(\lsixteensl)$, and the chiral algebras of four-dimensional
$\mathcal N=2$ superconformal field theories, whose vacuum characters are Schur indices.
\textbf{The library does not prove this theorem for any vertex algebra} (finding~3 below): this chapter explains it, and
then records exactly which parts of it are formalised, and under which hypotheses.

\paragraph{The findings of this chapter, in brief.} Read these before anything else; each is
explained where it occurs.
\begin{enumerate}
\item \textbf{Quasi-lisseness is not formalised.} No associated variety, no symplectic leaf and no
  part of Etingof--Schedler exists in the library. Every theorem is stated under the hypothesis
  \lean{EtingofFinite W}, which \emph{is} the finiteness of $\lsixteenHP(R_V)$, i.e.\ the
  \emph{conclusion} of Etingof--Schedler. So ``quasi-lisse'' in the file names means ``Etingof-finite''
  in the statements. This is honest (every docstring says it) and mathematically it makes the
  theorems \emph{stronger} (they assume less), but no theorem mentions symplectic leaves.
  \textbf{What is formalised instead} is the following, and nothing more: for
  $L_{-4/3}(\lsixteensl)$, (a) it is not $C_2$-cofinite and its zeroth Poisson homology $\lsixteenHP(R_V)$ is
  finite-dimensional (``Etingof-finite''); (b) its character $q$-series converges on the upper half plane; (c) its
  space of genus-one conformal blocks is finite-dimensional. And for every VOA, $C_2$-cofinite $\Rightarrow$
  Etingof-finite. Wherever this chapter says ``quasi-lisse'' about a Lean statement it means ``Etingof-finite''; that a
  quasi-lisse VOA is Etingof-finite is the Etingof--Schedler theorem, \emph{known on paper, not formalised}.
\item \textbf{The hypothesis is discharged for a genuinely non-lisse example.} For
  $L_{-4/3}(\lsixteensl)$, the simple affine vertex algebra at the admissible level $-4/3$, the
  library proves both \lean{EtingofFinite} (\lean{VOA.Sl2Pres.}\allowbreak\lean{etingofFinite_}\allowbreak\lean{simpleSl2QL'}) and
  $\neg$\lean{C2Cofinite} (\lean{VOA.Sl2NonLisse.}\allowbreak\lean{not_C2Cofinite_}\allowbreak\lean{simpleSl2QL}), unconditionally
  (\texttt{\#print axioms}: Lean's three standard axioms only). So the strictness of ``$C_2$-cofinite $\Rightarrow$ Etingof-finite'' is a
  theorem. \textbf{Namespace warning:} the unconditional non-lisseness theorem is
  \lean{VOA.Sl2NonLisse.}\allowbreak\lean{not_C2Cofinite_}\allowbreak\lean{simpleSl2QL} (file \leanpath{AffineSl2NonLisse.lean}). A
  second declaration with the same short name, \lean{VOA.QuasiLisseSl2.}\allowbreak\lean{not_C2Cofinite_}\allowbreak\lean{simpleSl2QL}
  (file \leanpath{AKSl2Payoff.lean}), \emph{formerly} took a presentation \lean{S} and an exactness hypothesis
  \lean{hex}. After this chapter's finding, a later revision of the library restated it without
  hypotheses, as the one-line proof \lean{:= Sl2NonLisse.not_C2Cofinite_simpleSl2QL}; the two twins now have
  \emph{identical} types, and so do the twins \lean{QuasiLisseSl2.}\allowbreak\lean{etingofFinite_}\allowbreak\lean{simpleSl2QL} and
  \lean{Sl2Pres.}\allowbreak\lean{etingofFinite_}\allowbreak\lean{simpleSl2QL'}. The conditional forms survive only as
  \lean{Sl2Pres.}\allowbreak\lean{not_C2Cofinite_}\allowbreak\lean{simpleSl2QL_}\allowbreak\lean{of_exact} and
  \lean{Sl2Pres.}\allowbreak\lean{quasiLisseWitness_}\allowbreak\lean{simpleSl2QL_}\allowbreak\lean{of_exact}, whose hypothesis
  \lean{hex : Sl2ExactSymbol} is unproved and redundant. Throughout this chapter we always write the namespace.
\item \textbf{What ``modularity'' means in the proved theorems. Arakawa--Kawasetsu Theorem~5.1 (an analytic MLDE
  for all characters of a quasi-lisse VOA) is proved in the form: for every Etingof-finite VOA, one monic MLDE of
  unbounded order for all abstract analytic genus-one blocks (Theorem~\ref{l16:thm:mlde-uncond}); quasi-lisseness is not
  formalised.} Its one concrete instance is $L_{-4/3}(\lsixteensl)$ at the vacuum, of unspecified order. The only MLDEs of
  \emph{specified order} the library proves for specific VOAs are the formal, coefficientwise second-order equations for the \emph{vacuum} characters of the seven Deligne-series
  VOAs, obtained by Zhu's recursion in Chapter~\ref{ch:L25}. At $L_{-4/3}(\lsixteensl)$ the
  library proves, unconditionally: (a) the character satisfies a monic linear ODE in $q\,d/dq$ whose
  coefficients are power series with subexponentially growing coefficients
  (\lean{character_}\allowbreak\lean{formalODE_}\allowbreak\lean{simpleSl2QL_}\allowbreak\lean{uncond}); (b) the character $q$-series converges on the
  upper half plane; (c) the space of genus-one conformal blocks is finite-dimensional (Zhu's
  Theorem~5.3.2). \textbf{Statement (a) is not an MLDE and is not AK Theorem~5.1, although the source calls
  it ``AK equation (2), formal form'': by Proposition~\ref{l16:prop:weak} (on paper) every convergent
  $q$-series with nonzero leading term and zeros of bounded multiplicity in the unit disc satisfies such an
  equation, so (a) is near-vacuous --- it carries almost no information beyond the convergence (b).} A
  \emph{genuine} MLDE (iterated Serre derivatives, coefficients modular forms of the forced weights,
  Theorem~\ref{l16:thm:mlde}) is proved \textbf{conditionally only on \lean{EtingofFinite W}}. Its two former
  side inputs are now theorems: the Arakawa--Kawasetsu weight grading \lean{AKWeightGrading W} holds for every VOA of the
  library (CFT type), \lean{akWeightGrading W} (Definition~\ref{l16:def:akgrading}), and the Ramanujan/Serre closure
  \lean{E46SerreClosed} is \lean{e46SerreClosed} (Chapter~\ref{ch:L14}). The resulting statement is
  \lean{IsSqPrimary.}\allowbreak\lean{exists_}\allowbreak\lean{monicMLDEOn_}\allowbreak\lean{all_analyticBlocks_}\allowbreak\lean{of_etingofFinite_uncond}
  (conditional on \lean{EtingofFinite} despite the suffix), and at $L_{-4/3}(\lsixteensl)$ it gives the unconditional
  \lean{exists_}\allowbreak\lean{monicMLDEOn_}\allowbreak\lean{all_analyticBlocks_}\allowbreak\lean{simpleSl2QL}: one monic MLDE of
  some unspecified order $s\ge1$ satisfied by the vacuum one-point function of every analytic genus-one block (AK's order~2 is
  not proved). Both are \texttt{\#print axioms}-clean (recorded by WS-RELEASE2); the aggregate build at the release commit is
  pending. Zhu-style
  $SL_2(\Z)$-invariance of the span of characters under \lean{EtingofFinite} needs two more
  hypotheses (\lean{hfib}, \lean{hrank}), discharged only for the trivial VOA; and at the flagship
  $L_{-4/3}(\lsixteensl)$ \lean{hrank} is \emph{expected to be false} for every family containing the adjoint module
  (a paper argument of the release audit, not machine-checked; Section~\ref{l16:sec:zhu}). So the library contains
  \emph{no} modularity theorem for any non-lisse VOA.
\item \textbf{Lean defects found} (listed in Section~\ref{l16:sec:defects}): the hypothesis
  \lean{hex : Sl2ExactSymbol} of several theorems was redundant (now removed from the flagship ODE theorems); a dozen
  docstrings said that non-lisseness of $L_{-4/3}(\lsixteensl)$ is unproved or that only the
  trivial VOA is Etingof-finite (the declaration docstrings have since been corrected); two theorems share the short name
  \lean{not_C2Cofinite_}\allowbreak\lean{simpleSl2QL} (namespaces \lean{VOA.Sl2NonLisse} and \lean{VOA.QuasiLisseSl2};
  formerly with different hypotheses, now with identical types); the headline docstring of
  \lean{Sl2Pres.}\allowbreak\lean{character_}\allowbreak\lean{formalODE_}\allowbreak\lean{simpleSl2QL_}\allowbreak\lean{uncond}
  (\leanpath{AKSl2Instance.lean}, ll.~97--105 at the release audit) calls the near-vacuous formal equation ``THE
  HEADLINE'' and ``the Arakawa--Kawasetsu character equation''; and the citation attached to
  \lean{Sl2ExactSymbol} (``Arakawa's associated-variety theorem'') names a strictly weaker result.
\item \textbf{What may be said about this chapter, in one sentence} (independent release audit, 2026-10-27,
  \texttt{Notes/\allowbreak Textbook/\allowbreak referee/\allowbreak RELEASE\_AUDIT\_L16.md}): \emph{Lean proves that
  $L_{-4/3}(\lsixteensl)$ is not $C_2$-cofinite but has finite-dimensional zeroth Poisson homology, and that for every
  VOA with this finiteness property characters of ordinary modules converge on the upper half plane and the space of
  genus-one conformal blocks is finite-dimensional; the Arakawa--Kawasetsu MLDE is formalised conditionally only on
  \lean{EtingofFinite} (its grading and Serre-closure inputs are theorems), with an unbounded order, for Zhu's abstract
  analytic one-point blocks, and holds unconditionally for $L_{-4/3}(\lsixteensl)$ at the vacuum, order unspecified.}
  (The audit's sentence, written before \lean{akWeightGrading} landed, said the MLDE was instantiated for no VOA; that is
  superseded.) It would be false to say that ``the Arakawa--Kawasetsu theorem for quasi-lisse VOAs'' is formalised:
  quasi-lisseness is not formalised, and \lean{EtingofFinite} stands in for it.
\end{enumerate}

%======================================================================
\section{The mathematics: lisse, quasi-lisse, and the Arakawa--Kawasetsu theorem}\label{l16:sec:math}
%======================================================================

\subsection{The associated variety}

Let $V$ be a vertex operator algebra; in the library every VOA is of CFT type
($V=\bigoplus_{n\ge0}V_n$, $V_0=\C\lsixteenvac$, $\dim V_n<\infty$; Chapter~\ref{ch:L12},
Section~\ref{l12:sec:recall}). Zhu's subspace $C_2(V)$ is the span of all $u_{(-2)}v$, and
$R_V = V/C_2(V)$. Chapter~\ref{ch:L12} (Theorem~\ref{l12:thm:poisson}) shows that $R_V$ is a
commutative, associative, unital algebra under $[u]\cdot[v] = [u_{(-1)}v]$, and a Poisson algebra under
$\{[u],[v]\} = [u_{(0)}v]$: the bracket is a Lie bracket and a derivation of the product in each
argument. If $V$ is strongly generated by fields $a^1,\dots,a^r$ (every state is a linear combination of
monomials $a^{i_1}_{(-n_1)}\cdots a^{i_r}_{(-n_r)}\lsixteenvac$ with $n_j\ge1$, i.e.\ normally ordered products of the
fields $a^i$ and their derivatives $\partial^{n-1}a^i\leftrightarrow a^i_{(-n)}\lsixteenvac$), then $R_V$ is generated as a commutative
algebra by the classes $[a^i]$, so it is a finitely generated commutative $\C$-algebra: the coordinate ring
of an affine scheme. The \emph{associated variety} of Arakawa is the reduced affine variety
\begin{equation*}
X_V = \operatorname{Specm} R_V \subseteq \C^r ,
\end{equation*}
the set of maximal ideals of $R_V$ (equivalently, of points of the reduced affine scheme), i.e.\ the common zero set of
the relations among the generators. It inherits a Poisson structure from $R_V$. \emph{None of this geometry is formalised}:
the library works with the algebra $R_V$ and assumes finiteness of $\lsixteenHP(R_V)$ directly
(Definition~\ref{l16:def:etingof}).

\begin{example}[Three associated varieties]\label{l16:ex:three}
\begin{enumerate}
\item If $V$ is \emph{lisse} ($C_2$-cofinite, $\dim R_V<\infty$), $X_V$ is a finite set; since $R_V$ is
  graded by conformal weight with $(R_V)_0 = \C$, the only point is the origin: $X_V=\{0\}$.
\item For the Heisenberg (free boson) VOA $M(1)$ generated by one field $a$ of weight $1$,
  $R_{M(1)} = \C[x]$, $x=[a]$, with zero bracket (since $a_{(0)}a = 0$). So $X_{M(1)} = \C$, a line.
\item For the universal affine vertex algebra $V^k(\mathfrak g)$, PBW gives $R_{V^k(\mathfrak g)} = \C[\mathfrak g^*]$
  (symmetric algebra of $\mathfrak g$) with the Kirillov--Kostant bracket, so $X = \mathfrak g^*$. For the simple
  quotient $L_k(\lsixteensl)$ at an admissible level $k$, Arakawa proved $X_{L_k} = \lsixteenNil$, the
  nilpotent cone $\{h^2+4ef = 0\}\subset\lsixteensl^*\cong\C^3$, a two-dimensional cone isomorphic to
  $\C^2/\Z_2$.
\end{enumerate}
\end{example}

\subsection{Symplectic leaves, quasi-lisseness, Poisson homology}

The Poisson bracket of $R_V$ makes $X_V$ a Poisson variety. Each function $g$ generates a Hamiltonian vector
field $\{g,\cdot\}$; following these vector fields partitions (the smooth locus of each stratum of) $X_V$ into
\emph{symplectic leaves}, the maximal connected immersed submanifolds on which the Poisson structure is
non-degenerate. For the nilpotent cone there are exactly two leaves: the origin and the regular nilpotent orbit
$\lsixteenNil\setminus\{0\}$. For the line with zero bracket every point is its own leaf: infinitely many.

\begin{definition}[Quasi-lisse; Arakawa--Kawasetsu 2018]\label{l16:def:quasilisse}
A finitely strongly generated VOA $V$ is \emph{quasi-lisse} if $X_V$ has finitely many symplectic leaves.
\end{definition}

\begin{physics}
For the chiral algebra of a four-dimensional $\mathcal N=2$ theory, $X_V$ is (conjecturally) the Higgs branch, a
hyperk\"ahler cone stratified by finitely many symplectic leaves (the strata of partial Higgsing); quasi-lisse is the
two-dimensional shadow of ``the moduli space of vacua is a finite union of nice strata''.
\end{physics}

Lisse implies quasi-lisse (a point has one leaf). $M(1)$ is not quasi-lisse. Admissible $L_k(\lsixteensl)$ is
quasi-lisse but not lisse.

\begin{definition}[Zeroth Poisson homology, Poisson traces]\label{l16:def:hp0}
For a Poisson algebra $R$, let $\{R,R\}\subseteq R$ be the \emph{linear span} of all brackets $\{a,b\}$. The
zeroth Poisson homology is the vector space $\lsixteenHP(R) = R/\{R,R\}$. A \emph{Poisson trace} is a linear
functional $\varphi:R\to\C$ with $\varphi(\{a,b\})=0$ for all $a,b$; the Poisson traces form the dual space
$\lsixteenHP(R)^*$.
\end{definition}

\begin{physics}
A torus one-point function of a CFT is $\langle v\rangle_\tau = \operatorname{tr}_{\mathcal H}\,o(v)\,q^{L_0-c/24}$, the
expectation value of the zero mode of the field $v$ in the thermal ensemble at modular parameter $\tau$. It is a linear
functional of $v$ which, by cyclicity of the trace, vanishes on commutators $[o(a),o(v)] = o(a_{(0)}v)$, i.e.\ on the symmetry
variations of $v$ (this is condition (i) of a genus-one block, Definition~\ref{l16:def:formalblock}). In the semiclassical
limit $R_V$ replaces the operators, $u_{(0)}$ becomes the Hamiltonian vector field $\{u,\cdot\}$, and the trace becomes a
functional on phase-space functions invariant under all Hamiltonian flows: a Poisson trace. So Poisson traces are the
classical limits of torus one-point functions (``classical torus amplitudes'').
\end{physics}

\begin{theorem}[Etingof--Schedler 2010]\label{l16:thm:ES}
If $R$ is a finitely generated Poisson algebra whose variety $\operatorname{Specm}R$ has finitely many
symplectic leaves, then $\dim\lsixteenHP(R)<\infty$.
\end{theorem}

\begin{physics}
Heuristically: an invariant functional is determined by its behaviour on each leaf, where Hamiltonian flows act
transitively, so finitely many leaves allow only finitely many independent invariant functionals. (This heuristic is not the
proof, which is $D$-module theoretic.)
\end{physics}

This is a deep theorem (its proof uses $D$-modules: the Poisson traces are the solutions of a holonomic
$D$-module). It is \emph{not} formalised in the library and is never used as a black box either; instead its
conclusion is the hypothesis of every theorem. We shall call a VOA \emph{Etingof-finite} if
$\dim\lsixteenHP(R_V)<\infty$.

\subsection{The Arakawa--Kawasetsu theorem}

Recall (Chapter~\ref{ch:L14}) that the Serre derivative of weight $k$ is
$\vartheta_k f = \frac{1}{2\pi i}f' - \frac{k}{12}E_2 f$; it maps modular forms of weight $k$ to modular forms of
weight $k+2$. Its iterates are $\vartheta^{0}_{(k)} f = f$ and $\vartheta^{i+1}_{(k)}f =
\vartheta_{k+2i}\vartheta^{i}_{(k)}f$.

\begin{theorem}[Arakawa--Kawasetsu 2018, Theorem~5.1 and its proof]\label{l16:thm:AK}
Let $V$ be a quasi-lisse VOA (conformal, of CFT type). There are $s\ge1$ and modular forms $p_i$ of weight
$2(s-i)$ for $SL_2(\Z)$, i.e.\ $p_i\in\C[E_4,E_6]$ homogeneous, such that for every ordinary $V$-module $M$ the
normalised character $\chi_M(\tau) = \operatorname{tr}_M q^{L_0-c/24}$, $q = e^{2\pi i\tau}$, satisfies
\begin{equation}\label{l16:eq:mlde}
\vartheta^{s}_{(0)}\chi_M + \sum_{i=0}^{s-1}p_i(\tau)\,\vartheta^{i}_{(0)}\chi_M = 0 .
\end{equation}
Consequently $V$ has finitely many ordinary irreducible modules, and the characters $\chi_M$ lie in the
finite-dimensional, $SL_2(\Z)$-invariant solution space of~\eqref{l16:eq:mlde}.
\end{theorem}

Arakawa and Kawasetsu's proof runs as follows. (i) Etingof--Schedler gives
$\dim\lsixteenHP(R_V)<\infty$ (their Theorem~2.2). (ii) Over the ring $\mathcal A = \C[G_4,G_6]$ (or
$\C[G_2,G_4,G_6]$), the module $\lsixteenVA/\lsixteenOq$ of Zhu --- $\lsixteenVA = \mathcal A\otimes V$ modulo
the relations that all genus-one one-point functions satisfy --- is finitely generated (their
Proposition~5.2); in Zhu's proof this was where $C_2$-cofiniteness entered, and the new observation is that the
relations $a[0]v$, which Zhu never used, have leading term $a_{(0)}v$, the Poisson bracket. (iii) Hence the
tower $L[-2]^i\lsixteenvac$ satisfies a monic relation modulo $\lsixteenOq$ (Lemma~5.3). (iv) The
$\omega$-recursion of Zhu, $Z(L[-2]v) = q\frac{d}{dq}Z(v) + \sum_{m\ge0}G_{2m+2}Z(L[2m]v)$ for every torus one-point
function $Z$ (condition (iii) of Definition~\ref{l16:def:formalblock}), converts $L[-2]$ into a Serre derivative up to lower-order terms (Lemmas~5.4--5.6). (v)
Substituting gives~\eqref{l16:eq:mlde}.

\begin{physics}
\textbf{Associated varieties are Higgs branches.} Beem, Lemos, Liendo, Peelaers, Rastelli and van Rees (2015)
showed that every four-dimensional $\mathcal N=2$ superconformal field theory contains a protected sector,
the \emph{Schur operators}, whose operator product expansion on a plane is that of a two-dimensional chiral
algebra (a VOA) $V_{4d}$, with central charge $c_{2d} = -12\,c_{4d}$. Its vacuum character is the \emph{Schur
limit of the superconformal index}. Beem and Rastelli (2018) conjectured, and verified in many examples,
that the associated variety $X_{V_{4d}}$ is the \emph{Higgs branch} of the four-dimensional theory --- the
moduli space of vacua on which the scalars of the hypermultiplets get expectation values. Higgs branches are
hyperk\"ahler cones with finitely many symplectic leaves, so these chiral algebras are quasi-lisse but (whenever
the Higgs branch is not a point) not lisse. The Arakawa--Kawasetsu theorem then predicts that the Schur index
satisfies a modular linear differential equation, which Beem and Rastelli found and used. The example of this
chapter, $L_{-4/3}(\lsixteensl)$ with $c = -6$, is the chiral algebra of the Argyres--Douglas theory
$(A_1,D_3)\cong(A_1,A_3)$, whose $c_{4d} = 1/2$ and whose Higgs branch is $\C^2/\Z_2$, the nilpotent cone of
$\lsixteensl$. Such theories are non-rational in the two-dimensional sense (the chiral algebra has infinitely
many non-ordinary modules) yet have finitely many ordinary modules, whose characters lie in the finite-dimensional,
modular-invariant solution space of an MLDE. (The span of the ordinary characters alone need not be invariant; for
$L_{-4/3}(\lsixteensl)$ it is expected not to be, by the paper argument in the warning after
Theorem~\ref{l16:thm:zhuEt}.) Poisson traces are the classical shadow of torus one-point functions: a one-point function
$v\mapsto\operatorname{tr}_M o(v)q^{L_0-c/24}$ kills $a[0]v$, whose leading term is a Poisson bracket.
\end{physics}

\subsection{What the library formalises: a map}

\begin{center}\small
\begin{tabular}{p{0.38\textwidth}p{0.55\textwidth}}
\toprule
Mathematics & Library\\
\midrule
associated variety, symplectic leaves, quasi-lisse & \emph{not defined}\\
Etingof--Schedler (Thm.~\ref{l16:thm:ES}) & \emph{not proved}; its conclusion is the hypothesis \lean{EtingofFinite W}\\
$C_2$-cofinite $\Rightarrow$ Etingof-finite & \lean{EtingofFinite_}\allowbreak\lean{of_C2Cofinite}, unconditional\\
AK Prop.~5.2 (finite generation) & \lean{exists_}\allowbreak\lean{finite_}\allowbreak\lean{gen_sup_}\allowbreak\lean{blockRelSpanG_}\allowbreak\lean{of_etingofFinite} ($\C[G_2,G_4,G_6]$) and \lean{exists_}\allowbreak\lean{finite_}\allowbreak\lean{gen_sup_}\allowbreak\lean{blockRelSpanE_}\allowbreak\lean{of_etingofFinite} ($\C[E_4,E_6]$), from \lean{EtingofFinite}\\
AK Lemma 5.3, 5.5 & proved, from \lean{EtingofFinite}; Lemma 5.5 for square-bracket primary vectors\\
AK eq.~(2) in formal $q$-series form (\textbf{near-vacuous}, Prop.~\ref{l16:prop:weak}) & \lean{character_}\allowbreak\lean{formalODE_}\allowbreak\lean{of_etingofFinite}: $(q\frac{d}{dq})$-basis, subexponential coefficients; \emph{not an MLDE}\\
AK weight grading & \lean{akWeightGrading W : AKWeightGrading W}, for every VOA of the library (CFT type)\\
AK Thm.~5.1 (genuine MLDE) & \lean{IsSqPrimary.}\allowbreak\lean{exists_}\allowbreak\lean{monicMLDEOn_}\allowbreak\lean{all_analyticBlocks_}\allowbreak\lean{of_etingofFinite_uncond}, conditional only on \lean{EtingofFinite W} (grading and Serre closure discharged); order $s\ge1$ existential, unbounded; for Zhu's abstract analytic blocks\\
AK example $L_{-4/3}(\lsixteensl)$ & \lean{exists_}\allowbreak\lean{monicMLDEOn_}\allowbreak\lean{all_analyticBlocks_}\allowbreak\lean{simpleSl2QL}, unconditional, vacuum insertion, order unspecified (AK's order~2 \emph{not} proved)\\
Zhu modularity under Etingof-finiteness & \lean{exists_}\allowbreak\lean{zhuModularity_}\allowbreak\lean{of_sqAnalyticRank_}\allowbreak\lean{le_irr_}\allowbreak\lean{etingofFinite}, conditional on \lean{hfib}, \lean{hrank}; discharged only for the trivial VOA, and \lean{hrank} expected false at $L_{-4/3}(\lsixteensl)$ (paper argument)\\
$L_{-4/3}(\lsixteensl)$ Etingof-finite & \lean{Sl2Pres.}\allowbreak\lean{etingofFinite_}\allowbreak\lean{simpleSl2QL'}, unconditional\\
$L_{-4/3}(\lsixteensl)$ not lisse & \lean{Sl2NonLisse.}\allowbreak\lean{not_C2Cofinite_}\allowbreak\lean{simpleSl2QL}, unconditional\\
$X_{L_{-4/3}} = \lsixteenNil$ (Arakawa) & \emph{not proved}; a stronger scheme-theoretic form is the named \lean{Sl2ExactSymbol}, now unneeded\\
\bottomrule
\end{tabular}
\end{center}

%======================================================================
\section{The Poisson algebra as a $\C$-algebra, and quotients}\label{l16:sec:quot}
%======================================================================

Chapter~\ref{ch:L12} builds the commutative ring \lean{A₂CommRing W} and the Lie ring \lean{A₂LieRing W} on
\lean{A₂ W} $=R_V$, writes \lean{mk₂ W : 𝒱 →ₗ[ℂ] A₂ W} for the projection $u\mapsto[u]$, \lean{A₂mul} and
\lean{A₂bracket} for the descended bilinear maps, and proves the Leibniz rule. Two small pieces of
infrastructure are added in \leanpath{VOA/\allowbreak Modules/\allowbreak PoissonQuotient.lean}.

\begin{definition}[$R_V$ as a $\C$-algebra]\label{l16:def:algebra}
$R_V$ is a $\C$-algebra with structure map $c\mapsto c[\lsixteenvac]$.
\end{definition}
\noindent File \leanpath{VOA/\allowbreak Modules/\allowbreak PoissonQuotient.lean}.
\begin{leancode}
theorem A₂_smul_mul (W : VertexOperatorAlgebra 𝒱) (c : ℂ) (p q : A₂ W) :
    (c • p) * q = c • (p * q) := by

theorem A₂_mul_smul (W : VertexOperatorAlgebra 𝒱) (c : ℂ) (p q : A₂ W) :
    p * (c • q) = c • (p * q) :=

noncomputable instance A₂Algebra (W : VertexOperatorAlgebra 𝒱) : Algebra ℂ (A₂ W) :=
  Algebra.ofModule (A₂_smul_mul W) (A₂_mul_smul W)

theorem A₂_algebraMap_apply (W : VertexOperatorAlgebra 𝒱) (c : ℂ) :
    algebraMap ℂ (A₂ W) c = c • (1 : A₂ W) := rfl
\end{leancode}

\begin{dictionary}
\lean{c • p} is the scalar multiplication that \lean{A₂ W} inherits as a quotient vector space; \lean{p * q} is the
product of \lean{A₂CommRing}. \lean{Algebra.ofModule} is Mathlib's constructor that turns a ring which is
already a vector space, with the two interchange laws $(cp)q = c(pq) = p(cq)$, into a $\C$-algebra; it reuses
the existing scalar multiplication, so there is no second, possibly different, $\C$-action (no ``diamond''),
and \lean{A₂_algebraMap_apply} is true by \lean{rfl} (by unfolding). Same as the textbook. Why it matters:
with an \lean{Algebra} instance, Mathlib's \lean{MvPolynomial.aeval} defines algebra homomorphisms out of a
polynomial ring by choosing images of the variables, which is how the symbol map
$\C[e,h,f]\to R_V$ of Section~\ref{l16:sec:pres} is built.
\end{dictionary}

\textbf{Proof on paper.} The product is bilinear, hence $\C$-linear in each slot. \qed

\begin{physics}
Scalars multiply states; the classical limit $R_V$ of the chiral algebra is an algebra of functions on the
classical phase space, and constants $c\cdot 1$ are the constant functions (multiples of the identity
operator).
\end{physics}

\begin{status}
Proved; \texttt{\#print axioms} of the four declarations: \lean{[propext, Classical.}\allowbreak\lean{choice, Quot.}\allowbreak\lean{sound]} (the
file's own audit block, compiled in the aggregate build).
\end{status}

\begin{theorem}[$R$ is functorial for quotients]\label{l16:thm:quot}
Let $I\subseteq V$ be a graded ideal and $V/I$ the quotient VOA. The map $[u]\mapsto[u+I]$ is a well-defined
surjective homomorphism of $\C$-algebras $R_V\to R_{V/I}$ which also preserves the Poisson bracket.
\end{theorem}
\noindent File \leanpath{VOA/\allowbreak Modules/\allowbreak PoissonQuotient.lean}.
\begin{leancode}
noncomputable def A₂quot (hI : W.IsIdeal I) (hgr : IsGradedIdeal W I) :
    A₂ W →ₐ[ℂ] A₂ (quotientVOA W I hI hgr) :=

@[simp] theorem A₂quot_mk₂ (hI : W.IsIdeal I) (hgr : IsGradedIdeal W I) (u : 𝒱) :
    A₂quot hI hgr (mk₂ W u) = mk₂ (quotientVOA W I hI hgr) (I.mkQ u) := rfl

theorem A₂quot_surjective (hI : W.IsIdeal I) (hgr : IsGradedIdeal W I) :
    Function.Surjective (A₂quot hI hgr) := A₂quotₗ_surjective hI hgr

theorem A₂quot_bracket (hI : W.IsIdeal I) (hgr : IsGradedIdeal W I) (p q : A₂ W) :
    A₂quot hI hgr (A₂bracket W p q)
      = A₂bracket (quotientVOA W I hI hgr) (A₂quot hI hgr p) (A₂quot hI hgr q) :=
\end{leancode}

\begin{dictionary}
\lean{→ₐ[ℂ]} is a homomorphism of $\C$-algebras. \lean{quotientVOA W I hI hgr} is the quotient VOA of
Chapter~\ref{ch:L04} (Section~\ref{l04:sec:quot}); its two arguments \lean{hI} ($I$ is a two-sided ideal) and
\lean{hgr} ($I$ is the sum of its intersections with the weight spaces) are exactly what is needed for the
quotient to be a VOA. \lean{I.mkQ} is the quotient map $V\to V/I$. The four statements are: the map exists
and is an algebra map (multiplicative, unital, linear); it sends $[u]$ to $[u+I]$; it is surjective; it is a
Poisson map. Same as the textbook. Could it be vacuous? It is a construction, and the formula
\lean{A₂quot_mk₂} pins it down.
\end{dictionary}

\textbf{Proof on paper.} Modes of the quotient are computed on representatives, $[u]_{(n)}[v] = [u_{(n)}v]$; so
$C_2(V)$ maps into $C_2(V/I)$ and the composite $V\to V/I\to R_{V/I}$ factors through $R_V$. The same formula at
$n=-1$ and $n=0$ gives compatibility with product and bracket; surjectivity is inherited from $V\to V/I$. \qed

\begin{physics}
Setting null states to zero (passing to a quotient) can only add relations among the classical fields: the
classical phase space of the quotient is a closed Poisson subvariety, $X_{V/I}\subseteq X_V$. This is how the
nilpotent cone appears inside $\lsixteensl^*$ for $L_k(\lsixteensl)$.
\end{physics}

\begin{status}
Proved; axioms \lean{[propext, Classical.}\allowbreak\lean{choice, Quot.}\allowbreak\lean{sound]} (re-run for this chapter on \lean{VOA.QuotVOA.A₂quot}
with \texttt{lake env lean -DmaxSynthPendingDepth=3}).
\end{status}

%======================================================================
\section{Poisson traces and the Etingof--Schedler hypothesis}\label{l16:sec:traces}
%======================================================================

\begin{definition}[Bracket span, $\lsixteenHP$, Poisson traces]\label{l16:def:poissonSpan}
$\{R_V,R_V\}$ is the linear span of all $\{p,q\}$, $p,q\in R_V$; $\lsixteenHP(R_V) = R_V/\{R_V,R_V\}$; the Poisson
traces are the linear functionals on $R_V$ vanishing on $\{R_V,R_V\}$.
\end{definition}
\noindent File \leanpath{VOA/\allowbreak Modules/\allowbreak PoissonTrace.lean}.
\begin{leancode}
noncomputable def poissonSpan (W : VertexOperatorAlgebra 𝒱) : Submodule ℂ (A₂ W) :=
  bilinSpan (A₂bracket W)

theorem poissonSpan_def (W : VertexOperatorAlgebra 𝒱) :
    poissonSpan W = Submodule.span ℂ {x : A₂ W | ∃ p q : A₂ W, x = A₂bracket W p q} := rfl

noncomputable abbrev HP₀ (W : VertexOperatorAlgebra 𝒱) : Type u := A₂ W ⧸ poissonSpan W

noncomputable def poissonTraces (W : VertexOperatorAlgebra 𝒱) :
    Submodule ℂ (Module.Dual ℂ (A₂ W)) :=
  vanishingFunctionals ℂ (A₂ W) (poissonSpan W)

theorem mem_poissonTraces {W : VertexOperatorAlgebra 𝒱} {φ : Module.Dual ℂ (A₂ W)} :
    φ ∈ poissonTraces W ↔ ∀ p q : A₂ W, φ (A₂bracket W p q) = 0 :=
\end{leancode}

\begin{dictionary}
\lean{bilinSpan B} (file \leanpath{VOA/\allowbreak Modularity/\allowbreak TraceSpace.lean}, Chapter~\ref{ch:L15}) is the
span of the range of a bilinear map $B$; \lean{poissonSpan_def} unfolds it to ``the span of all
$\{p,q\}$'', by \lean{rfl}. \lean{Module.Dual ℂ X} is the space of linear functionals $X\to\C$;
\lean{vanishingFunctionals ℂ X S} is the subspace of those vanishing on $S$, so \lean{poissonTraces} is literally
$\{R_V,R_V\}^{\perp}$, and \lean{mem_poissonTraces} says so. The library also proves
\lean{poissonTracesEquiv : poissonTraces W ≃ₗ[ℂ] Module.}\allowbreak\lean{Dual ℂ (HP₀ W)} (traces = dual of $\lsixteenHP$) and
\lean{poissonSpan_eq_map} (the span downstairs is the image of the span of the representatives
$u_{(0)}v$ upstairs). Note that $\lsixteenHP$ is defined as a \emph{vector space}: $\{R,R\}$ is not an ideal, and no
ring structure is claimed. All the same as Definition~\ref{l16:def:hp0}.
\end{dictionary}

\begin{physics}
A Poisson trace is a ``classical torus one-point function'': a linear functional on classical observables
invariant under all Hamiltonian flows. On a symplectic leaf, invariance under all Hamiltonian flows forces a
functional to be (up to the measure) constant along the leaf, which is the heuristic reason why finitely many
leaves give finitely many independent traces.
\end{physics}

\begin{status}
Definitions and membership lemmas; axioms \lean{[propext, Classical.}\allowbreak\lean{choice, Quot.}\allowbreak\lean{sound]} (file docstring and
aggregate audit).
\end{status}

\begin{definition}[Etingof-finite]\label{l16:def:etingof}
$V$ is \emph{Etingof-finite} if $\dim_\C\lsixteenHP(R_V)<\infty$.
\end{definition}
\noindent File \leanpath{VOA/\allowbreak Modules/\allowbreak PoissonTrace.lean}.
\begin{leancode}
def EtingofFinite (W : VertexOperatorAlgebra 𝒱) : Prop := FiniteDimensional ℂ (HP₀ W)
\end{leancode}

\begin{dictionary}
\lean{FiniteDimensional ℂ (HP₀ W)}: the vector space $R_V/\{R_V,R_V\}$ is finite-dimensional. This is exactly
Definition~\ref{l16:def:etingof}. \textbf{Is it quasi-lisseness?} No. Quasi-lisseness
(Definition~\ref{l16:def:quasilisse}) implies it by Etingof--Schedler (Theorem~\ref{l16:thm:ES}); the converse is
not claimed anywhere and we do not know it to be true in general. Every theorem of this chapter assumes
\lean{EtingofFinite}, not quasi-lisseness, so each is (mathematically) a \emph{stronger} statement than the
literature's version with ``quasi-lisse'' --- provided its conclusion matches, which is the question for the
later sections. The docstring of \lean{EtingofFinite} states all of this in capitals; it is a deliberate
choice that removes all algebraic geometry from the formalisation. \textbf{Could it be vacuous?} It holds for
every $C_2$-cofinite VOA (next theorem) and for $L_{-4/3}(\lsixteensl)$ (Theorem~\ref{l16:thm:presL}); it fails
for the free boson on paper (Remark~\ref{l16:rem:boson}).
\end{dictionary}

\begin{warning}[title={``Quasi-lisse'' is a file name, not a definition}]
No declaration of the library defines quasi-lisseness, associated varieties or symplectic leaves (a search of
\leanpath{VOA/} for declarations whose names contain \texttt{QuasiLisse}, \texttt{AssociatedVariety},
\texttt{Leaf} or \texttt{Symplectic} finds none; the word ``symplectic'' occurs only in comments and in the
unrelated symplectic-fermion files of Chapter~\ref{ch:L17}). The
files \leanpath{QuasiLisseFinite.lean} and \leanpath{AffineSl2QuasiLisse.lean} prove statements about
\lean{EtingofFinite}. A referee should read ``quasi-lisse'' in a Lean name as ``Etingof-finite''.
\end{warning}

\begin{physics}
Physically, $\dim\lsixteenHP<\infty$ says that the chiral algebra has finitely many independent ``classical
torus amplitudes''. For a four-dimensional $\mathcal N=2$ theory it is the statement that the Higgs branch
chiral ring has finite-dimensional Poisson homology, true because the Higgs branch is a hyperk\"ahler cone with
finitely many leaves.
\end{physics}

\begin{status}
Definition. Satisfiable (Theorems~\ref{l16:thm:c2etingof} and~\ref{l16:thm:presL}).
\end{status}

\begin{theorem}[Lisse implies Etingof-finite; the trace space]\label{l16:thm:c2etingof}
(i) If $V$ is $C_2$-cofinite then $V$ is Etingof-finite. (ii) If $V$ is Etingof-finite, the space of Poisson
traces on $R_V$ is finite-dimensional. (iii) If the bracket of $R_V$ vanishes identically, Etingof-finite is
equivalent to $C_2$-cofinite.
\end{theorem}
\noindent File \leanpath{VOA/\allowbreak Modules/\allowbreak PoissonTrace.lean}.
\begin{leancode}
theorem EtingofFinite_of_C2Cofinite (W : VertexOperatorAlgebra 𝒱) (h : C2Cofinite W) :
    EtingofFinite W := by

theorem finiteDimensional_poissonTraces (W : VertexOperatorAlgebra 𝒱) (h : EtingofFinite W) :
    FiniteDimensional ℂ (poissonTraces W) := by

theorem EtingofFinite_iff_C2Cofinite_of_poissonSpan_eq_bot (W : VertexOperatorAlgebra 𝒱)
    (h : poissonSpan W = ⊥) : EtingofFinite W ↔ C2Cofinite W := by
\end{leancode}

\begin{dictionary}
\lean{C2Cofinite W} is \lean{FiniteDimensional ℂ (A₂ W)} (Chapter~\ref{ch:L12}, Definition~\ref{l12:def:c2}).
\lean{poissonSpan W = ⊥} says the bracket span is zero. Statements (i)--(iii) exactly. These are easy facts; their
role is sanity: (i) shows the hypothesis is not stronger than Zhu's, (iii) shows where it is not weaker.
\end{dictionary}

\textbf{Proof on paper.} (i) $\lsixteenHP(R_V)$ is a quotient of $R_V$. (ii) The traces are the dual of the
finite-dimensional space $\lsixteenHP(R_V)$. (iii) If $\{R_V,R_V\}=0$ then $\lsixteenHP(R_V)=R_V$. \qed

\begin{remark}[The free boson is not a witness]\label{l16:rem:boson}
For the free boson $M(1)$, $R_{M(1)}\cong\C[x]$ with zero bracket, so by (iii) $M(1)$ is Etingof-finite iff it is
$C_2$-cofinite, and it is not (Chapter~\ref{ch:L12}, Example~\ref{l12:ex:c2}). The library does \emph{not} prove
\lean{poissonSpan (heisenbergVOA k hk) = ⊥} (the closing note of \leanpath{PoissonTrace.lean} explains that the
identification of $R_{M(1)}$ with $\C[x]$ \emph{as a ring} is missing), so ``$M(1)$ is not Etingof-finite'' is an
argument on paper only. Geometrically: $X_{M(1)}$ is a line with zero Poisson structure, with infinitely many
leaves.
\end{remark}

\begin{physics}
A lisse chiral algebra has a zero-dimensional phase space, so trivially finitely many classical amplitudes. A
free boson has a continuous family of vacua (the expectation value of $a$, a flat direction or ``modulus''), and
the classical traces $\varphi(x^n)$ are independent: infinitely many. Its characters $1/\eta$ are modular for a
different reason (it is a free theory), not because of Etingof-finiteness.
\end{physics}

\begin{status}
Proved, unconditional; axioms of \lean{VOA.EtingofFinite_}\allowbreak\lean{of_C2Cofinite} and
\lean{VOA.finiteDimensional_}\allowbreak\lean{poissonTraces}: \lean{[propext, Classical.}\allowbreak\lean{choice, Quot.}\allowbreak\lean{sound]} (re-run for this
chapter).
\end{status}

%======================================================================
\section{The finiteness half: Arakawa--Kawasetsu Proposition 5.2}\label{l16:sec:prop52}
%======================================================================

\subsection{The augmented relation space}

The first step is to move the Etingof--Schedler hypothesis from $R_V$ up to $V$.

\begin{definition}[$O_2^P(V)$]\label{l16:def:O2P}
$O_2^P(V) = C_2(V) + \operatorname{span}\{u_{(0)}v : u,v\in V\}\subseteq V$.
\end{definition}
\noindent File \leanpath{VOA/\allowbreak Modules/\allowbreak QuasiLisseFinite.lean}.
\begin{leancode}
noncomputable def poissonUp (W : VertexOperatorAlgebra 𝒱) : Submodule ℂ 𝒱 :=
  Submodule.span ℂ {x : 𝒱 | ∃ u v : 𝒱, x = poissonBracket W u v}

noncomputable def O₂P (W : VertexOperatorAlgebra 𝒱) : Submodule ℂ 𝒱 := O₂ W ⊔ poissonUp W
\end{leancode}

\begin{dictionary}
\lean{poissonBracket W u v} is \lean{W.smul u 0 v} $=u_{(0)}v$ (Chapter~\ref{ch:L12}); \lean{⊔} of subspaces is
their sum. Same as Definition~\ref{l16:def:O2P}.
\end{dictionary}

\begin{physics}
$O_2^P(V)$ consists of states that are classically derivatives or symmetry variations; both are invisible to torus
one-point functions at leading order.
\end{physics}

\begin{theorem}[AK Theorem 2.2 read upstairs]\label{l16:thm:O2P}
$V$ is Etingof-finite iff $V/O_2^P(V)$ is finite-dimensional. Moreover $O_2^P(V)$ is a graded subspace, and
for $a\in V_p$, $b\in V_q$ one has $a_{(0)}b\in V_{p+q-1}$.
\end{theorem}
\noindent File \leanpath{VOA/\allowbreak Modules/\allowbreak QuasiLisseFinite.lean}.
\begin{leancode}
theorem map_mkQ_O₂P (W : VertexOperatorAlgebra 𝒱) :
    Submodule.map (O₂ W).mkQ (O₂P W) = poissonSpan W := by

theorem etingofFinite_iff_quot_O₂P (W : VertexOperatorAlgebra 𝒱) :
    EtingofFinite W ↔ FiniteDimensional ℂ (𝒱 ⧸ O₂P W) := by
\end{leancode}

\begin{dictionary}
\lean{Submodule.}\allowbreak\lean{map (O₂ W).}\allowbreak\lean{mkQ (O₂P W)} is the image of $O_2^P(V)$ in $R_V$; \lean{map_mkQ_O₂P} says it is
$\{R_V,R_V\}$. Then the third isomorphism theorem $(V/C_2)/(O_2^P/C_2)\cong V/O_2^P$ gives the equivalence.
The gradedness statements are \lean{O₂P_eq_span_O₂PhomGen} and \lean{piHom_maps_O₂P} (projections to weight
spaces preserve $O_2^P$), and \lean{brk_mem_grading} (the bracket lowers total weight by one). Same as AK's
Theorem 2.2 combined with their remark that $\lsixteenHP(R_V)$ is graded.
\end{dictionary}

\textbf{Proof on paper.} The image of $C_2(V)$ in $R_V$ is $0$ and the image of $u_{(0)}v$ is $\{[u],[v]\}$,
so the image of $O_2^P$ is $\{R_V,R_V\}$; then apply the isomorphism theorem. For weights: $a_{(n)}b$ has
weight $p+q-n-1$, so $a_{(0)}b\in V_{p+q-1}$ and $a_{(-2)}b\in V_{p+q+1}$; both generating families are made of
homogeneous vectors, hence the spans are graded. \qed

\begin{physics}
$V/O_2^P(V)$ is the space of states modulo ``derivatives'' (the $u_{(-2)}v$, which involve $\partial u$) and
modulo ``charge commutators'' $u_{(0)}v$ (the infinitesimal symmetry transformation of $v$ generated by the charge
$\oint u$). A one-point function on the torus kills both kinds of state up to modular corrections, so its value is
determined by finitely many numbers when this quotient is finite.
\end{physics}

\begin{status}
Proved, unconditional; \lean{VOA.etingofFinite_}\allowbreak\lean{iff_quot_}\allowbreak\lean{O₂P}: \lean{[propext, Classical.}\allowbreak\lean{choice, Quot.}\allowbreak\lean{sound]}
(re-run for this chapter).
\end{status}

\subsection{Mode-zero triangularity}

Zhu's genus-one theory uses the \emph{square-bracket modes} $a[m]$ of a homogeneous vector $a$ of weight
$k$ (Chapter~\ref{ch:L15}):
\begin{equation*}
a[m] = \Res_z\,Y(a,z)\,(1+z)^{k-1}\bigl(\log(1+z)\bigr)^m = \sum_{j\ge0}c^{(k,m)}_j\,a_{(m+j)} ,
\end{equation*}
where $c^{(k,m)}_j$ is the coefficient of $z^j$ in $(1+z)^{k-1}\bigl(\log(1+z)/z\bigr)^m$. In Lean this is
\lean{vSqZ W a k m}, with coefficients \lean{sqBracketCoeffZ k m i} (file
\leanpath{VOA/\allowbreak Modularity/\allowbreak OnePointFunctions.lean}, Chapter~\ref{ch:L15}).
The leading coefficient is $c^{(k,m)}_0 = 1$. At $m=0$:
$a[0]b = a_{(0)}b + \sum_{i\ge1}c^{(k,0)}_i a_{(i)}b$.

\begin{theorem}[Mode-$0$ triangularity]\label{l16:thm:triang}
For $a\in V_p$ and $b\in V_q$, $a[0]b - a_{(0)}b$ lies in $\bigoplus_{m<p+q-1}V_m$.
\end{theorem}
\noindent File \leanpath{VOA/\allowbreak Modules/\allowbreak QuasiLisseFinite.lean}.
\begin{leancode}
noncomputable def sqDefect0 (W : VertexOperatorAlgebra 𝒱) (a : 𝒱) (p : ℕ) (b : 𝒱) : 𝒱 :=
  ∑ᶠ i : ℤ, (if i = 0 then (0 : ℂ) else sqBracketCoeffZ p 0 i) • mode (W.Y a) i b

theorem vSqZ_zero_eq (W : VertexOperatorAlgebra 𝒱) (a : 𝒱) (p : ℕ) (b : 𝒱) :
    vSqZ W a p 0 b = mode (W.Y a) 0 b + sqDefect0 W a p b := by

theorem sqDefect0_mem_of_lower (W : VertexOperatorAlgebra 𝒱) {p q n : ℤ} {a b : 𝒱}
    (ha : a ∈ W.grading p) (hb : b ∈ W.grading q)
    (hpq : p + q - 1 = n) (T : Submodule ℂ 𝒱)
    (hT : ∀ m : ℤ, m < n → W.grading m ≤ T) :
    sqDefect0 W a p.toNat b ∈ T := by
\end{leancode}

\begin{dictionary}
\lean{mode (W.Y a) i b} is $a_{(i)}b$. \lean{sqDefect0} is the tail $\sum_{i\ne0}c_i a_{(i)}b$ (the \lean{∑ᶠ} has
finite support by \lean{sqDefect0_}\allowbreak\lean{support_}\allowbreak\lean{finite}). \lean{sqDefect0_mem_of_lower} is the lower-weight statement in
the form used by an induction: any subspace $T$ containing all $V_m$, $m<p+q-1$, contains the tail. \lean{p.toNat} is
$p$ as a natural number (harmless: $p\ge0$ whenever $a\ne0$). Same as Theorem~\ref{l16:thm:triang}.
\end{dictionary}

\textbf{Proof on paper.} $c^{(p,0)}_j$ is the coefficient of $z^j$ in $(1+z)^{p-1}$: it is $1$ at $j=0$ and $0$ for
$j<0$, and $a_{(i)}b\in V_{p+q-i-1}$, which for $i\ge1$ is below $p+q-1$. \qed

\begin{physics}
The square-bracket modes are the modes of the field on the \emph{cylinder} (or torus) coordinate $w$, $z = e^w-1$.
This statement says: the charge $\oint a$ measured on the torus equals the charge on the plane up to corrections
of lower energy. Hence ``torus one-point functions kill $a[0]v$'' (global symmetry Ward identity on the torus)
becomes ``they kill Poisson brackets modulo lower energy''.
\end{physics}

\begin{status}
Proved, unconditional (file audit, \lean{[propext, Classical.}\allowbreak\lean{choice, Quot.}\allowbreak\lean{sound]}).
\end{status}

\subsection{Proposition 5.2}

Let $\mathcal G = \C[G_2,G_4,G_6]\subseteq\C[[q]]$ be the algebra generated by Zhu's normalised Eisenstein series
$G_{2j}(q) = -\frac{B_{2j}}{(2j)!} + \frac{2}{(2j-1)!}\sum_{n\ge1}\sigma_{2j-1}(n)q^n$ (\lean{GSeries}, Chapter~\ref{ch:L14};
the Lean coefficient sequence \lean{zhuG} is quoted in Definition~\ref{l16:def:formalblock}), and $\lsixteenVA = \mathcal G\otimes_\C V$ (\lean{𝒱G 𝒱}, with
\lean{ιG v} $=1\otimes v$). Zhu's relation module $\lsixteenOq\subseteq\lsixteenVA$ (\lean{blockRelSpanG W}) is the
$\mathcal G$-span of
\begin{equation*}
a[0]v \qquad\text{and}\qquad a[-2]v + \sum_{m\ge0}(2m+3)\,G_{2m+4}\otimes a[2m+2]v ,
\end{equation*}
for homogeneous $a$ of weight $\ge1$ and $v\in V$; these are exactly the relations satisfied by every genus-one
one-point function (Chapter~\ref{ch:L15}).

\begin{leancode}
noncomputable def blockRelSpanG : Submodule (↥GSeries) (𝒱G 𝒱) :=
  Submodule.span (↥GSeries)
    ({x : 𝒱G 𝒱 | ∃ (ka : ℕ) (a v : 𝒱), 1 ≤ ka ∧ a ∈ W.grading (ka : ℤ) ∧
        x = ιG (vSqZ W a ka 0 v)} ∪
     {x : 𝒱G 𝒱 | ∃ (ka K : ℕ) (a v : 𝒱), 1 ≤ ka ∧ a ∈ W.grading (ka : ℤ) ∧
        (∀ m : ℕ, K ≤ m → vSqZ W a ka (2 * (m + 1) - 2 : ℤ) v = 0) ∧
        x = ιG (vSqZ W a ka (-2) v) + ∑ m ∈ Finset.range K,
              coefG m • ιG (vSqZ W a ka (2 * (m + 2) - 2 : ℤ) v)})
\end{leancode}
\noindent File \leanpath{VOA/\allowbreak Modularity/\allowbreak BlockRelSpanG.lean} (Chapter~\ref{ch:L15}).
Here \lean{coefG m} is $(2(m+2)-1)G_{2(m+2)}$, and the cut-off $K$ beyond which all $a[2m]v$ vanish makes the
sum finite.

\begin{theorem}[Arakawa--Kawasetsu Proposition 5.2]\label{l16:thm:prop52}
If $V$ is Etingof-finite, there is a finite set $T\subseteq\lsixteenVA$ with
$\lsixteenVA = \mathcal G\cdot T + \lsixteenOq$; i.e.\ $\lsixteenVA/\lsixteenOq$ is a finitely generated
$\mathcal G$-module.
\end{theorem}
\noindent File \leanpath{VOA/\allowbreak Modules/\allowbreak QuasiLisseFinite.lean}.
\begin{leancode}
theorem exists_finite_gen_sup_blockRelSpanG_of_etingofFinite (hzero : VacuumGradeZero W)
    (hE : EtingofFinite W) :
    ∃ T : Set (𝒱G 𝒱), T.Finite ∧
      (⊤ : Submodule (↥GSeries) (𝒱G 𝒱))
        = Submodule.span (↥GSeries) T ⊔ blockRelSpanG W := by
\end{leancode}

\begin{dictionary}
\lean{VacuumGradeZero W} is \lean{W.grading 0 ≤ Submodule.span ℂ ({W.vacuum} : Set 𝒱)}, i.e.\ $V_0\subseteq\C\lsixteenvac$
(file \leanpath{VOA/\allowbreak Modularity/\allowbreak BlockZeroMode.lean}). It is \textbf{free}: every library VOA has
$V_0=\C\lsixteenvac$ as an axiom, and \lean{QuasiLisseSl2.}\allowbreak\lean{vacuumGradeZero_}\allowbreak\lean{of_vacuum_}\allowbreak\lean{grade} proves
\lean{VacuumGradeZero W} for every \lean{W} (Section~\ref{l16:sec:payoff}). \lean{hE} is the only real hypothesis.
\lean{⊤ = span T ⊔ blockRelSpanG W} says every element of $\lsixteenVA$ is a $\mathcal G$-combination of $T$ plus a
relation, which is finite generation of the quotient. The docstring explains why the quotient module itself is not
named in the statement: a typeclass mismatch prevents restating \lean{Module.}\allowbreak\lean{Finite (↥GSeries) (𝒱G 𝒱 ⧸ blockRelSpanG W)}
outside its home file; the formulation with the finite generating family is equivalent. Same as AK Proposition~5.2
with ``quasi-lisse'' replaced by the weaker Etingof-finiteness and AK's $\mathcal A$ replaced by
$\C[G_2,G_4,G_6]$. (An $E_4,E_6$-version, \lean{exists_}\allowbreak\lean{finite_}\allowbreak\lean{gen_sup_}\allowbreak\lean{blockRelSpanE_}\allowbreak\lean{of_etingofFinite}, is in
Section~\ref{l16:sec:mlde}.)
\end{dictionary}

\textbf{Proof on paper.} By Theorem~\ref{l16:thm:O2P} choose finitely many homogeneous $s_1,\dots,s_g$ of positive
weight with $V = \operatorname{span}\{\lsixteenvac,s_i\} + O_2^P(V)$ (lift a basis of the finite-dimensional graded
space $V/O_2^P$). Let $T = \{1\otimes\lsixteenvac,1\otimes s_i\}$ and $N = \mathcal G T + \lsixteenOq$. We prove
$1\otimes V_n\subseteq N$ by strong induction on $n$. A vector of weight $n$ is a combination of the $s_i$ of weight $n$,
of homogeneous generators $a_{(-2)}b$ with $p+q+1 = n$, and of homogeneous generators $a_{(0)}b$ with $p+q-1 = n$.
\begin{itemize}
\item For $a_{(-2)}b$: $a[-2]b = a_{(-2)}b + (\text{weight}<n)$, and modulo $\lsixteenOq$,
  $a[-2]b\equiv-\sum(2m+3)G_{2m+4}a[2m+2]b$, each $a[2m+2]b$ of weight $<n$. This is Zhu's original step.
\item For $a_{(0)}b$ (the new step): by Theorem~\ref{l16:thm:triang} $a_{(0)}b = a[0]b - (\text{weight}<n)$, and
  $a[0]b\in\lsixteenOq$ exactly, with no Eisenstein tail.
\item If $a$ has weight $0$, then $a = c\lsixteenvac$ and $\lsixteenvac_{(0)}b = 0$, $\lsixteenvac_{(-2)}b = 0$.
\end{itemize}
Lower weights are in $N$ by induction. Hence $1\otimes V\subseteq N$, and since $1\otimes V$ generates $\lsixteenVA$ over
$\mathcal G$, $N = \lsixteenVA$. \qed

\begin{physics}
The torus one-point functions of a theory form a module over the ring of quasi-modular forms; this says that, if
there are finitely many classical amplitudes, every torus one-point function is determined by finitely many
``basic'' ones with modular-form coefficients. It is the algebraic shadow of ``finitely many conformal blocks on the
torus''.
\end{physics}

\begin{status}
Proved; \lean{VOA.exists_}\allowbreak\lean{finite_}\allowbreak\lean{gen_sup_}\allowbreak\lean{blockRelSpanG_}\allowbreak\lean{of_etingofFinite}: \lean{[propext, Classical.}\allowbreak\lean{choice, Quot.}\allowbreak\lean{sound]}
(re-run). Conditional only on \lean{hE}; \lean{hzero} is free.
\end{status}

\begin{theorem}[Growth of one-point functions from Etingof-finiteness]\label{l16:thm:growth}
If $V$ is Etingof-finite and $M$ is an ordinary module of conformal weight $h$ (that is, $M_h\ne0$ and every nonzero
weight space $M_\alpha$ has $\alpha\in h+\N$, Chapter~\ref{ch:L12}), then for every $v\in V$ the
coefficients of the torus one-point function $Z_M(v;q) = \operatorname{tr}_M o(v)q^{L_0-c/24}$ grow subexponentially:
for each $r>1$ there is $C$ with $|[q^{h-c/24+n}]Z_M(v;q)|\le C r^n$.
\end{theorem}
\noindent File \leanpath{VOA/\allowbreak Modules/\allowbreak QuasiLisseFinite.lean}.
\begin{leancode}
theorem onePointSubexp_of_etingofFinite {M : Type u} [AddCommGroup M] [Module ℂ M]
    (N : OrdVModule W M) {hw : ℂ} (hcw : HasConformalWeight N hw)
    (hzero : VacuumGradeZero W) (hE : EtingofFinite W) :
    OnePointSubexp N (hw - W.centralCharge / 24) := by
\end{leancode}
\noindent Files \leanpath{VOA/\allowbreak Modularity/\allowbreak AnalyticBlock.lean},
\leanpath{VOA/\allowbreak Modularity/\allowbreak ConvergenceBridge.lean}.
\begin{leancode}
def OnePointSubexp (N : OrdVModule W M) (h : ℂ) : Prop :=
  ∀ v : 𝒱, Subexponential (fun n : ℕ => ‖torusOnePoint N v (h + (n : ℂ))‖)

def Subexponential (a : ℕ → ℝ) : Prop :=
  ∀ r : ℝ, 1 < r → ∃ C : ℝ, ∀ n : ℕ, a n ≤ C * r ^ n
\end{leancode}

\begin{dictionary}
\lean{OrdVModule W M} is an ordinary module (Chapter~\ref{ch:L12}, Definition~\ref{l12:def:ordmodule});
\lean{HasConformalWeight N hw} says $M_{hw}\ne0$ and every nonzero weight lies in $hw+\N$
(file \leanpath{VOA/\allowbreak Modules/\allowbreak Rationality.lean}, Chapter~\ref{ch:L12}). A formal $q$-series is a
function \lean{FormalQSeries := ℂ → ℂ}, $e\mapsto[q^e]$ (Chapter~\ref{ch:L14}); \lean{torusOnePoint N v} is the
formal series $e\mapsto\operatorname{tr}_{M_{e+c/24}}o(v)$. \lean{Subexponential a} is the growth condition of
Theorem~\ref{l16:thm:growth}, which is equivalent to: $\sum a_nq^n$ converges on the open unit disc. So the
conclusion is: every one-point function converges for $|q|<1$, i.e.\ on the upper half plane. Same as Zhu's
convergence theorem, with $C_2$-cofinite weakened to Etingof-finite (Zhu's version is
\lean{onePointSubexp_of_c2}, Chapter~\ref{ch:L15}).
\end{dictionary}

\textbf{Proof on paper.} Let $S(v) = Z_M(v;q)$. By Theorem~\ref{l16:thm:prop52} and the fact that $S$ kills
$\lsixteenOq$ (Zhu's genus-one relations, \lean{IsFormalBlock}), the $\mathcal G$-span of all $S(v)$ is spanned by
finitely many $y_1,\dots,y_g$. Zhu's $\omega$-recursion $S(L[-2]v) = q\frac{d}{dq}S(v) +
\sum_{m}G_{2m+2}S(L[2m]v)$ shows that $q\frac{d}{dq}$ maps this span into its span over the larger ring of
subexponential power series; so $q\frac{d}{dq}y = Ay$ for a $g\times g$ matrix $A$ with subexponential entries. In
coefficients, $(h+n)y_n = \sum_j A_jy_{n-j}$, and a discrete Gr\"onwall argument shows the $y_n$, hence all
coefficients of all $S(v)$, grow subexponentially. \qed

\begin{physics}
Convergence of the torus partition function $\operatorname{tr}q^{L_0-c/24}$ for $|q|<1$, i.e.\ at every positive
temperature: the density of states grows slower than any exponential in the energy (for a CFT,
like $e^{2\pi\sqrt{c_{\mathrm{eff}}n/6}}$). For non-unitary or quasi-lisse theories this is not automatic, and here
it follows from finitely many classical amplitudes.
\end{physics}

\begin{status}
Proved; \lean{VOA.onePointSubexp_}\allowbreak\lean{of_etingofFinite}: \lean{[propext, Classical.}\allowbreak\lean{choice, Quot.}\allowbreak\lean{sound]} (re-run).
Non-vacuity: \lean{trivial_}\allowbreak\lean{onePointSubexp_}\allowbreak\lean{adjoint} instantiates it at the trivial VOA, and Section~\ref{l16:sec:payoff}
at $L_{-4/3}(\lsixteensl)$.
\end{status}

\begin{warning}[title={Stale ``no witness'' documentation (documentation only)}]
When this chapter was first written, the module docstring of \leanpath{QuasiLisseFinite.lean} (Status item~4 and the note of \S7) said that ``every theorem
here is currently instantiated only at VOAs that are already $C_2$-cofinite'', that the trivial VOA is ``the
\emph{only} unconditional instance'', and that $L_{-4/3}(\lsixteensl)$ is ``conditional on a
\lean{Sl2Presentation simpleSl2QL} bundle that is not constructed''. All three claims were already false at that time:
the bundle is constructed (Theorem~\ref{l16:thm:presL}), \lean{EtingofFinite} is unconditional for
$L_{-4/3}(\lsixteensl)$, and that algebra is proved not $C_2$-cofinite (Theorem~\ref{l16:thm:nonlisse}). The same stale
claims occurred in the headers of \leanpath{AKModularODE.lean}, \leanpath{AKUniformODE.lean} and
\leanpath{AKZhuModularity.lean} (its ``honest statement'' item~1). These docstrings have since been corrected (by a later revision of the library;
see the update in Section~\ref{l16:sec:nonlisse}).
\end{warning}

%======================================================================
\section{The character equation in formal form}\label{l16:sec:formal}
%======================================================================

\subsection{The formal objects}

We recall the formal $q$-series vocabulary of Chapter~\ref{ch:L14} that appears in the statements.

\noindent Files \leanpath{VOA/\allowbreak Modules/\allowbreak Character.lean},
\leanpath{VOA/\allowbreak Modularity/\allowbreak ZhuBlockConditions.lean},
\leanpath{VOA/\allowbreak Modularity/\allowbreak ZhuRecursion.lean},
\leanpath{VOA/\allowbreak Modularity/\allowbreak SubexpSeriesRing.lean},
\leanpath{VOA/\allowbreak Modularity/\allowbreak OnePointFunctions.lean}.
\begin{leancode}
abbrev FormalQSeries : Type := ℂ → ℂ

noncomputable def gradedDim (N : OrdVModule W M) : FormalQSeries :=
  fun e => (Module.finrank ℂ (N.weight (e + W.centralCharge / 24)) : ℂ)

noncomputable def qDeriv (χ : FormalQSeries) : FormalQSeries := fun e => e * χ e

noncomputable def qConv (f : ℕ → ℂ) (g : FormalQSeries) : FormalQSeries :=
  fun e => ∑ᶠ n : ℕ, f n * g (e - (n : ℂ))

def SubexpPS : Subalgebra ℂ (PowerSeries ℂ) where
  carrier := {F : PowerSeries ℂ | Subexponential fun n : ℕ => ‖PowerSeries.coeff n F‖}

noncomputable def torusPartition (N : OrdVModule W M) : FormalQSeries := gradedDim N
\end{leancode}

\begin{dictionary}
A formal $q$-series $\chi$ is the coefficient function $e\mapsto[q^e]\chi$, $e\in\C$. \lean{gradedDim N} is the
character $\chi_M = q^{-c/24}\sum_\alpha\dim M_\alpha\,q^\alpha$: the coefficient of $q^e$ is $\dim M_{e+c/24}$.
\lean{torusPartition N} is the same series under its physical name (\lean{rfl}). \lean{qDeriv} is the Euler operator
$\theta = q\frac{d}{dq}$, $q^e\mapsto e\,q^e$; \lean{qDeriv^[i]} is $\theta^i$. \lean{qConv f g} is the product of the
ordinary power series $\sum_nf_nq^n$ with $g$: $[q^e](fg) = \sum_nf_n[q^{e-n}]g$ (a finite sum on the series used,
whose exponents lie in $h+\N$). \lean{SubexpPS} is the subalgebra of $\C[[q]]$ of series with subexponential
coefficients, i.e.\ series converging on the open unit disc; it contains $\mathcal G=\C[G_2,G_4,G_6]$ (the
Eisenstein coefficients grow polynomially). We write $\mathcal S$ for \lean{SubexpPS}.
\end{dictionary}

\subsection{Square-bracket primaries and the Virasoro tower}

Zhu's square-bracket Virasoro operators are $L[n] = \omega[n+1] - \frac{c}{24}\lsixteenvac[n+1]$ (\lean{Lsq W n},
Chapter~\ref{ch:L15}); they satisfy the Virasoro algebra with the same central charge $c$.

\begin{definition}[Square-bracket primary]\label{l16:def:sqprimary}
$a\in V$ is square-bracket primary of weight $\Delta$ if $L[0]a = \Delta a$ and $L[n]a = 0$ for $n\ge1$. (Here $L[n]$ are the
Virasoro modes for Zhu's second VOA structure $Y[a,w] = Y(a,e^w-1)e^{w\,\mathrm{wt}\,a}$ on $V$, i.e.\ in the cylinder
coordinate; ordinary primaries, $L_na = 0$ for $n\ge1$, are a different notion in general. The vacuum is primary in both
senses.)
\end{definition}
\noindent File \leanpath{VOA/\allowbreak Modularity/\allowbreak AKModularODE.lean}.
\begin{leancode}
def IsSqPrimary (W : VertexOperatorAlgebra 𝒱) (a : 𝒱) (Δ : ℂ) : Prop :=
  Lsq W 0 a = Δ • a ∧ ∀ n : ℤ, 1 ≤ n → Lsq W n a = 0

theorem vacuum_isSqPrimary (W : VertexOperatorAlgebra 𝒱) : IsSqPrimary W W.vacuum 0 :=
\end{leancode}

\begin{dictionary}
Exactly Definition~\ref{l16:def:sqprimary}, which is the hypothesis of AK Lemma~5.5 (``Let $a$ be a primary vector of
weight $\Delta$, that is $L[0]a = \Delta a$ and $L[n]a = 0$ for $n\ge1$''). It is the \emph{square-bracket} notion, not
the ordinary one; the vacuum is primary of weight $0$ in both senses, which is all the character needs.
\end{dictionary}

\begin{physics}
A square-bracket primary is a primary field for the Virasoro algebra on the cylinder (torus) coordinate; the vacuum
(identity operator) is the primary whose one-point function is the partition function itself.
\end{physics}

\begin{theorem}[The Virasoro tower over a primary]\label{l16:thm:tower}
Let $a$ be square-bracket primary of weight $\Delta$, and consider the \emph{tower} $L[-2]^na$ ($a$ dressed by $n$ insertions
of the cylinder stress tensor). For all $n,j\ge0$,
$L[2j]\,L[-2]^n a = \gamma_{\Delta,c}(j,n)\,L[-2]^{n-j}a$, where $\gamma$ is defined by $\gamma(j,0) = \Delta\delta_{j0}$,
$\gamma(0,n+1) = \gamma(0,n)+2$, $\gamma(j+1,n+1) = (2j+4)\gamma(j,n) + \delta_{j0}\frac c2 + \gamma(j+1,n)$; and
$\gamma(j,n) = 0$ for $j>n$. In closed form $\gamma(0,n) = \Delta+2n$ and $\gamma(1,n) = n(4\Delta+4n-4+\frac c2)$.
\end{theorem}
\noindent File \leanpath{VOA/\allowbreak Modularity/\allowbreak AKModularODE.lean}.
\begin{leancode}
noncomputable def towerScalar (Δ c : ℂ) : ℕ → ℕ → ℂ
  | j, 0 => if j = 0 then Δ else 0
  | 0, (n + 1) => towerScalar Δ c 0 n + 2
  | (j + 1), (n + 1) =>
      (2 * (j : ℂ) + 4) * towerScalar Δ c j n
        + (if j = 0 then c / 2 else 0) + towerScalar Δ c (j + 1) n

theorem Lsq_tower (W : VertexOperatorAlgebra 𝒱) {a : 𝒱} {Δ : ℂ} (ha : IsSqPrimary W a Δ) :
    ∀ (n j : ℕ), Lsq W (2 * (j : ℤ)) ((⇑(Lsq W (-2)))^[n] a)
      = towerScalar Δ W.centralCharge j n • ((⇑(Lsq W (-2)))^[n - j] a) := by
\end{leancode}

\begin{dictionary}
\lean{(⇑(Lsq W (-2)))^[n] a} is $L[-2]^na$ (\lean{^[n]} is $n$-fold iteration of a function). \lean{n - j} is truncated
subtraction in $\N$; for $j>n$ the scalar is $0$ (\lean{towerScalar_}\allowbreak\lean{eq_zero_}\allowbreak\lean{of_lt}), so the truncation is harmless.
The closed forms are \lean{towerScalar_zero_left} and \lean{towerScalar_one_left}. Same as AK's
``$L[2k-2]L[-2]^{i-1}a = c_kL[-2]^{i-k}a$'' with $j = k-1$, $n = i-1$; the Lean scalar is computed rather than left
implicit, and it depends on $i$ as well as $k$ (AK suppress the dependence).
\end{dictionary}

\textbf{Proof on paper.} Induction on $n$. For $n=0$ this is primarity. For the step write
$L[2j]L[-2]y = L[-2]L[2j]y + (2j+2)L[2j-2]y + \delta_{j1}\frac{c}{12}(8-2)y$ with $y = L[-2]^na$ and use the induction
hypothesis on both terms; the case $j=0$ is $L[0]L[-2]y = L[-2](L[0]+2)y$. Collecting gives the recursion. \qed

\begin{physics}
$L[-2]$ inserts the stress tensor on the torus; $L[-2]^na$ is $a$ dressed by $n$ stress tensors. The statement is the
torus Ward identity's algebraic core: moving a positive Virasoro mode through the dressing produces only the same
tower with explicit coefficients (involving $c$ through the Schwarzian/central term).
\end{physics}

\begin{status}
Proved, unconditional; \lean{VOA.Lsq_tower}: \lean{[propext, Classical.}\allowbreak\lean{choice, Quot.}\allowbreak\lean{sound]} (re-run).
\end{status}

\subsection{Lemma 5.5: from $L[-2]$ to $q\,d/dq$}

\begin{definition}[Formal genus-one block; Zhu Def.~5.1.1, formal form]\label{l16:def:formalblock}
A linear map $S:V\to\{\text{formal }q\text{-series}\}$ is a \emph{formal block} if for every homogeneous $a$ of weight
$k_a\ge1$ and every $v\in V$:
(i) $S(a[0]v) = 0$;
(ii) $S(a[-2]v) + \sum_{m\ge0}(2m+1)\,G_{2m+2}\star S(a[2m]v) = 0$;
(iii) $S(L[-2]v) = \theta S(v) + \sum_{m\ge0}G_{2m+2}\star S(L[2m]v)$,
where $\star$ is the product of a power series with a formal $q$-series (\lean{qConv}) and $G_{2j}$ is Zhu's normalised
Eisenstein series (below).
\end{definition}
\noindent File \leanpath{VOA/\allowbreak Modularity/\allowbreak ZhuBlockConditions.lean} (Chapter~\ref{ch:L15}).
\begin{leancode}
structure IsFormalBlock (W : VertexOperatorAlgebra 𝒱) (S : 𝒱 →ₗ[ℂ] FormalQSeries) : Prop where
  vanish : ∀ (ka : ℕ), 1 ≤ ka → ∀ a ∈ W.grading (ka : ℤ), ∀ v : 𝒱, S (vSqZ W a ka 0 v) = 0
  negTwo : ∀ (ka : ℕ), 1 ≤ ka → ∀ a ∈ W.grading (ka : ℤ), ∀ v : 𝒱,
    S (vSqZ W a ka (-2) v)
      + ∑ᶠ m : ℕ, ((2 * (m + 1) - 1 : ℕ) : ℂ)
          • qConv (zhuG (m + 1)) (S (vSqZ W a ka (2 * (m + 1) - 2 : ℤ) v)) = 0
  omega : ∀ v : 𝒱, S (Lsq W (-2) v)
    = qDeriv (S v) + ∑ᶠ m : ℕ, qConv (zhuG (m + 1)) (S (Lsq W (2 * (m + 1) - 2 : ℤ) v))
\end{leancode}
\noindent File \leanpath{VOA/\allowbreak Modularity/\allowbreak ZhuNormalization.lean} (Chapter~\ref{ch:L14}).
\begin{leancode}
noncomputable def zhuG (j : ℕ) : ℕ → ℂ :=
  fun m =>
    if m = 0 then -(((bernoulli (2 * j) : ℚ) : ℂ) / (Nat.factorial (2 * j) : ℂ))
    else (2 / (Nat.factorial (2 * j - 1) : ℂ)) * lambertCoeff (2 * j - 1) m
\end{leancode}

\begin{dictionary}
\lean{vSqZ W a ka m} is the square-bracket mode $a[m]$; \lean{Lsq W n} is $L[n]$; \lean{qDeriv} is $\theta$; \lean{qConv f}
multiplies by the power series with coefficients $f$. \lean{zhuG j} is the coefficient sequence of Zhu's
$G_{2j}(q) = -\frac{B_{2j}}{(2j)!} + \frac{2}{(2j-1)!}\sum_{n\ge1}\sigma_{2j-1}(n)q^n$ (\lean{lambertCoeff (2j-1) n}
$=\sigma_{2j-1}(n)$), i.e.\ the Eisenstein series normalised so that $G_{2j} = \frac{2\zeta(2j)}{(2\pi i)^{2j}}E_{2j}$; the
index shift $m\mapsto m+1$ makes the three fields exactly (i)--(iii) above, with $(2m+1)G_{2m+2}$ in (ii). The \lean{∑ᶠ} sums
are finite because $a[n]v = 0$ and $L[n]v = 0$ for $n\gg0$. Same as Zhu's Definition 5.1.1 in formal $q$-series form (Zhu
uses holomorphic functions on $\mathbb H$; that is the analytic version, Definition~\ref{l16:def:analyticblock}).
$\mathcal G = \C[G_2,G_4,G_6]$ (\lean{GSeries}) is the algebra generated by the series \lean{PowerSeries.}\allowbreak\lean{mk (zhuG j)},
$j = 1,2,3$; it contains every $G_{2j}$.
\end{dictionary}

\begin{physics}
A genus-one block is a candidate torus one-point function $v\mapsto\langle v\rangle_\tau$: (i) is the Ward identity for the
global symmetry charges on the torus, (ii) expresses that ``derivative'' insertions are total derivatives up to Eisenstein
contact terms, and (iii) is the conformal Ward identity --- inserting the stress tensor equals differentiating in the modular
parameter $\tau$ ($\theta = \frac1{2\pi i}\partial_\tau$), plus contact terms.
\end{physics}

\begin{status}
Definition; satisfiable: \lean{torusOnePointHom_}\allowbreak\lean{isFormalBlock} (every module's trace functions form a formal block).
\end{status}

The basic example is $S = Z_M(\cdot;q)$
(\lean{torusOnePointHom N}, a block by \lean{torusOnePointHom_}\allowbreak\lean{isFormalBlock}). \lean{QSeriesAt h} is the space of
series supported on $h+\N$, and \lean{psSMul F χ} is \lean{qConv} with the coefficient sequence of $F\in\C[[q]]$.

\begin{theorem}[AK Lemma 5.5, for primary vectors]\label{l16:thm:huni}
Let $S$ be a formal block supported on $h+\N$ and $a$ square-bracket primary. For each $i$ there are
$\eta_{i,0},\dots,\eta_{i,i}\in\mathcal S$ with $\eta_{i,i} = 1$ and
$S(L[-2]^ia) = \sum_{j=0}^{i}\eta_{i,j}\,\theta^jS(a)$.
\end{theorem}
\noindent File \leanpath{VOA/\allowbreak Modularity/\allowbreak AKModularODE.lean}.
\begin{leancode}
theorem huni_of_sqPrimary (W : VertexOperatorAlgebra 𝒱)
    {S : 𝒱 →ₗ[ℂ] FormalQSeries} {hbase : ℂ} (hsupp : ∀ v, S v ∈ QSeriesAt hbase)
    (hS : IsFormalBlock W S) {a : 𝒱} {Δ : ℂ} (ha : IsSqPrimary W a Δ) (i : ℕ) :
    ∃ hc : ℕ → PowerSeries ℂ, (∀ j, hc j ∈ SubexpPS) ∧ hc i = 1 ∧ (∀ j, i < j → hc j = 0) ∧
      S ((⇑(Lsq W (-2)))^[i] a)
        = ∑ j ∈ Finset.range (i + 1), psSMul (hc j) (qDeriv^[j] (S a)) := by
\end{leancode}

\begin{dictionary}
\lean{hc j} is $\eta_{i,j}$; ``\lean{hc i = 1} and \lean{hc j = 0} for $j>i$'' is unitriangularity. The coefficients
are in $\mathcal S$, not in AK's $\C[G_4,G_6]$ (see Section~\ref{l16:sec:weak}), and in the Euler basis $\theta^j$, not
the Serre basis. They are produced by an existence statement inside the span of the block values, so they may depend on
$S$; AK's depend only on $(a,\Delta,c)$. \textbf{Why primarity is essential.} The file \leanpath{UniCollapse.lean}
(Chapter~\ref{ch:L15}) refutes the same conclusion quantified over all $v$, via $v = L[-1]x$, which is not primary; the
present statement is for a fixed primary $a$, which is the only case AK use.
\end{dictionary}

\textbf{Proof on paper.} Induction on $i$. By (iii) and Theorem~\ref{l16:thm:tower},
$S(L[-2]^{i+1}a) = \theta S(L[-2]^ia) + \sum_{m}G_{2m+2}\gamma(m,i)S(L[-2]^{i-m}a)$. By induction each
$S(L[-2]^{i'}a)$, $i'\le i$, is a unitriangular $\mathcal S$-combination of $\theta^jS(a)$; applying $\theta$ uses the
Leibniz rule $\theta(\eta\chi) = (\theta\eta)\chi + \eta\theta\chi$, and $\theta\eta\in\mathcal S$. The leading
coefficient stays $1$. \qed

\begin{physics}
Inserting the stress tensor on the torus is the same as differentiating with respect to the modular parameter
(``$T(z)$ is $\partial_\tau$'') up to contact terms with Eisenstein coefficients. This is the conformal Ward identity on
the torus.
\end{physics}

\begin{status}
Proved, unconditional; \lean{VOA.huni_of_sqPrimary}: \lean{[propext, Classical.}\allowbreak\lean{choice, Quot.}\allowbreak\lean{sound]} (re-run).
\end{status}

\subsection{Lemma 5.3 and the formal character equation}

\begin{theorem}[AK Lemma 5.3, block form]\label{l16:thm:lemma53}
If $V$ is Etingof-finite and $S$ is a formal block, then for every $v$ there are $s\ge1$ and $G_0,\dots,G_{s-1}\in\mathcal G$
with $S(L[-2]^sv) = \sum_{i<s}G_i\,S(L[-2]^iv)$.
\end{theorem}
\noindent File \leanpath{VOA/\allowbreak Modularity/\allowbreak AKModularODE.lean}.
\begin{leancode}
theorem exists_monic_Lsq_relation_of_etingofFinite {S : 𝒱 →ₗ[ℂ] FormalQSeries} {hbase : ℂ}
    (hsupp : ∀ v, S v ∈ QSeriesAt hbase) (hS : IsFormalBlock W S)
    (hzero : VacuumGradeZero W) (hE : EtingofFinite W) (v : 𝒱) :
    ∃ (s : ℕ) (G : ℕ → PowerSeries ℂ), 1 ≤ s ∧ (∀ i, G i ∈ GSeries) ∧
      S ((⇑(Lsq W (-2)))^[s] v)
        = ∑ i ∈ Finset.range s, psSMul (G i) (S ((⇑(Lsq W (-2)))^[i] v)) := by
\end{leancode}

\begin{dictionary}
\lean{G i ∈ GSeries}: the coefficients are in $\mathcal G = \C[G_2,G_4,G_6]$. The quantifier order is
$\forall S\,\forall v\,\exists s\,\exists G$: $s$ and $G$ may depend on the block $S$ (hence on the module). AK's
Lemma~5.3 gives one relation in $\lsixteenVA$ independent of the module; that is the uniform version
(Theorem~\ref{l16:thm:uniform}).
\end{dictionary}

\textbf{Proof on paper.} $\mathcal G$ is Noetherian (a finitely generated $\C$-algebra) and the $\mathcal G$-span of the
values of $S$ is finitely generated (Theorem~\ref{l16:thm:prop52}, as in the proof of Theorem~\ref{l16:thm:growth}),
hence Noetherian. The ascending chain $\operatorname{span}_{\mathcal G}\{S(L[-2]^iv): i\le r\}$ stabilises, so some
$S(L[-2]^{n+1}v)$ is a combination of the earlier ones. \qed

\begin{physics}
Finitely many independent torus amplitudes force the stress-tensor tower to close: dressing a state by enough stress
tensors produces nothing new.
\end{physics}

\begin{status}
Proved, conditional on \lean{hE} (\lean{hzero} free); \lean{[propext, Classical.}\allowbreak\lean{choice, Quot.}\allowbreak\lean{sound]} (re-run).
\end{status}

\begin{theorem}[Formal character equation (``AK equation (2), formal form'' in the source); \emph{not} AK Theorem~5.1, near-vacuous]\label{l16:thm:formalODE}
Let $V$ be Etingof-finite and $M$ an ordinary module with a conformal weight. There are $s\ge1$ and power series
$g_0,\dots,g_{s-1}\in\mathcal S$ such that the character satisfies
\begin{equation}\label{l16:eq:formalODE}
\theta^s\chi_M + \sum_{i<s}g_i\,\theta^i\chi_M = 0 .
\end{equation}
\textbf{Caution.} This is not a modular linear differential equation and not Arakawa--Kawasetsu's
Theorem~5.1: $\theta$ is not the Serre derivative, the $g_i$ are arbitrary convergent power series rather than
modular forms, and $s,g$ depend on $M$. By Proposition~\ref{l16:prop:weak} the conclusion holds for every
convergent $q$-series with nonzero leading coefficient and zeros of bounded multiplicity in the unit disc, so it
says almost nothing beyond convergence of $\chi_M$ (Theorem~\ref{l16:thm:growth}).
\end{theorem}
\noindent File \leanpath{VOA/\allowbreak Modularity/\allowbreak AKModularODE.lean}.
\begin{leancode}
theorem character_formalODE_of_etingofFinite (W : VertexOperatorAlgebra 𝒱) (N : OrdVModule W M)
    {hw : ℂ} (hcw : HasConformalWeight N hw) (hzero : VacuumGradeZero W) (hE : EtingofFinite W) :
    ∃ (s : ℕ) (g : ℕ → ℕ → ℂ), 1 ≤ s ∧ (∀ i, Subexponential (fun k => ‖g i k‖)) ∧
      ∀ e : ℂ, (qDeriv^[s] (torusPartition N)) e
        + ∑ i ∈ Finset.range s, qConv (g i) (qDeriv^[i] (torusPartition N)) e = 0 := by

theorem character_formalODE_adjoint_of_etingofFinite (W : VertexOperatorAlgebra 𝒱)
    (hne : W.vacuum ≠ 0) (hzero : VacuumGradeZero W) (hE : EtingofFinite W) :
    ∃ (s : ℕ) (g : ℕ → ℕ → ℂ), 1 ≤ s ∧ (∀ i, Subexponential (fun k => ‖g i k‖)) ∧
      ∀ e : ℂ, (qDeriv^[s] (torusPartition (adjointModule W))) e
        + ∑ i ∈ Finset.range s,
            qConv (g i) (qDeriv^[i] (torusPartition (adjointModule W))) e = 0 :=
\end{leancode}

\begin{dictionary}
\lean{g i : ℕ → ℂ} is the coefficient sequence of $g_i$, and \lean{Subexponential (fun k => ‖g i k‖)} says
$g_i\in\mathcal S$. The equation is asserted coefficientwise, \lean{∀ e : ℂ}, i.e.\ as an identity of formal
series. \lean{adjointModule W} is $V$ as a module over itself (Chapter~\ref{ch:L12},
Definition~\ref{l12:def:adjoint}); \lean{hne} (nonzero vacuum) makes its conformal weight $0$. \textbf{Comparison with
AK.} AK's~\eqref{l16:eq:mlde} uses (1) iterated Serre derivatives, (2) coefficients $p_i\in\C[E_4,E_6]$ of weight
$2(s-i)$, (3) one equation for all $M$. The Lean statement has (1) the Euler operator $\theta$ instead --- this alone
would be harmless, since $\vartheta^i_{(0)} = \theta^i + (\text{lower terms with coefficients in }\C[E_2,E_4,E_6])$ is a
unitriangular change of basis; (2) coefficients in $\mathcal S$, an enormous ring; (3) $s$ and $g$ depending on $M$. The
statement is therefore much \textbf{weaker} than AK's equation, and the docstrings say so (``This is NOT
Arakawa--Kawasetsu Theorem 5.1 and NOT a modular linear differential equation''). How much weaker is the subject of
Section~\ref{l16:sec:weak}.
\end{dictionary}

\textbf{Proof on paper.} Apply Theorem~\ref{l16:thm:lemma53} to $S = Z_M(\cdot;q)$ and $v=\lsixteenvac$, and substitute
Theorem~\ref{l16:thm:huni} (the vacuum is primary of weight $0$) into the monic relation: since $\eta_{i,i}=1$ the result is
monic in $\theta^s$, with coefficients $g_j = \eta_{s,j} - \sum_iG_i\eta_{i,j}\in\mathcal S$. Finally
$Z_M(\lsixteenvac;q) = \chi_M$. \qed

\begin{physics}
In AK's form this would be the statement that the Schur index (vacuum character) of a four-dimensional $\mathcal N=2$
theory satisfies an MLDE. In the formal form proved here it is the statement that the stress-tensor tower closes on
the torus with coefficients that are convergent functions of $q$ --- a statement about the vacuum character's
analytic behaviour at $|q|<1$ rather than about its modular properties.
\end{physics}

\begin{status}
Proved; \lean{VOA.character_}\allowbreak\lean{formalODE_}\allowbreak\lean{of_etingofFinite}: \lean{[propext, Classical.}\allowbreak\lean{choice, Quot.}\allowbreak\lean{sound]} (re-run).
\textbf{Content warning:} near-vacuous (Proposition~\ref{l16:prop:weak}, argued on paper); \textbf{not} AK Theorem~5.1
and not an MLDE.
Conditional on \lean{hE} (\lean{hzero} free). Instantiated at the trivial VOA (\lean{trivial_}\allowbreak\lean{character_}\allowbreak\lean{formalODE})
and, unconditionally, at $L_{-4/3}(\lsixteensl)$ (Section~\ref{l16:sec:payoff}).
\end{status}

\subsection{Uniformity in the module}

\begin{theorem}[AK Lemma 5.3 in $\lsixteenVA$; one order for all modules]\label{l16:thm:uniform}
Let $V$ be Etingof-finite and $a\in V$. (i) There are $s\ge1$ and $G_i\in\mathcal G$ such that \emph{every} formal block
$S$ satisfies $S(L[-2]^sa) = \sum_{i<s}G_iS(L[-2]^ia)$. (ii) There is $s\ge1$ such that for every ordinary module $M$ with a
conformal weight there are $g_i\in\mathcal S$ with $\theta^s\chi_M + \sum_{i<s}g_i\theta^i\chi_M = 0$.
\end{theorem}
\noindent File \leanpath{VOA/\allowbreak Modularity/\allowbreak AKUniformODE.lean}.
\begin{leancode}
theorem exists_uniform_monic_Lsq_relation_of_etingofFinite
    (hzero : VacuumGradeZero W) (hE : EtingofFinite W) (a : 𝒱) :
    ∃ (s : ℕ) (G : ℕ → PowerSeries ℂ), 1 ≤ s ∧ (∀ i, G i ∈ GSeries) ∧
      ∀ (hbase : ℂ) (S : 𝒱 →ₗ[ℂ] FormalQSeries) (hsupp : ∀ v, S v ∈ QSeriesAt hbase),
        IsFormalBlock W S →
          S ((⇑(Lsq W (-2)))^[s] a)
            = ∑ i ∈ Finset.range s, psSMul (G i) (S ((⇑(Lsq W (-2)))^[i] a)) := by

theorem uniform_order_character_formalODE_of_etingofFinite (W : VertexOperatorAlgebra 𝒱)
    (hzero : VacuumGradeZero W) (hE : EtingofFinite W) :
    ∃ s : ℕ, 1 ≤ s ∧
      ∀ {M : Type u} [AddCommGroup M] [Module ℂ M] (N : OrdVModule W M) (hw : ℂ),
        HasConformalWeight N hw →
          ∃ g : ℕ → ℕ → ℂ, (∀ i, Subexponential (fun k => ‖g i k‖)) ∧
            ∀ e : ℂ, (qDeriv^[s] (torusPartition N)) e
              + ∑ i ∈ Finset.range s,
                  qConv (g i) (qDeriv^[i] (torusPartition N)) e = 0 := by
\end{leancode}

\begin{dictionary}
Quantifier order is the content. In (i), $\exists s\,\exists G$ precedes $\forall S$: the relation is uniform, exactly
AK's Lemma~5.3. In (ii), $\exists s$ precedes $\forall M$ but $\exists g$ follows it: \textbf{only the order is uniform}, the
coefficients depend on $M$. The docstring is explicit: ``one order, many equations'', not ``one equation''. The reason
is that Lemma~5.5 (Theorem~\ref{l16:thm:huni}) produces block-dependent coefficients. The universe of \lean{M} is that of
\lean{𝒱}. So (ii) is weaker than AK's closing sentence (all characters solve one equation).
\end{dictionary}

\textbf{Proof on paper.} (i) Run the Noetherian argument of Theorem~\ref{l16:thm:lemma53} in the finitely generated
module $\lsixteenVA/\lsixteenOq$ itself, which is independent of any block, to get
$1\otimes L[-2]^sa - \sum G_i\otimes L[-2]^ia\in\lsixteenOq$; every block kills $\lsixteenOq$. (ii) Combine (i) with
Theorem~\ref{l16:thm:huni} as before. \qed

\begin{physics}
AK's statement ``all characters lie in the solution space of one MLDE'' is what turns ``finitely many ordinary modules''
into a theorem: an order-$s$ equation has an $s$-dimensional solution space. With module-dependent coefficients this
conclusion is lost.
\end{physics}

\begin{status}
Proved, conditional on \lean{hE}; both \lean{[propext, Classical.}\allowbreak\lean{choice, Quot.}\allowbreak\lean{sound]} (re-run). Instantiated at the
trivial VOA (\lean{trivial_}\allowbreak\lean{uniform_}\allowbreak\lean{order_}\allowbreak\lean{character_}\allowbreak\lean{formalODE}) and at $L_{-4/3}(\lsixteensl)$.
\end{status}

\subsection{How weak is the formal equation?}\label{l16:sec:weak}

The following observation is ours; it is argued on paper and not machine-checked.

\begin{proposition}[The formal equation is almost a growth statement]\label{l16:prop:weak}
Let $\chi = q^h F(q)$ with $F = \sum_{n\ge0}F_nq^n$, $F_0\ne0$, and suppose $F$ converges on the open unit disc $D$ and its
zeros in $D$ have multiplicity at most $m$. Then for $s = m+1$ there are $g_0,\dots,g_{s-1}\in\mathcal S$ with
$\theta^s\chi + \sum_{i<s}g_i\theta^i\chi = 0$.
\end{proposition}

\begin{proof}[Proof on paper]
Since $\theta(q^hF) = q^h(h+\theta)F$, it suffices to solve $(h+\theta)^sF\in\sum_{i<s}\mathcal O(D)\,(h+\theta)^iF$, where
$\mathcal O(D)$ is the ring of holomorphic functions on $D$; the span on the right equals
$\sum_{i<s}\mathcal O(D)\theta^iF$. The functions $F,\theta F,\dots,\theta^{s-1}F$ have no common zero in $D$: at $q=0$,
$F(0)=F_0\ne0$; at $q_0\ne0$, $\theta = q\,d/dq$ differs from $d/dq$ by a unit, and a zero of $F$ of order $\le m<s$ is not
a zero of its $m$-th derivative. In $\mathcal O(D)$ a finitely generated ideal without common zeros is the unit ideal (a
classical theorem of Wedderburn and Helmer: $\mathcal O(D)$ is a B\'ezout domain, and the generator of a finitely generated
ideal is a gcd, which here has no zeros). So $1 = \sum_{i<s}u_i\theta^iF$ with $u_i\in\mathcal O(D)$, and
$\theta^sF = \sum_i(u_i\theta^sF)\theta^iF$. Finally, a function holomorphic on $D$ has Taylor coefficients with
$|a_n|\le C_rr^n$ for every $r>1$ (Cauchy's estimate on the circle of radius $1/r$), i.e.\ lies in $\mathcal S$. Rewriting
$(h+\theta)^s$ in terms of $\theta^i$ changes coefficients by constants.
\end{proof}

\begin{physics}
A linear ODE in $q$ with coefficients merely holomorphic in the disc constrains nothing physical: the content of an MLDE is
the \emph{modularity} and finite dimensionality of the coefficient space, which ties the low-temperature expansion to the
high-temperature (Cardy) behaviour.
\end{physics}

\begin{warning}[title={The formal ODE theorems say little beyond convergence (moderate)}]
By Proposition~\ref{l16:prop:weak}, the conclusion of \lean{character_}\allowbreak\lean{formalODE_}\allowbreak\lean{of_etingofFinite} (and of its
instances and uniform-order variant) holds for \emph{every} convergent $q$-series with nonzero leading coefficient
and zeros of bounded multiplicity in the unit disc, whether or not it has anything to do with a vertex algebra or with
modularity. Together with Theorem~\ref{l16:thm:growth}, which already gives convergence, the formal ODE theorems add
essentially only the bounded-multiplicity information. The library's docstrings correctly say the result is ``not an
MLDE'' and ``genuinely weaker'' on the coefficient axis, but they present it as ``AK equation~(2)'', and the
Lean names (\lean{character_}\allowbreak\lean{formalODE_}\allowbreak\lean{...}) invite over-reading. A referee should treat these theorems as growth
statements. The genuine MLDE is Section~\ref{l16:sec:mlde}. (Also: the uniform order of
Theorem~\ref{l16:thm:uniform}(ii) is upward closed --- applying $\theta$ to an order-$s$ equation with coefficients in
$\mathcal S$ gives one of order $s+1$ --- so ``one order'' is likewise weak.)
\end{warning}

%======================================================================
\section{The genuine modular differential equation}\label{l16:sec:mlde}
%======================================================================

The file \leanpath{AKMLDE.lean} keeps the three rungs that \leanpath{AKModularODE.lean} had to merge apart, and works
analytically on the upper half plane with genuine modular forms.

\subsection{Modular coefficients and Serre derivatives}

\noindent Files \leanpath{VOA/\allowbreak Modularity/\allowbreak LevelOneStructure.lean},
\leanpath{VOA/\allowbreak Modularity/\allowbreak EisensteinRamanujan.lean},
\leanpath{VOA/\allowbreak Modularity/\allowbreak SerreDerivative.lean} (Chapter~\ref{ch:L14}).
\begin{leancode}
noncomputable def E46 : Subalgebra ℂ (UpperHalfPlane → ℂ) :=
  Algebra.adjoin ℂ {(⇑ModularForm.E₄ : UpperHalfPlane → ℂ), (⇑ModularForm.E₆ : UpperHalfPlane → ℂ)}

noncomputable def e46C (f : E46) (z : ℂ) : ℂ :=
  (f : UpperHalfPlane → ℂ) (UpperHalfPlane.ofComplex z)

def E46Weight (k : ℤ) (f : E46) : Prop :=
  ∃ F : ModularForm 𝒮ℒ k, (f : UpperHalfPlane → ℂ) = F

noncomputable def serreDerivC (k : ℤ) (f : ℂ → ℂ) (z : ℂ) : ℂ :=
  (2 * (Real.pi : ℂ) * Complex.I)⁻¹ * _root_.deriv f z
    + (k : ℂ) * ((2 * (Real.pi : ℂ) * Complex.I)⁻¹) ^ 2 * (G2C z * f z)

def E46SerreClosed : Prop :=
  ∀ (k : ℤ) (f : E46), E46Weight k f →
    ∃ g : E46, E46Weight (k + 2) g ∧
      ∀ z ∈ imGT 0, e46C g z = serreDerivC k (e46C f) z
\end{leancode}

\begin{dictionary}
\lean{E46} is the algebra of functions on $\mathbb H$ generated by $E_4,E_6$ (Mathlib's normalised Eisenstein series);
\lean{e46C} extends such a function to $\C$ by Mathlib's retraction (values off $\mathbb H$ are never used: every
statement is restricted to \lean{imGT 0} $=\{\operatorname{Im}z>0\}$). \lean{E46Weight k f}: $f$ is (the function of) a
modular form of weight $k$ for $SL_2(\Z)$ (\lean{𝒮ℒ}), holomorphic at the cusp. \lean{serreDerivC k f} is
$\frac{1}{2\pi i}f' + k(2\pi i)^{-2}G_2f$; with the normalisation $E_2 = G_2/(2\zeta(2)) = 3G_2/\pi^2$ recorded in
the source, this is $\frac1{2\pi i}f' - \frac{k}{12}E_2f$, the Serre derivative $\vartheta_k$ (Chapter~\ref{ch:L14}).
\lean{E46SerreClosed} says: $\vartheta_k$ maps weight-$k$ elements of $\C[E_4,E_6]$ to weight-$(k+2)$ elements. This is true
(Ramanujan's identities $\vartheta_4E_4 = -\frac13E_6$, $\vartheta_6E_6 = -\frac12E_4^2$ and the Leibniz rule), and it is
\textbf{proved} in the library: \lean{e46SerreClosed} (\leanpath{VOA/\allowbreak Modularity/\allowbreak EisensteinSerreClosure.lean},
Chapter~\ref{ch:L14}), via \lean{serreForm}, which shows that the Serre derivative of a level-one modular form is a
level-one modular form. Mathlib defines the derivative of modular forms but does not prove this preservation. The
Arakawa--Kawasetsu MLDE theorems of this chapter keep \lean{E46SerreClosed} as a hypothesis \lean{hclose} in their
statements; the primed versions
\lean{IsSqPrimary.}\allowbreak\lean{exists_}\allowbreak\lean{monicMLDEOn_}\allowbreak\lean{all_analyticBlocks_}\allowbreak\lean{of_etingofFinite'}
and \lean{IsSqPrimary.}\allowbreak\lean{exists_}\allowbreak\lean{monicMLDEOn_}\allowbreak\lean{all_analyticBlocks_}\allowbreak\lean{of_fin_gen'}
discharge it.
\end{dictionary}

\begin{definition}[Monic MLDE on $\mathbb H$]\label{l16:def:mlde}
For $f:\mathbb H\to\C$, $f$ satisfies a monic MLDE of order $s\ge1$ and weight $w$ with coefficients $p_0,\dots,p_{s-1}$
if each $p_i$ is a modular form of weight $2(s-i)$ and $\vartheta^s_{(w)}f + \sum_{i<s}p_i\vartheta^i_{(w)}f = 0$ on $\mathbb H$.
\end{definition}
\noindent File \leanpath{VOA/\allowbreak Modularity/\allowbreak AKMLDE.lean}.
\begin{leancode}
noncomputable def serreIterC (wt : ℤ) : ℕ → (ℂ → ℂ) → (ℂ → ℂ)
  | 0, f => f
  | i + 1, f => serreDerivC (wt + 2 * (i : ℤ)) (serreIterC wt i f)

def MonicMLDEOn (wt : ℤ) (f : ℂ → ℂ) (s : ℕ) (p : ℕ → E46) : Prop :=
  1 ≤ s ∧
    (∀ i < s, E46Weight (2 * ((s - i : ℕ) : ℤ)) (p i)) ∧
    ∀ z ∈ imGT 0, serreIterC wt s f z
      + ∑ i ∈ Finset.range s, e46C (p i) z * serreIterC wt i f z = 0
\end{leancode}

\begin{dictionary}
\lean{serreIterC wt i f} is $\vartheta^i_{(w)}f$ with the weight raised by $2$ at each step. \lean{MonicMLDEOn} is exactly
Definition~\ref{l16:def:mlde}; the weights of the coefficients are forced, as in AK's~\eqref{l16:eq:mlde}. This is a
genuine MLDE: unlike \lean{SubexpPS}, the coefficient space $M_{2(s-i)}$ is finite-dimensional, and the solution space of
such an equation is $SL_2(\Z)$-invariant.
\end{dictionary}

\begin{physics}
This is the form in which Schur indices of four-dimensional $\mathcal N=2$ theories, and characters of rational CFTs
(Mathur--Mukhi--Sen), are classified: a monic MLDE whose coefficients are a finite number of constants times monomials
in $E_4,E_6$.
\end{physics}

\begin{status}
Definitions. \lean{E46SerreClosed} is a \textbf{theorem} (\lean{e46SerreClosed}, Chapter~\ref{ch:L14}).
\end{status}

\subsection{The weight grading interface}

To make the coefficients of AK's relation modular forms of the \emph{right weight}, one needs the relation module to be
homogeneous for the total weight (modular weight of the coefficient plus square-bracket weight of the vector).

\begin{definition}[AK weight grading]\label{l16:def:akgrading}
Let $\lsixteenVA^{E} = \C[E_4,E_6]\otimes V$ and let its total-weight-$n$ subspace be spanned by $f\otimes v$ with $f$ of
modular weight $p$, $L[0]v = qv$, $p+q=n$. An \emph{AK weight grading} consists of: a direct-sum decomposition of $\C[E_4,E_6]$
into its weight spaces; a direct-sum decomposition of $\lsixteenVA^E$ into total-weight spaces; and the statement that the
relation module $O^E_q(V)$ (\lean{blockRelSpanE W}, the $E_4,E_6$-version of $\lsixteenOq$) is homogeneous.
\end{definition}
\noindent File \leanpath{VOA/\allowbreak Modularity/\allowbreak AKGrading.lean}.
\begin{leancode}
noncomputable def e46WeightSpace (k : ℤ) : Submodule ℂ E46 where
  carrier := {f | E46Weight k f}

noncomputable def akTotalWeightSpace (W : VertexOperatorAlgebra 𝒱) (n : ℤ) :
    Submodule ℂ (𝒱E 𝒱) :=
  Submodule.span ℂ {x | ∃ (p q : ℤ) (f : E46) (v : 𝒱), p + q = n ∧
    f ∈ e46WeightSpace p ∧ v ∈ sqWeightSpace W (q : ℂ) ∧ x = f ⊗ₜ[ℂ] v}

structure AKWeightGrading (W : VertexOperatorAlgebra 𝒱) where
  coeffDecomposition : DirectSum.Decomposition e46WeightSpace
  totalDecomposition : DirectSum.Decomposition (akTotalWeightSpace W)
  relationHomogeneous :
    letI := totalDecomposition
    DirectSum.SetLike.IsHomogeneous (akTotalWeightSpace W)
      ((blockRelSpanE W).restrictScalars ℂ)
\end{leancode}

\begin{dictionary}
\lean{𝒱E 𝒱} is \lean{↥E46 ⊗[ℂ] 𝒱} (file \leanpath{VOA/\allowbreak Modularity/\allowbreak BlockRelSpanE.lean});
\lean{sqWeightSpace W q} $=\ker(L[0]-q)$ (file \leanpath{VOA/\allowbreak Modularity/\allowbreak SquareBracketGrading.lean}).
\lean{DirectSum.}\allowbreak\lean{Decomposition A} is Mathlib's data of an internal direct-sum decomposition into the family $A$, with the
projections; \lean{IsHomogeneous} says a submodule is the sum of its intersections with the pieces. The structure is
\emph{data plus a proposition}, not a theorem. The first two fields hold for every VOA of CFT type: modular
forms of different weights are linearly independent, and $L[0]$ is diagonalisable on $V$ with eigenvalues in $\N$ (Zhu's
theorem that $(V,Y[\,],\tilde\omega)$ is again a VOA). The third field, \lean{relationHomogeneous}, is the real
VOA-side content: it says that Zhu's relation module is homogeneous for the \emph{total} weight. Arakawa
and Kawasetsu use a grading of this kind as a standard fact, but it is not a triviality, because the relation generators
$a[0]v$ are indexed by the \emph{ordinary} weight of $a\in V_{k_a}$ while the grading uses the square-bracket weight. So
the third field needs an argument (Zhu covariance of the square-bracket modes). \textbf{All three fields are now proved,
for every VOA of the library} (Theorem~\ref{l16:thm:akgrading} below); the earlier statement of this book that
\lean{AKWeightGrading} had no instance, not even for the trivial VOA, recorded a gap in the library, not a mathematical
obstruction, and is superseded.
\end{dictionary}

\begin{physics}
The weight grading is the statement that torus amplitudes are quasi-modular forms of definite weight, the weight of an
insertion being its scaling dimension on the cylinder. It is the bookkeeping that makes ``modular form of weight
$2(s-i)$'' meaningful.
\end{physics}

\begin{status}
Definition (structure with three fields). Inhabited for every VOA of the library (Theorem~\ref{l16:thm:akgrading}).
\end{status}

\begin{theorem}[The AK weight grading exists]\label{l16:thm:akgrading}
For every vertex operator algebra $V$ of the library (in particular of CFT type: $V_0=\C\lsixteenvac$ and positive grading),
\lean{AKWeightGrading V} is inhabited.
\end{theorem}
\noindent File \leanpath{VOA/\allowbreak Modularity/\allowbreak AKGradingTotal.lean} (section variable
\lean{(W : VertexOperatorAlgebra 𝒱)}).
\begin{leancode}
noncomputable def akWeightGrading : AKWeightGrading W where
  coeffDecomposition := e46Decomposition
  totalDecomposition := akTotalDecomposition W
  relationHomogeneous := by
    letI := akTotalDecomposition W
    rw [blockRelSpanE_eq_span]
    exact isHomogeneous_span _ (relHomE_homog W)
\end{leancode}

\textbf{Proof on paper.} (A) \emph{Coefficients.} $E_4$, $E_6$ and the constants are homogeneous, so the weight spaces
span $\C[E_4,E_6]$. They are independent: if $\sum_kf_k=0$ with $f_k$ modular of weight $k$, evaluate at
$\gamma_n z$ with $\gamma_n=\left(\begin{smallmatrix}1&0\\n&1\end{smallmatrix}\right)$ to get
$\sum_k(nz+1)^kf_k(z)=0$ for all $n$; the polynomial $\sum_kf_k(z)X^k$ has infinitely many roots, so every $f_k(z)=0$
(\lean{e46WeightSpace_isInternal}, file \leanpath{AKGradingInstance.lean}). (B) \emph{Total weight.} The tensors
$f\otimes v$ split into homogeneous pieces because $L[0]$ is diagonalisable with eigenvalues in $\N$; independence holds
because the total-weight space of weight $n$ lies in the $n$-eigenspace of $\mathrm{Wt}\otimes1+1\otimes L[0]$
(\lean{akTotalWeightSpace_isInternal}). (C) \emph{Relations.} Zhu covariance ($a[j]$ maps $V_{[K]}\otimes V_{[n]}$ to
$V_{[K+n-j-1]}$) shows that each generator $a[0]v$ or $a[-2]v+\sum(2k-1)G_{2k}a[2k-2]v$ is a sum of total-homogeneous
relations; the weight-zero component of $a$ is a multiple of $\lsixteenvac$ (CFT type), and $\lsixteenvac[j]=0$ for $j\neq-1$
(\lean{sqMode_vacuum_eq_zero}). So Zhu's relation module is the $\C$-span of homogeneous elements
(\lean{blockRelSpanE_eq_span}), hence homogeneous. \qed

\begin{status}
\texttt{\#print axioms VOA.akWeightGrading}: \lean{[propext, Classical.}\allowbreak\lean{choice, Quot.}\allowbreak\lean{sound]},
recorded in \texttt{Notes/\allowbreak release/\allowbreak WS-RELEASE2.md} (commit \texttt{34b74b85}); independently
audited as faithful (\texttt{Notes/\allowbreak Textbook/\allowbreak referee/\allowbreak RELEASE\_AUDIT\_QUEUE2.md}).
The aggregate \texttt{lake build} at the release commit is pending. Scope: ``every VOA'' means every
\lean{VertexOperatorAlgebra} in the library's sense, which builds in CFT type.
\end{status}

\begin{physics}
On the torus a field of cylinder scaling dimension $q$ (its $L[0]$-weight) contributes like a modular form of weight $q$,
and the Ward identities of Definition~\ref{l16:def:formalblock} relate insertions of different dimension with Eisenstein
series as coefficients. The theorem says these Ward identities never mix total weights: every relation among torus
one-point functions splits into relations each of definite total weight (modular weight of the coefficient plus scaling
dimension of the insertion). This is the dimensional analysis physicists take for granted when they write the
coefficient of $\vartheta^{(i)}$ in an order-$s$ MLDE as a modular form of weight $2(s-i)$; here it is checked for every
VOA of CFT type, with no finiteness assumption. By itself it produces no differential equation.
\end{physics}

\subsection{The conditional MLDE}

\begin{definition}[Analytic genus-one block; Zhu Def.~5.1.1]\label{l16:def:analyticblock}
A map $F:V\to\{\text{functions }\C\to\C\}$ is an \emph{analytic block} if it is linear in $v$, each $F(v)$ is holomorphic on
$\mathbb H$, and for homogeneous $a$ of weight $k_a\ge1$, all $v$ and all $z\in\mathbb H$: (i) $F(a[0]v) = 0$; (ii)
$F(a[-2]v)(z) + \sum_{m\ge0}(2m+1)G_{2m+2}(z)F(a[2m]v)(z) = 0$; (iii)
$F(L[-2]v)(z) = \frac1{2\pi i}F(v)'(z) + \sum_{m\ge0}G_{2m+2}(z)F(L[2m]v)(z)$, with $G_{2j}(z)$ the Eisenstein
\emph{functions} (the sums of Definition~\ref{l16:def:formalblock} evaluated at $q = e^{2\pi iz}$).
\end{definition}
\noindent File \leanpath{VOA/\allowbreak Modularity/\allowbreak AnalyticBlock.lean} (Chapter~\ref{ch:L15}).
\begin{leancode}
structure IsAnalyticBlock (W : VertexOperatorAlgebra 𝒱) (F : 𝒱 → ℂ → ℂ) : Prop where
  holo : ∀ v : 𝒱, DifferentiableOn ℂ (F v) (imGT 0)
  addv : ∀ (v w : 𝒱) (z : ℂ), F (v + w) z = F v z + F w z
  smulv : ∀ (c : ℂ) (v : 𝒱) (z : ℂ), F (c • v) z = c * F v z
  vanish : ∀ (ka : ℕ), 1 ≤ ka → ∀ a ∈ W.grading (ka : ℤ), ∀ v : 𝒱, F (vSqZ W a ka 0 v) = 0
  negTwo : ∀ (ka : ℕ), 1 ≤ ka → ∀ a ∈ W.grading (ka : ℤ), ∀ (v : 𝒱), ∀ z ∈ imGT 0,
    F (vSqZ W a ka (-2) v) z
      + ∑ᶠ m : ℕ, ((2 * (m + 1) - 1 : ℕ) : ℂ)
          * (zhuGC (m + 1) z * F (vSqZ W a ka (2 * (m + 1) - 2 : ℤ) v) z) = 0
  omega : ∀ (v : 𝒱), ∀ z ∈ imGT 0,
    F (Lsq W (-2) v) z
      = (2 * (Real.pi : ℂ) * Complex.I)⁻¹ * _root_.deriv (F v) z
        + ∑ᶠ m : ℕ, zhuGC (m + 1) z * F (Lsq W (2 * (m + 1) - 2 : ℤ) v) z
\end{leancode}

\begin{dictionary}
\lean{DifferentiableOn ℂ (F v) (imGT 0)}: holomorphic on $\mathbb H$. \lean{zhuGC j z} is the function $G_{2j}(z)$
(Chapter~\ref{ch:L14}); \lean{_root_.deriv} is the complex derivative in $z$, so
$(2\pi i)^{-1}F'$ is $\theta F$ in the variable $q = e^{2\pi iz}$. Exactly Definition~\ref{l16:def:analyticblock}; the
difference from Definition~\ref{l16:def:formalblock} is only that functions replace formal series. Note that \lean{vanish}
is required for all $z$ (the function is identically $0$), including off $\mathbb H$, which is harmless. The space of all
analytic blocks is \lean{analyticBlocks W}.
\end{dictionary}

\begin{physics}
The same Ward identities as Definition~\ref{l16:def:formalblock}, now for honest functions of the modular parameter: these
are the torus conformal blocks of the chiral algebra.
\end{physics}

\begin{status}
Definition; trace functions of modules with subexponential growth are instances (\lean{traceBlock_}\allowbreak\lean{isAnalyticBlock}).
\end{status}

\begin{theorem}[AK Theorem 5.1, conditional form]\label{l16:thm:mlde}
Let $V$ be Etingof-finite, $a\in V$ square-bracket primary of integer weight $w$, and assume an AK weight grading for $V$
and \lean{E46SerreClosed}. Then there are $s\ge1$ and $p_0,\dots,p_{s-1}$ such that for \emph{every} analytic genus-one
block $F$ the function $F(a)$ satisfies the monic MLDE of Definition~\ref{l16:def:mlde} of weight $w$ with these
coefficients. Both side hypotheses are theorems: \lean{AKWeightGrading W} is \lean{akWeightGrading W}
(Theorem~\ref{l16:thm:akgrading}) and \lean{E46SerreClosed} is \lean{e46SerreClosed};
\lean{IsSqPrimary.}\allowbreak\lean{exists_}\allowbreak\lean{monicMLDEOn_}\allowbreak\lean{all_analyticBlocks_}\allowbreak\lean{of_etingofFinite'}
discharges the second and Theorem~\ref{l16:thm:mlde-uncond} below discharges both.
\end{theorem}
\noindent File \leanpath{VOA/\allowbreak Modularity/\allowbreak AKMLDE.lean}.
\begin{leancode}
theorem IsSqPrimary.exists_monicMLDEOn_all_analyticBlocks_of_etingofFinite
    {W : VertexOperatorAlgebra 𝒱} {a : 𝒱} {wt : ℤ}
    (ha : IsSqPrimary W a (wt : ℂ)) (hgrade : AKWeightGrading W)
    (hclose : E46SerreClosed) (hzero : VacuumGradeZero W) (hE : EtingofFinite W) :
    ∃ (s : ℕ) (p : ℕ → E46), 1 ≤ s ∧
      ∀ (F : 𝒱 → ℂ → ℂ), IsAnalyticBlock W F → MonicMLDEOn wt (F a) s p := by
\end{leancode}

\begin{dictionary}
\lean{IsAnalyticBlock W F} (file \leanpath{VOA/\allowbreak Modularity/\allowbreak AnalyticBlock.lean},
Chapter~\ref{ch:L15}) says $F:V\to(\C\to\C)$ is linear in $v$, each $F(v)$ is holomorphic on $\mathbb H$, and Zhu's three
conditions hold with the Eisenstein \emph{functions}: $F(a[0]v)=0$, the $[-2]$-identity, and
$F(L[-2]v) = \frac1{2\pi i}F(v)' + \sum_mG_{2m+2}F(L[2m]v)$. The trace functions $v\mapsto Z_M(v;\tau)$ of irreducible modules with a
conformal weight and subexponential one-point functions are analytic blocks (\lean{traceBlock_}\allowbreak\lean{isAnalyticBlock}, file
\leanpath{VOA/\allowbreak Modularity/\allowbreak SlashAmbientObstruction.lean}, Chapter~\ref{ch:L15}), and the growth
hypothesis is supplied by Theorem~\ref{l16:thm:growth}; so at $a = \lsixteenvac$ the theorem would give one MLDE for all
such characters. (No library theorem composes the two; the composition is immediate.) \textbf{Quantifier order:} $\exists s\,\exists p$ before $\forall F$ --- one equation
for all blocks, exactly AK's uniformity. \textbf{Hypotheses:} \lean{hzero} free; \lean{hE} the Etingof hypothesis;
\lean{hclose} is the classical Ramanujan closure of $\C[E_4,E_6]$ under the Serre derivative, \textbf{proved}
(\lean{e46SerreClosed}) and discharged in the primed form; \lean{hgrade} is \textbf{proved} for every VOA
(\lean{akWeightGrading}, Theorem~\ref{l16:thm:akgrading}). Only \lean{hE} remains.
\textbf{Comparison with AK Theorem 5.1:} same conclusion (with ``quasi-lisse'' weakened to Etingof-finite and the
character replaced by any analytic block evaluated at a primary insertion, which is more general). The conclusion is
not vacuous in the sense of Section~\ref{l16:sec:weak}: the coefficient spaces are finite-dimensional.
\end{dictionary}

\textbf{Proof on paper.} (1) Proposition~5.2 over $\C[E_4,E_6]$
(\lean{exists_}\allowbreak\lean{finite_}\allowbreak\lean{gen_sup_}\allowbreak\lean{blockRelSpanE_}\allowbreak\lean{of_etingofFinite}): Zhu's $G_2$ never needs to appear, because the only
relation in which $G_2$ would occur is $a[-2]v + 3G_4a[2]v+\cdots + G_2(\ldots)a[0]v$, and the $a[0]v$ term is itself a
relation. (2) Noetherianity of $\C[E_4,E_6]$ gives a uniform monic relation
$L[-2]^sa - \sum c_iL[-2]^ia\in O^E_q(V)$ (\lean{exists_}\allowbreak\lean{uniform_}\allowbreak\lean{monic_}\allowbreak\lean{Lsq_relationE_}\allowbreak\lean{mem_of_}\allowbreak\lean{etingofFinite}); the weight
grading lets one project it onto total weight $w+2s$, which forces $c_i$ to have modular weight $2(s-i)$. (3) Every analytic
block evaluates this relation to $F(L[-2]^sa) = \sum c_iF(L[-2]^ia)$. (4) Using \lean{E46SerreClosed} and the
$\omega$-recursion, in which the $G_2L[0]$-term combines with the derivative into a Serre derivative, one shows by
induction (\lean{IsSqPrimary.}\allowbreak\lean{exists_}\allowbreak\lean{uniformSerreExpansionUpTo}) that
$F(L[-2]^ia) = \vartheta^i_{(w)}F(a) + \sum_{j<i}h_{ij}\vartheta^j_{(w)}F(a)$ with $h_{ij}$ modular of weight $2(i-j)$,
independently of $F$. (5) Substitute. \qed

\begin{physics}
This is the theorem physicists use: the Schur index, and more generally every torus one-point function of a primary,
satisfies an MLDE whose order and coefficients are fixed by the chiral algebra alone. In the library it holds for every
Etingof-finite chiral algebra (a stand-in for quasi-lisse; Etingof--Schedler is not formalised), for Zhu's abstract
one-point blocks, with an order that is not bounded; at the $(A_1,A_3)$ example $L_{-4/3}(\lsixteensl)$ it is
unconditional, but AK's order~2 is not proved.
\end{physics}

\begin{status}
Proved; \texttt{\#print axioms} of \lean{VOA.IsSqPrimary.}\allowbreak\lean{exists_}\allowbreak\lean{monicMLDEOn_}\allowbreak\lean{all_analyticBlocks_}\allowbreak\lean{of_etingofFinite}:
\lean{[propext, Classical.}\allowbreak\lean{choice, Quot.}\allowbreak\lean{sound]} (re-run). \textbf{Conditional} on \lean{hE}, \lean{hgrade},
\lean{hclose}; \lean{hclose} is discharged by \lean{e46SerreClosed} in \lean{IsSqPrimary.}\allowbreak\lean{exists_}\allowbreak\lean{monicMLDEOn_}\allowbreak\lean{all_analyticBlocks_}\allowbreak\lean{of_etingofFinite'}
(\leanpath{EisensteinSerreClosure.lean}); \lean{hgrade} is discharged by \lean{akWeightGrading} in
Theorem~\ref{l16:thm:mlde-uncond}, which leaves \lean{hE} as the only hypothesis.
\end{status}

\begin{theorem}[AK Theorem 5.1 under Etingof-finiteness; the $L_{-4/3}(\lsixteensl)$ instance]\label{l16:thm:mlde-uncond}
(i) Let $V$ be Etingof-finite and $a\in V$ square-bracket primary of integer weight $w$. Then there are $s\ge1$ and
coefficients $p_i\in\C[E_4,E_6]$, modular of weight $2(s-i)$, such that for every analytic genus-one block $F$ the
function $F(a)$ satisfies $\vartheta^{(s)}F(a)+\sum_{i<s}p_i\vartheta^{(i)}F(a)=0$ on $\mathbb H$ (iterated Serre
derivatives). (ii) For $V=L_{-4/3}(\lsixteensl)$ and $a=\lsixteenvac$ this holds with no hypothesis.
\end{theorem}
\noindent File \leanpath{VOA/\allowbreak Modularity/\allowbreak AKMLDEUnconditional.lean}.
\begin{leancode}
theorem IsSqPrimary.exists_monicMLDEOn_all_analyticBlocks_of_etingofFinite_uncond
    {W : VertexOperatorAlgebra 𝒱} {a : 𝒱} {wt : ℤ}
    (ha : IsSqPrimary W a (wt : ℂ)) (hE : EtingofFinite W) :
    ∃ (s : ℕ) (p : ℕ → E46), 1 ≤ s ∧
      ∀ (F : 𝒱 → ℂ → ℂ), IsAnalyticBlock W F → MonicMLDEOn wt (F a) s p :=
  ha.exists_monicMLDEOn_all_analyticBlocks_of_etingofFinite' (akWeightGrading W)
    (VOA.QuasiLisseSl2.vacuumGradeZero_of_vacuum_grade W) hE

theorem exists_monicMLDEOn_all_analyticBlocks_simpleSl2QL :
    ∃ (s : ℕ) (p : ℕ → E46), 1 ≤ s ∧
      ∀ F, IsAnalyticBlock simpleSl2QL F →
        MonicMLDEOn 0 (F simpleSl2QL.vacuum) s p :=
  IsSqPrimary.exists_monicMLDEOn_all_analyticBlocks_of_etingofFinite_uncond
    (wt := 0) (by exact_mod_cast vacuum_isSqPrimary simpleSl2QL) VOA.Sl2Pres.etingofFinite_simpleSl2QL'
\end{leancode}

\begin{dictionary}
\textbf{Read the name with care.} Despite the suffix \lean{_uncond}, (i) is \emph{conditional on \lean{EtingofFinite W}}
(finite zeroth Poisson homology, the conclusion of Etingof--Schedler); quasi-lisseness is not formalised. The order $s$
comes from Noetherianity and is \emph{existential and unbounded}. ``Analytic block'' means Zhu's abstract one-point
conditions (\lean{IsAnalyticBlock}; the zero function is one). Module traces are blocks by a separate lemma,
\lean{torusOnePoint_isAnalyticBlock}, whose growth input \lean{OnePointSubexp} is supplied under \lean{EtingofFinite} by
\lean{onePointSubexp_of_etingofFinite} (\leanpath{Modules/\allowbreak QuasiLisseFinite.lean}). (ii) is genuinely
unconditional, with $s$ unspecified: \textbf{AK's order~2 for $L_{-4/3}(\lsixteensl)$ is not proved}. That the vacuum
character of $L_{-4/3}(\lsixteensl)$ is such a block follows from \lean{torusOnePoint_isAnalyticBlock},
\lean{onePointSubexp_adjoint_of_etingofFinite} and the support lemma, but no single library theorem states ``the character
satisfies a monic MLDE''. Do not confuse (ii) with \lean{Sl2Pres.}\allowbreak\lean{character_}\allowbreak\lean{formalODE_}\allowbreak\lean{simpleSl2QL_}\allowbreak\lean{uncond},
the near-vacuous formal $q\frac{d}{dq}$-equation of Section~\ref{l16:sec:weak}.
\end{dictionary}

\begin{physics}
For a quasi-lisse chiral algebra (in four-dimensional terms, a theory whose Higgs branch is a finite union of symplectic
leaves) the characters of modules are expected to solve a modular linear differential equation fixed by the chiral
algebra alone, as for rational theories. Part (i) proves this for every torus one-point function of a primary, with
Etingof-finiteness standing in for quasi-lisseness; the order $s$ is only shown to exist. Part (ii) is the first
unconditional instance: $L_{-4/3}(\lsixteensl)$ is the chiral algebra of the $(A_1,A_3)$ Argyres--Douglas theory under the
SCFT/VOA correspondence of Beem et al., so its vacuum character, the Schur index of that theory, lies in the solution
space of one monic MLDE (once one composes with the trace-function lemmas listed above). The second-order equation of
Arakawa--Kawasetsu and the identification of the remaining solutions with module characters are not proved.
\end{physics}

\begin{status}
Proved, conditional on \lean{EtingofFinite} for (i), unconditional for (ii). \texttt{\#print axioms} of both:
\lean{[propext, Classical.}\allowbreak\lean{choice, Quot.}\allowbreak\lean{sound]}, recorded in
\texttt{Notes/\allowbreak release/\allowbreak WS-RELEASE2.md} (commit \texttt{34b74b85}); audited as stated, with the
caveats above, in \texttt{RELEASE\_AUDIT\_QUEUE2.md}. The aggregate \texttt{lake build} at the release commit is pending.
\end{status}

%======================================================================
\section{Zhu modularity and conformal blocks under Etingof-finiteness}\label{l16:sec:zhu}
%======================================================================

\begin{theorem}[Zhu's Theorem 5.3.2 under Etingof-finiteness]\label{l16:thm:blocks}
If $V$ is Etingof-finite, the space of analytic genus-one blocks (restricted to $\mathbb H$) is finite-dimensional.
\end{theorem}
\noindent File \leanpath{VOA/\allowbreak Modularity/\allowbreak AKZhuModularity.lean}.
\begin{leancode}
theorem finiteDimensional_analyticBlocks_of_etingofFinite (hzero : VacuumGradeZero W)
    (hE : EtingofFinite W) :
    FiniteDimensional ℂ (Submodule.map restrictImGT (analyticBlocks W)) := by
\end{leancode}

\begin{dictionary}
\lean{analyticBlocks W} is the subspace of all \lean{F} with \lean{IsAnalyticBlock W F}
(\leanpath{VOA/\allowbreak Modularity/\allowbreak BlockSpace.lean}); \lean{restrictImGT} extends by zero off $\mathbb H$
(\leanpath{VOA/\allowbreak Modularity/\allowbreak BlockRank.lean}), so two blocks agreeing on $\mathbb H$ are identified. Same as
Zhu's Theorem~5.3.2 with $C_2$ weakened. The proof avoids the two places where the $C_2$-route needs $C_2$-cofiniteness
(finite-dimensionality of Zhu's algebra $A(V)$ and finite generation of the $E$-relation module).
\end{dictionary}

\textbf{Proof on paper.} At a fixed $z_0\in\mathbb H$, evaluating the relations of $\lsixteenOq$ at $z_0$ gives numerical
relations; the spanning argument of Theorem~\ref{l16:thm:prop52} with numbers in place of $\mathcal G$ shows that $V$ is
spanned by a finite set modulo these, so a block is determined at $z_0$ by finitely many values; by the
$\omega$-recursion and holomorphy, the block is determined on $\mathbb H$ by finitely many initial data. \qed

\begin{physics}
Finitely many independent torus conformal blocks: a quasi-lisse chiral algebra on the torus has a finite-dimensional
space of ``partition functions with insertions'', even though the theory may have infinitely many non-ordinary sectors.
\end{physics}

\begin{status}
Proved, conditional on \lean{hE}; \lean{[propext, Classical.}\allowbreak\lean{choice, Quot.}\allowbreak\lean{sound]} (re-run).
\end{status}

\begin{theorem}[Zhu modularity under Etingof-finiteness, conditional]\label{l16:thm:zhuEt}
Let $V$ be Etingof-finite and $(M_i)_{i\in\iota}$ a finite family of irreducible modules. If the trace functions are
independent on each $T$-eigenvalue fibre (\lean{hfib}) and the restricted analytic blocks together with the $S$-slashes of
the traces span a space of dimension at most $|\iota|$ (\lean{hrank}), then there is an $SL_2(\Z)$-representation
$\rho$ with $Z_{M_i}(v;A\tau) = (c\tau+d)^n\sum_j\rho(A)_{ij}Z_{M_j}(v;\tau)$ for all $v$ of $L[0]$-weight $n$.
\end{theorem}
\noindent File \leanpath{VOA/\allowbreak Modularity/\allowbreak AKZhuModularity.lean}.
\begin{leancode}
theorem exists_zhuModularity_of_sqAnalyticRank_le_irr_etingofFinite
    (hE : EtingofFinite W)
    (hfib : FibreIndependent Mfam (fun i => conformalWeightOf (Mfam i)))
    (hrank : Module.rank ℂ
        (sqAnalyticSlashAmbient Mfam (fun i => conformalWeightOf (Mfam i)))
      ≤ (Fintype.card ι : Cardinal)) :
    ∃ D : ModularData W ι, ZhuModularityAnalytic D :=
\end{leancode}

\begin{dictionary}
\lean{Mfam : ι → IrrModule W} is the family; \lean{conformalWeightOf} picks its conformal weights; \lean{ModularData},
\lean{ZhuModularityAnalytic}, \lean{FibreIndependent}, \lean{sqAnalyticSlashAmbient} are the objects of Zhu's theorem in
Chapter~\ref{ch:L15} (\lean{ZhuModularityAnalytic D} is the displayed transformation law for every $A\in SL_2(\Z)$).
The statement is \lean{exists_}\allowbreak\lean{zhuModularity_}\allowbreak\lean{of_sqAnalyticRank_}\allowbreak\lean{le_irr_}\allowbreak\lean{c2} of Chapter~\ref{ch:L15} with \lean{C2Cofinite}
replaced by \lean{EtingofFinite}; a kernel-checked \lean{example} in the same file shows that the $C_2$ theorem, re-proved
through this one, has literally the same type. \textbf{Hypotheses:} \lean{hfib}, \lean{hrank} are open inputs, and for
the motivating non-lisse examples \lean{hrank} is expected to be \emph{false} (warning below).
The file's honesty block records that \lean{hrank} is discharged only at the trivial VOA, and that Etingof-finiteness
\emph{cannot} supply it on this route (it bounds $V/O_2^P$, not Zhu's algebra $A(V) = V/O(V)$, since $a_{(0)}b\notin O(V)$
in general).
\end{dictionary}

\textbf{Proof on paper.} Zhu's argument (Chapter~\ref{ch:L15}) uses $C_2$-cofiniteness at exactly one place on this route,
the growth of one-point functions; Theorem~\ref{l16:thm:growth} supplies it. \qed

\begin{physics}
The characters of the ordinary sectors transform into each other under the modular group --- the statement that the torus
partition function of the full CFT can be made modular invariant. For quasi-lisse theories AK only predict that characters
lie in an invariant solution space; whether the span of characters itself is invariant is a separate question, and here it is
reduced to \lean{hfib} and \lean{hrank}. For the Argyres--Douglas example the answer is expected to be no (warning below):
the single ordinary character, with $c=-6$, cannot be modular invariant by itself.
\end{physics}

\begin{status}
Proved; \lean{[propext, Classical.}\allowbreak\lean{choice, Quot.}\allowbreak\lean{sound]} (re-run). Conditional on \lean{hE}, \lean{hfib}, \lean{hrank}. Fully
discharged only at the trivial VOA (\lean{trivial_}\allowbreak\lean{exists_}\allowbreak\lean{zhuModularity_}\allowbreak\lean{etingofFinite}), which is lisse.
Expected \emph{not} to be instantiable at $L_{-4/3}(\lsixteensl)$ (warning below; paper argument). This is Zhu's theorem with
\lean{C2Cofinite} replaced by \lean{EtingofFinite} in one step; it is not a modularity theorem for quasi-lisse VOAs. The file also gives the
character-convergence corollary \lean{qSummable_}\allowbreak\lean{gradedDim_}\allowbreak\lean{of_etingofFinite} and a holomorphic $c=24$ corollary
\lean{char_SL2Z_}\allowbreak\lean{invariant_}\allowbreak\lean{of_sqAnalyticRank_}\allowbreak\lean{of_etingofFinite} (conditional on a rank bound \lean{hle}).
\end{status}

\begin{warning}[title={\texttt{hrank} is expected to fail at $L_{-4/3}(\mathfrak{sl}_2)$ (paper argument, release audit)}]
The following argument is due to the release audit (\texttt{Notes/\allowbreak Textbook/\allowbreak referee/\allowbreak
RELEASE\_AUDIT\_L16.md}, finding~5). It is argued on paper and \textbf{not machine-checked}.
\begin{enumerate}
\item At $k=-4/3=-2+2/3$ the admissible highest weights have $\lambda_1\in\{0,-2/3,-4/3\}$, and only $\lambda_1 = 0$ gives a
  finite-dimensional top space. So (Adamovi\'c--Milas; Arakawa) $L = L_{-4/3}(\lsixteensl)$ has exactly one ordinary
  irreducible module, $L$ itself.
\item \lean{FibreIndependent} forbids repeating a module, so a family containing the adjoint module is a singleton, and
  \lean{ZhuModularityAnalytic} would make $\chi_L(\tau)$ a weight-$0$ modular function for a character $\rho$ of
  $SL_2(\Z)$.
\item $c=-6$, so the leading power of $\chi_L$ is $q^{-c/24} = q^{1/4}$, which forces $\rho(T) = i$; hence $\rho = \varepsilon^6$
  with $\varepsilon$ the multiplier of $\eta$. Then $\chi_L/\eta^6$ would be a holomorphic modular form of level one and
  weight $-3$, nonzero at the cusp. No nonzero modular form of odd weight for $SL_2(\Z)$ exists ($-I$ acts by $(-1)^k$).
\end{enumerate}
Hence \lean{hrank} must fail for every family containing the adjoint module of $L_{-4/3}(\lsixteensl)$. This agrees with
what Arakawa and Kawasetsu actually prove: the \emph{solution space} of the MLDE is $SL_2(\Z)$-stable, not the span of
the ordinary characters. So for the motivating class of examples the hypotheses of Theorem~\ref{l16:thm:zhuEt} are not
merely open but expected to be false, and the theorem must not be described as ``modularity for quasi-lisse VOAs''.
\end{warning}

%======================================================================
\section{The example: the classical Poisson model of $L_{-4/3}(\mathfrak{sl}_2)$}\label{l16:sec:model}
%======================================================================

A theorem with hypothesis \lean{EtingofFinite W} is only interesting if the hypothesis holds for some $V$ that is
\emph{not} $C_2$-cofinite (otherwise Zhu's theorem already applies). The canonical example is the simple affine vertex
algebra $L_k(\lsixteensl)$ at an admissible level. Recall (Chapter~\ref{ch:L09}, where \lean{AdmissibleLevel} is defined)
that $k$ is \emph{admissible} for $\widehat{\lsixteensl}$ if $k+2 = p/q$ with $p,q$ coprime positive integers and $p\ge2$;
for $q\ge2$ these levels are non-integral and $L_k$ is not lisse. The library takes $(p,q) = (2,3)$: $\gcd(2,3)=1$,
$k = -2+\frac23 = -\frac43$, $k+2 = \frac23\ne0$ (non-critical, which is what \lean{hkQL} below records), and
$c = 3k/(k+2) = 3(-\frac43)/\frac23 = -6$. The library's proof of admissibility is
\lean{Admissible.}\allowbreak\lean{admissibleLevel_}\allowbreak\lean{neg_four_}\allowbreak\lean{thirds} (\leanpath{AdmissibleLevel.lean}, Chapter~\ref{ch:L09}). This section does the commutative algebra on the classical side, unconditionally; the
next two transport it to the vertex algebra.

\subsection{The symbol of the singular vector}

By Chapter~\ref{ch:L08} (Theorem~\ref{l08:thm:sing43}) $V^{-4/3}(\lsixteensl)$ has the singular vector
\begin{equation*}
s = 20\,e_{-3}\lsixteenvac - 12\,e_{-2}h_{-1}\lsixteenvac - 6\,h_{-2}e_{-1}\lsixteenvac + 36\,e_{-1}e_{-1}f_{-1}\lsixteenvac
+ 9\,e_{-1}h_{-1}h_{-1}\lsixteenvac
\end{equation*}
of weight $3$, and (Theorem~\ref{l08:thm:symbol}) its class $[s]$ in $R_{V^k} = V^k/C_2(V^k)$ (the image of $s$ under the quotient map, called its
\emph{symbol}) is $9[e](4[e][f]+[h]^2) = 9\,e\,\lsixteenCas$. The reason: a current mode $x_{-n}$ with $n\ge2$ acting on a state
$w$ is $(x_{-1}\lsixteenvac)_{(-n)}w$, and $u_{(-n)}w\in C_2(V)$ for $n\ge2$ (for $n=2$ by definition, for $n>2$ by
$u_{(-n)} = \frac1{n-1}(L_{-1}u)_{(-n+1)}$ and induction, Chapter~\ref{ch:L12}, Theorem~\ref{l12:thm:poisson}); so the first three terms of
$s$ have zero class, while $x_{-1}w = (x_{-1}\lsixteenvac)_{(-1)}w$ has class $[x]\cdot[w]$, giving
$[s] = 36[e]^2[f] + 9[e][h]^2$. Here $\lsixteenCas = h^2+4ef$ is the Casimir. In
$\operatorname{Sym}(\lsixteensl) = \C[e,h,f]$ with the Kirillov--Kostant bracket $\{e,f\} = h$, $\{h,e\} = 2e$,
$\{h,f\} = -2f$, the Casimir is Poisson-central, so $\{x, e\lsixteenCas\} = \{x,e\}\lsixteenCas$, and the Poisson ideal
generated by $e\lsixteenCas$ is
\begin{equation*}
\lsixteenCas\cdot\mathfrak m = (\lsixteenCas\,e,\ \lsixteenCas\,h,\ \lsixteenCas\,f),\qquad \mathfrak m = (e,h,f).
\end{equation*}

\begin{definition}[The classical model]\label{l16:def:model}
$\operatorname{Sym}(\lsixteensl) = \C[X_0,X_1,X_2]$ with $e = X_0$, $h = X_1$, $f = X_2$; $\lsixteenCas = h^2+4ef$; the three
derivations $\operatorname{ad}_e = -2e\,\partial_h + h\,\partial_f$, $\operatorname{ad}_h = 2e\,\partial_e - 2f\,\partial_f$,
$\operatorname{ad}_f = -h\,\partial_e + 2f\,\partial_h$ (the brackets $\{e,\cdot\}$, $\{h,\cdot\}$, $\{f,\cdot\}$); the ideal
$\lsixteenCas\cdot\mathfrak m$; and the model algebra $\lsixteenR = \operatorname{Sym}(\lsixteensl)/\lsixteenCas\cdot\mathfrak m$.
\end{definition}
\noindent File \leanpath{VOA/\allowbreak Examples/\allowbreak AffineSl2QuasiLisse.lean}.
\begin{leancode}
abbrev SymSl2 : Type := MvPolynomial (Fin 3) ℂ

def eSym : SymSl2 := X 0
def hSym : SymSl2 := X 1
def fSym : SymSl2 := X 2

def casSym : SymSl2 := hSym ^ 2 + 4 * eSym * fSym

def adOf (a : SymSl2) (i : Fin 3) (b : SymSl2) (j : Fin 3) : SymSl2 →ₗ[ℂ] SymSl2 :=
  (LinearMap.mulLeft ℂ a).comp (pderiv i).toLinearMap
    + (LinearMap.mulLeft ℂ b).comp (pderiv j).toLinearMap

def adE : SymSl2 →ₗ[ℂ] SymSl2 := adOf (-2 * eSym) 1 hSym 2
def adH : SymSl2 →ₗ[ℂ] SymSl2 := adOf (2 * eSym) 0 (-2 * fSym) 2
def adF : SymSl2 →ₗ[ℂ] SymSl2 := adOf (-hSym) 0 (2 * fSym) 1

def casIdeal : Ideal SymSl2 := Ideal.span (Set.range fun i : Fin 3 => casSym * X i)

def casSub : Submodule ℂ SymSl2 := Submodule.restrictScalars ℂ casIdeal

abbrev NilpAlg : Type := SymSl2 ⧸ casSub
\end{leancode}

\begin{dictionary}
\lean{MvPolynomial (Fin 3) ℂ} is $\C[X_0,X_1,X_2]$; \lean{pderiv i} is $\partial/\partial X_i$; \lean{adOf a i b j} is the
first-order operator $p\mapsto a\,\partial_ip + b\,\partial_jp$. So \lean{adE} $=-2e\partial_h + h\partial_f$, etc., exactly
Definition~\ref{l16:def:model}; \lean{adE_f}, \lean{adE_h}, \lean{adH_e}, \ldots\ (in
\leanpath{AffineSl2Presentation.lean}) check the nine values on generators against $\lsixteensl$. \textbf{Index warning}
(from the source): the polynomial variables are ordered $(e,h,f)$, whereas the affine colours of Chapter~\ref{ch:L08} are
$(e,f,h)$ = $(0,1,2)$; the two \lean{Fin 3}'s are different. \lean{casIdeal} is the ideal generated by
$\lsixteenCas X_i$, i.e.\ $\lsixteenCas\cdot\mathfrak m$ --- \textbf{not} the principal ideal $(\lsixteenCas)$, whose
generator has degree $2$. \lean{casSub} is the same set viewed as a vector subspace, and \lean{NilpAlg} is $\lsixteenR$ as a
vector space (no ring structure on the quotient is used). The library proves that \lean{casIdeal} is a Poisson ideal
(\lean{adE_casIdeal}, \lean{adH_casIdeal}, \lean{adF_casIdeal}) and that the Casimir is central (\lean{adE_cas}, \ldots).
\end{dictionary}

\begin{warning}[title={A retracted objection, recorded in the source}]
An earlier version of the source claimed that the hypothesis ``the symbol ideal dies'' was unsatisfiable because it would
force $R_V$ to be reduced; the error was to read \lean{casIdeal} as $(\lsixteenCas)$. As the docstring of
\lean{Sl2Presentation} now explains, $\lsixteenCas\notin\lsixteenCas\cdot\mathfrak m$, and indeed $R_{L_k}$ is
non-reduced: $[\lsixteenCas]\ne0$ but $[\lsixteenCas]^2 = 0$, i.e.\ the Sugawara vector, which is proportional to the
Casimir, is nilpotent in $R_V$. As a vector space $\lsixteenR = \C[\lsixteenNil]\oplus\C\lsixteenCas$, a one-dimensional
thickening of the coordinate ring of the nilpotent cone (not formalised).
\end{warning}

\begin{physics}
$e,h,f$ are the classical moment maps of the $SU(2)$ flavour symmetry. In the four-dimensional theory they are the
Higgs-branch moment-map operators; $\lsixteenCas = 0$ is the F-term (hyperk\"ahler moment map) relation cutting out
$\C^2/\Z_2$ as the nilpotent cone, $\{h^2+4ef = 0\}$. The nilpotence of $[\lsixteenCas]$ is the chiral-algebra statement
that the stress tensor is, classically, a nilpotent element: the Sugawara construction identifies $T$ with the Casimir of
the currents, and in Beem--Rastelli's dictionary this is why the 2d stress tensor is not a Higgs-branch chiral ring
element.
\end{physics}

\begin{status}
Definitions, plus the Poisson-ideal lemmas, all \lean{[propext, Classical.}\allowbreak\lean{choice, Quot.}\allowbreak\lean{sound]}.
\end{status}

\begin{theorem}[The classical quasi-lisse phenomenon, unconditionally]\label{l16:thm:model}
(i) $\operatorname{Sym}(\lsixteensl) = \lsixteenCas\cdot\mathfrak m + \{\lsixteensl,\operatorname{Sym}(\lsixteensl)\} +
\operatorname{span}\{1,h^2,ef\}$, where $\{\lsixteensl,\operatorname{Sym}\}$ is the sum of the ranges of the three
derivations. (ii) $\lsixteenR$ is infinite-dimensional. (iii)
$\operatorname{Sym}(\lsixteensl)/(\lsixteenCas\cdot\mathfrak m + \{\lsixteensl,\operatorname{Sym}\})$ is finite-dimensional; a
fortiori $\lsixteenHP(\lsixteenR)$ is.
\end{theorem}
\noindent File \leanpath{VOA/\allowbreak Examples/\allowbreak AffineSl2QuasiLisse.lean}.
\begin{leancode}
def adRange : Submodule ℂ SymSl2 :=
  LinearMap.range adE ⊔ LinearMap.range adH ⊔ LinearMap.range adF

def lowSpan : Submodule ℂ SymSl2 := Submodule.span ℂ {1, hSym ^ 2, eSym * fSym}

def bigSub : Submodule ℂ SymSl2 := casSub ⊔ adRange ⊔ lowSpan

theorem bigSub_eq_top : bigSub = ⊤ := by

theorem not_finiteDimensional_NilpAlg : ¬ FiniteDimensional ℂ NilpAlg := by

theorem finiteDimensional_quot_adRange :
    FiniteDimensional ℂ (SymSl2 ⧸ (casSub ⊔ adRange)) := by
\end{leancode}

\begin{dictionary}
\lean{LinearMap.range adE} is $\{e,\operatorname{Sym}\}$, etc.; \lean{adRange} is contained in the full bracket span, so (iii)
as stated (quotient by \lean{casSub ⊔ adRange}) is \emph{stronger} than finiteness of $\lsixteenHP(\lsixteenR)$, which the
library does not state separately for the model. Exactly Theorem~\ref{l16:thm:model}. Note that (ii) and (iii) together
are the whole point: an infinite-dimensional Poisson algebra with finite zeroth Poisson homology.
\end{dictionary}

\textbf{Proof on paper.} (i) It suffices to treat monomials $e^ah^cf^b$; induct on $a+b$. Since $\operatorname{ad}_h$ acts on
$e^ah^cf^b$ by $2(a-b)$, monomials with $a\ne b$ are brackets. If $a = b\ge1$, write $u = e^{a-1}h^cf^{a-1}$, so $e^ah^cf^a = u\cdot ef$. Since $4ef = \lsixteenCas - h^2$,
\begin{equation*}
e^ah^cf^a = \tfrac14\,u\,\lsixteenCas - \tfrac14\,u\,h^2 = \tfrac14\,\lsixteenCas\cdot u - \tfrac14\,e^{a-1}h^{c+2}f^{a-1}.
\end{equation*}
Unless $(a,c) = (1,0)$ (then $e^ah^cf^a = ef\in\operatorname{span}\{1,h^2,ef\}$), the monomial $u$ has positive degree, i.e.\
$u\in\mathfrak m$, so $\lsixteenCas\cdot u\in\lsixteenCas\cdot\mathfrak m$, and therefore
$e^ah^cf^a\equiv-\frac14e^{a-1}h^{c+2}f^{a-1}$ modulo $\lsixteenCas\cdot\mathfrak m$, a monomial with smaller $a+b$, to which
the induction hypothesis applies. Remain the powers $h^c$: $c = 0,2$ are in the span; $h =
\operatorname{ad}_e f$; for $c\ge3$, $\operatorname{ad}_e(fh^{c-1}) = h\cdot h^{c-1} + f\cdot(c-1)h^{c-2}(-2e) = h^c - 2(c-1)efh^{c-2}$, and by the
same computation with $u = h^{c-2}\in\mathfrak m$, $efh^{c-2} = \frac14\lsixteenCas h^{c-2} - \frac14h^c\equiv
-\frac14h^c$ modulo $\lsixteenCas\mathfrak m$, so $h^c - 2(c-1)(-\frac14)h^c =$ $(1+\frac{c-1}2)h^c$ is in the sum, and the coefficient is nonzero. (ii) The
algebra map $e\mapsto t$, $h,f\mapsto0$ to $\C[t]$ kills $\lsixteenCas$, hence $\lsixteenCas\mathfrak m$, and sends $e^n$ to
$t^n$; so the classes $[e^n]$ are linearly independent. (iii) By (i) the quotient is spanned by the images of
$1,h^2,ef$. \qed

\begin{physics}
(ii): the Higgs branch is not a point --- the moment map $e$ has arbitrarily high nonzero powers, the holomorphic functions
$\mu^n$ on $\C^2/\Z_2$. (iii): yet only three classical torus amplitudes survive. The leaves of $\lsixteenNil$: the Poisson bivector of
$\lsixteensl^*$ at a point $\xi$ has rank $\dim$ of the coadjoint orbit through $\xi$; it vanishes at $\xi=0$ (the origin is a
leaf by itself) and has rank $2$ on $\lsixteenNil\setminus\{0\}$, which is a single coadjoint orbit (the regular nilpotent
orbit, connected). So $\lsixteenNil$ has exactly two leaves, the situation of Theorem~\ref{l16:thm:ES}.
\end{physics}

\begin{status}
Proved, unconditional; the three theorems: \lean{[propext, Classical.}\allowbreak\lean{choice, Quot.}\allowbreak\lean{sound]} (re-run).
\end{status}

%======================================================================
\section{Transport to a vertex algebra: symbol presentations}\label{l16:sec:pres}
%======================================================================

\begin{definition}[An $\lsixteensl$ symbol presentation of $R_V$]\label{l16:def:pres}
A presentation for $V$ is a linear surjection $\sigma:\operatorname{Sym}(\lsixteensl)\to R_V$ that kills
$\lsixteenCas\cdot\mathfrak m$ and intertwines $\operatorname{ad}_x$ with $\{\sigma(x),\cdot\}$ for $x = e,h,f$. It is
\emph{exact} if $\ker\sigma\subseteq\lsixteenCas\cdot\mathfrak m$.
\end{definition}
\noindent File \leanpath{VOA/\allowbreak Examples/\allowbreak AffineSl2QuasiLisse.lean}.
\begin{leancode}
structure Sl2Presentation (W : VertexOperatorAlgebra 𝒱) where
  map : SymSl2 →ₗ[ℂ] A₂ W
  surjective : Function.Surjective map
  cas_vanishes : ∀ p ∈ casIdeal, map p = 0
  adE_compat : ∀ p, map (adE p) = A₂bracket W (map eSym) (map p)
  adH_compat : ∀ p, map (adH p) = A₂bracket W (map hSym) (map p)
  adF_compat : ∀ p, map (adF p) = A₂bracket W (map fSym) (map p)

def IsExact : Prop := ∀ p : SymSl2, S.map p = 0 → p ∈ casIdeal
\end{leancode}

\begin{dictionary}
Field by field, Definition~\ref{l16:def:pres}; \lean{IsExact} is defined inside \lean{namespace Sl2Presentation} with
\lean{S : Sl2Presentation W}. \textbf{Is it a disguised finiteness assumption?} No: none of the fields mentions
dimension, and the model $\lsixteenR$ itself (with its own bracket) satisfies all of them while being
infinite-dimensional. \textbf{Is it satisfiable for an interesting $V$?} Yes: it is constructed for $L_{-4/3}(\lsixteensl)$
(Theorem~\ref{l16:thm:presL}). Note that \lean{map} is only required to be linear, not multiplicative; the
constructed one is an algebra map.
\end{dictionary}

\begin{physics}
A presentation identifies the classical phase space of the chiral algebra as (a thickening of) a subvariety of
$\lsixteensl^*$ cut out by the null vector's symbol; exactness says no further classical relations are imposed.
\end{physics}

\begin{status}
Structure and predicate. Inhabited at $L_{-4/3}(\lsixteensl)$; \lean{IsExact} for that instance is not proved (but no longer
needed, Section~\ref{l16:sec:nonlisse}).
\end{status}

\begin{theorem}[Consequences of a presentation]\label{l16:thm:presCons}
Let $\sigma$ be a presentation for $V$. (i) $V$ is Etingof-finite (without $C_2$-cofiniteness). (ii) If $\sigma$ is exact,
$V$ is not $C_2$-cofinite, and $\sigma$ induces a linear isomorphism $\lsixteenR\cong R_V$.
\end{theorem}
\noindent File \leanpath{VOA/\allowbreak Examples/\allowbreak AffineSl2QuasiLisse.lean}.
\begin{leancode}
theorem etingofFinite (S : Sl2Presentation W) : EtingofFinite W := by

theorem not_C2Cofinite (hex : S.IsExact) : ¬ C2Cofinite W := by

def nilpEquiv (S : Sl2Presentation W) (hex : S.IsExact) : NilpAlg ≃ₗ[ℂ] A₂ W :=

theorem quasiLisseWitness (S : Sl2Presentation W) (hex : S.IsExact) :
    ¬ C2Cofinite W ∧ EtingofFinite W :=
\end{leancode}

\begin{dictionary}
All in \lean{namespace VOA.}\allowbreak\lean{QuasiLisseSl2.}\allowbreak\lean{Sl2Presentation}. \lean{≃ₗ[ℂ]} is a linear isomorphism. Exactly
Theorem~\ref{l16:thm:presCons}.
\end{dictionary}

\textbf{Proof on paper.} (i) $R_V = \sigma(\operatorname{Sym}) = \sigma(\lsixteenCas\mathfrak m + \{\lsixteensl,\operatorname{Sym}\}
+ \operatorname{span}\{1,h^2,ef\})$ by Theorem~\ref{l16:thm:model}(i); the first summand maps to $0$, the second into
$\{R_V,R_V\}$ by the intertwining, the third onto a space of dimension $\le3$. So $\dim\lsixteenHP(R_V)\le3$. (ii) Exactness
makes the induced map $\lsixteenR\to R_V$ injective; it is surjective; and $\lsixteenR$ is infinite-dimensional by
Theorem~\ref{l16:thm:model}(ii). \qed

\begin{physics}
Once the classical phase space is known to be (a thickening of) the nilpotent cone, both halves of ``quasi-lisse but not
lisse'' follow: a two-dimensional Higgs branch (not lisse) with finitely many leaves (finitely many classical amplitudes).
\end{physics}

\begin{status}
Proved; all \lean{[propext, Classical.}\allowbreak\lean{choice, Quot.}\allowbreak\lean{sound]} (re-run for \lean{etingofFinite}, \lean{not_C2Cofinite}). (i) is
conditional only on the presentation; (ii) also on exactness.
\end{status}

\subsection{The flagship algebra}

\begin{definition}[$V^{-4/3}(\lsixteensl)$ and $L_{-4/3}(\lsixteensl)$]\label{l16:def:flagship}
$V^{-4/3}(\lsixteensl)$ is the universal affine vertex operator algebra of $\lsixteensl$ at level $-4/3$, and
$L_{-4/3}(\lsixteensl)$ is its quotient by the maximal proper graded ideal.
\end{definition}
\noindent File \leanpath{VOA/\allowbreak Examples/\allowbreak AffineSl2QuasiLisse.lean}.
\begin{leancode}
theorem hkQL : (-4 / 3 : ℂ) + 2 ≠ 0 :=
  Admissible.AdmissibleLevel.add_two_ne_zero Admissible.admissibleLevel_neg_four_thirds

def affineSl2QL : VertexOperatorAlgebra (AffineVOA.WakModule (-4 / 3 : ℂ)) :=
  AffineVOA.affineSl2VOA_closed (-4 / 3 : ℂ) hkQL

def simpleSl2QL :
    VertexOperatorAlgebra (AffineVOA.WakModule (-4 / 3 : ℂ) ⧸ QuotVOA.maxGradedIdeal affineSl2QL) :=
  QuotVOA.maxQuotientVOA affineSl2QL
\end{leancode}

\begin{dictionary}
\lean{affineSl2VOA_}\allowbreak\lean{closed k hk} is the universal affine VOA $V^k(\lsixteensl)$ on the vacuum module \lean{WakModule k}
(despite its name, not a Wakimoto module), Chapter~\ref{ch:L08}, Definition~\ref{l08:def:Vk}; \lean{hkQL} is the
non-criticality $k+2\ne0$, obtained from admissibility of $-4/3$. \lean{maxQuotientVOA} is the quotient by the maximal proper
graded ideal (Chapter~\ref{ch:L04}). So \lean{simpleSl2QL} is $L_{-4/3}(\lsixteensl)$. It is definitionally the algebra
\lean{LQ43.L43} of Chapter~\ref{ch:L09} (the two differ only in the proof of non-criticality, and proofs are
irrelevant; \lean{affineSl2QL_}\allowbreak\lean{eq_voaSl2Adm} is \lean{rfl}), which is proved simple, nonzero, and a proper quotient of
$V^{-4/3}$, all without \lean{VOA.pbw}. So the example is the intended object and is not the zero algebra.
\end{dictionary}

\begin{physics}
$L_{-4/3}(\lsixteensl)$, $c=-6$: the chiral algebra of the $(A_1,D_3)\cong(A_1,A_3)$ Argyres--Douglas theory (Beem et al.,
Creutzig), with $c_{4d} = 1/2$ and Higgs branch $\C^2/\Z_2$. It is non-unitary and non-rational: in category $\mathcal O$ it has
three irreducible highest-weight modules (admissible weights $0,-2/3,-4/3$), of which only the vacuum module has
finite-dimensional conformal-weight spaces (standard admissible-level representation theory, Adamovi\'c--Milas; not
formalised).
\end{physics}

\begin{status}
Definitions; \lean{simpleSl2QL} nonzero by \lean{Sl2Pres.}\allowbreak\lean{vacuum_}\allowbreak\lean{simpleSl2QL_}\allowbreak\lean{ne_zero'} (\lean{pbw}-free).
\end{status}

\begin{theorem}[The presentation of $R_{L_{-4/3}}$ is constructed]\label{l16:thm:presL}
(i) For every non-critical $k$, the algebra map $\operatorname{Sym}(\lsixteensl)\to R_{V^k(\lsixteensl)}$, $x\mapsto[x_{-1}\lsixteenvac]$, is
surjective. (ii) In $R_{L_{-4/3}}$, $[e]\cdot\lsixteenCas = 0$, and the whole ideal $\lsixteenCas\cdot\mathfrak m$ maps to $0$. (iii)
These maps form an $\lsixteensl$ symbol presentation of $L_{-4/3}(\lsixteensl)$. (iv) Hence $L_{-4/3}(\lsixteensl)$ is
Etingof-finite.
\end{theorem}
\noindent File \leanpath{VOA/\allowbreak Examples/\allowbreak AffineSl2Presentation.lean}.
\begin{leancode}
def symbolAlg : SymSl2 →ₐ[ℂ] A₂ (affineSl2VOA_closed k hk) :=
  MvPolynomial.aeval ![eCl k hk, hCl k hk, fCl k hk]

theorem symbolAlg_surjective : Function.Surjective (symbolAlg k hk) := by

def symbolAlgL : SymSl2 →ₐ[ℂ] A₂ simpleSl2QL := MvPolynomial.aeval ![eL, hL, fL]

theorem eL_mul_casL : eL * casL = 0 := by

theorem cas_vanishes_L : ∀ p ∈ casIdeal, symbolAlgL p = 0 := by

def sl2PresentationSimpleSl2QL : Sl2Presentation simpleSl2QL where

theorem etingofFinite_simpleSl2QL' : EtingofFinite simpleSl2QL :=
  sl2PresentationSimpleSl2QL.etingofFinite
\end{leancode}

\begin{dictionary}
\lean{eCl k hk}, \lean{hCl}, \lean{fCl} are the classes $[e_{-1}\lsixteenvac]$, $[h_{-1}\lsixteenvac]$, $[f_{-1}\lsixteenvac]$ in
$R_{V^k}$ (\leanpath{AffineSl2Symbol.lean}, Chapter~\ref{ch:L08}); \lean{MvPolynomial.}\allowbreak\lean{aeval ![x,y,z]} is the algebra map
sending $X_0,X_1,X_2$ to $x,y,z$ (note the order $e,h,f$). \lean{eL}, \lean{hL}, \lean{fL} are their images in
$R_{L_{-4/3}}$ under the map of Theorem~\ref{l16:thm:quot}, and \lean{casL} $=4\,eL\,fL + hL^2$. The structure
\lean{sl2PresentationSimpleSl2QL} has \lean{map := symbolAlgL.}\allowbreak\lean{toLinearMap}, \lean{surjective := symbolAlgL_}\allowbreak\lean{surjective},
\lean{cas_vanishes := cas_}\allowbreak\lean{vanishes_}\allowbreak\lean{L} and the three compatibilities \lean{adE_compat_L}, \ldots. Exactly
Theorem~\ref{l16:thm:presL}. (i) is only the surjective half of $R_{V^k} = \operatorname{Sym}(\lsixteensl)$; injectivity is not
used.
\end{dictionary}

\textbf{Proof on paper.} (i) $V^k$ is spanned by word states $J^{a_1}_{-n_1-1}\cdots J^{a_r}_{-n_r-1}\lsixteenvac$
(\lean{wordStateA_span}, Chapter~\ref{ch:L08}). If some mode is $\le-2$ the word lies in $C_2$; otherwise it is the iterated
product $[J^{a_1}]\cdots[J^{a_r}]$ in $R_{V^k}$, because $x_{-1}v = (x_{-1}\lsixteenvac)_{(-1)}v$. (ii) $s$ lies in the maximal ideal
of $V^{-4/3}$ (Chapter~\ref{ch:L09}: \lean{LQ43.mkQ_}\allowbreak\lean{sl2Singular_}\allowbreak\lean{eq_zero}, via Li's theorem that the maximal ideal is the radical of
the invariant form), so $0 = [s] = 9[e]\lsixteenCas$ in $R_{L}$. Then, $\lsixteenCas$ being central,
$[h]\lsixteenCas = -\{[f],[e]\lsixteenCas\} = 0$ and $2[f]\lsixteenCas = \{[f],[h]\lsixteenCas\} = 0$. (iii) The
compatibilities hold on generators by the nine brackets of currents ($\{[x],[y]\} = [[x,y]]$) and propagate to all
polynomials by the Leibniz rule. (iv) Theorem~\ref{l16:thm:presCons}(i). \qed

\begin{physics}
The null vector at level $-4/3$ imposes, classically, $e\cdot\lsixteenCas = 0$; the $SU(2)$ symmetry then forces
$\lsixteenCas\cdot\mathfrak m = 0$. That is enough to leave only three classical torus amplitudes. In four-dimensional terms:
the Higgs-branch relation is implied by the chiral algebra's null state.
\end{physics}

\begin{status}
Proved, unconditional, \lean{pbw}-free; \lean{VOA.Sl2Pres.}\allowbreak\lean{symbolAlg_}\allowbreak\lean{surjective}, \lean{VOA.Sl2Pres.}\allowbreak\lean{eL_mul_}\allowbreak\lean{casL},
\lean{VOA.Sl2Pres.}\allowbreak\lean{sl2PresentationSimpleSl2QL}, \lean{VOA.Sl2Pres.}\allowbreak\lean{etingofFinite_}\allowbreak\lean{simpleSl2QL'}: \lean{[propext,
Classical.}\allowbreak\lean{choice, Quot.}\allowbreak\lean{sound]} (re-run).
\end{status}

\begin{definition}[The exactness citation]\label{l16:def:exact}
\lean{Sl2ExactSymbol} is the proposition that the constructed presentation of $R_{L_{-4/3}}$ is exact, i.e.\ that
$\ker(\operatorname{Sym}(\lsixteensl)\to R_{L_{-4/3}})\subseteq\lsixteenCas\cdot\mathfrak m$.
\end{definition}
\noindent File \leanpath{VOA/\allowbreak Examples/\allowbreak AffineSl2Presentation.lean}.
\begin{leancode}
def Sl2ExactSymbol : Prop := sl2PresentationSimpleSl2QL.IsExact

theorem not_C2Cofinite_simpleSl2QL_of_exact (hex : Sl2ExactSymbol) : ¬ C2Cofinite simpleSl2QL :=
  sl2PresentationSimpleSl2QL.not_C2Cofinite hex
\end{leancode}

\begin{dictionary}
\lean{Sl2ExactSymbol} says that the kernel of the constructed symbol map is contained in $\lsixteenCas\cdot\mathfrak m$, i.e.\
$R_{L_{-4/3}}\cong\operatorname{Sym}(\lsixteensl)/\lsixteenCas\cdot\mathfrak m$. It is a named, \textbf{unproved}
proposition; the library never assumes it globally, only as an explicit hypothesis.
\end{dictionary}

\begin{warning}[title={The citation attached to \texttt{Sl2ExactSymbol} names a weaker theorem (documentation only)}]
The docstring calls \lean{Sl2ExactSymbol} ``Arakawa's associated-variety theorem for $L_{-4/3}(\lsixteensl)$, in the form this
repository consumes'' and says that it is ``a lower bound on $A_2$'' supplied by $X_{L_k(\lsixteensl)} = \lsixteenNil$. Arakawa's
theorem is a statement about the \emph{reduced} variety: it gives $\sqrt{\ker\sigma} = (\lsixteenCas)$, not
$\ker\sigma = \lsixteenCas\cdot\mathfrak m$. The scheme-theoretic statement \lean{Sl2ExactSymbol} is true, but its proof
(on paper) uses different inputs: (a) the maximal ideal of $V^{-4/3}$ is generated by $s$ (a standard fact of admissible-level
$\widehat{\lsixteensl}$ representation theory, going back to Kac--Wakimoto and Malikov--Feigin--Fuchs), so by PBW it equals $U(\hat{\mathfrak g}_{<0})\cdot\operatorname{span}\{s,f_0s,f_0^2s\}$; (b) in
$R_{V^k} = \operatorname{Sym}(\lsixteensl)$ (which needs PBW injectivity) the image of this space is the ideal generated by the
$\lsixteensl$-orbit of $e\lsixteenCas$, which is $\lsixteenCas\cdot\mathfrak m$. Since non-lisseness is now proved without it
(Section~\ref{l16:sec:nonlisse}), the hypothesis has become \textbf{redundant} wherever it occurs (checked in Lean against the
library; the check is not part of this release).
\end{warning}

\begin{physics}
Exactness would say that the chiral algebra's classical phase space is \emph{exactly} the thickened nilpotent cone, with no
further nilpotent relations --- a scheme-theoretic refinement of ``associated variety = Higgs branch''.
\end{physics}

\begin{status}
\lean{Sl2ExactSymbol}: unproved named proposition. \lean{not_C2Cofinite_}\allowbreak\lean{simpleSl2QL_}\allowbreak\lean{of_exact}: proved from it,
\lean{[propext, Classical.}\allowbreak\lean{choice, Quot.}\allowbreak\lean{sound]}; superseded by Theorem~\ref{l16:thm:nonlisse}.
\end{status}

%======================================================================
\section{$L_{-4/3}(\mathfrak{sl}_2)$ is not lisse}\label{l16:sec:nonlisse}
%======================================================================

The file \leanpath{AffineSl2NonLisse.lean} proves the lower bound on $R_{L_{-4/3}}$ that exactness was meant to give, by a
direct and much weaker argument: only the tower $[e^r]$ is shown to survive.

\begin{theorem}[Non-lisseness of $L_{-4/3}(\lsixteensl)$]\label{l16:thm:nonlisse}
Let $E_r = e_{-1}^r\lsixteenvac$ and $F_r = f_{-1}^r\lsixteenvac$ in $V = V^{-4/3}(\lsixteensl)$, and $\langle\cdot,\cdot\rangle$ the
canonical invariant bilinear form. (i) $\langle F_r, x\rangle = 0$ for all $x\in C_2(V)$. (ii) $\langle F_r,E_r\rangle =
(-1)^r\,r!\,\prod_{j<r}(k-j)\ne0$ at $k=-4/3$. (iii) The class of $E_r$ in $R_{L}$, $L = L_{-4/3}(\lsixteensl)$, is nonzero,
and these classes are linearly independent. (iv) $L$ is not $C_2$-cofinite. (v) $L$ is not $C_2$-cofinite and is
Etingof-finite.
\end{theorem}
\noindent File \leanpath{VOA/\allowbreak Examples/\allowbreak AffineSl2NonLisse.lean}.
\begin{leancode}
theorem form_fHarmonic_O₂_zero (r : ℕ) {x : WakModule kk} (hx : x ∈ O₂ W43) :
    LQ43.formAdm.form (fHarmonic kk r) x = 0 := by

theorem form_fHarmonic_eHarmonic_ne_zero (r : ℕ) :
    LQ43.formAdm.form (fHarmonic kk r) (eHarmonic kk r) ≠ 0 := by

theorem eHarmonicL_not_mem_O₂ (r : ℕ) : eHarmonicL r ∉ O₂ L43 := by

theorem eHarmonicL_C₂_linearIndependent :
    LinearIndependent ℂ (fun r : ℕ => mk₂ L43 (eHarmonicL r)) := by

theorem not_C2Cofinite_simpleSl2QL : ¬ C2Cofinite L43 := by

theorem strictEtingofFinite_simpleSl2QL :
    ¬ C2Cofinite L43 ∧ EtingofFinite L43 :=
  ⟨not_C2Cofinite_simpleSl2QL, Sl2Pres.etingofFinite_simpleSl2QL'⟩
\end{leancode}

\begin{dictionary}
\lean{kk}, \lean{W43}, \lean{I43}, \lean{L43} are \emph{local notations} of the file for $-4/3$, \lean{affineSl2QL},
\lean{QuotVOA.}\allowbreak\lean{maxGradedIdeal affineSl2QL} and \lean{simpleSl2QL}; so \lean{VOA.Sl2NonLisse.}\allowbreak\lean{not_C2Cofinite_}\allowbreak\lean{simpleSl2QL} is literally
\lean{¬ C2Cofinite simpleSl2QL}. \lean{eHarmonic k r}, \lean{fHarmonic k r} are $E_r$, $F_r$
(\leanpath{AffineSl2HarmonicNorm.lean}), and \lean{eHarmonicL r} is the image of $E_r$ in $L$. \lean{LQ43.formAdm} is the
canonical invariant form on $V^{-4/3}$ of Chapter~\ref{ch:L09}. (ii) is \lean{form_fHarmonic_}\allowbreak\lean{eHarmonic}
(\leanpath{AffineSl2HarmonicNorm.lean}) evaluated at $k=-4/3$: $k-j\ne0$ because $k$ is not a non-negative integer. Exactly
Theorem~\ref{l16:thm:nonlisse}; (iv) is unconditional. \textbf{Name clash:} the full name is
\lean{VOA.Sl2NonLisse.}\allowbreak\lean{not_C2Cofinite_}\allowbreak\lean{simpleSl2QL}; the declaration \lean{VOA.QuasiLisseSl2.}\allowbreak\lean{not_C2Cofinite_}\allowbreak\lean{simpleSl2QL}
(\leanpath{AKSl2Payoff.lean}) has the same short name. At the snapshot it took a presentation \lean{S} and
\lean{hex : S.IsExact}; a Lean check (not part of this release) showed the two are equal as proofs, so the hypotheses were redundant, and the
library has since restated it without hypotheses, as \lean{¬ C2Cofinite simpleSl2QL} proved by
\lean{Sl2NonLisse.not_C2Cofinite_simpleSl2QL}. The two twins now have identical types.
\end{dictionary}

\textbf{Proof on paper.} Write $h_0$-charge for the eigenvalue of $h_0$ ($e$ has charge $2$, $f$ has $-2$, $h$ has $0$). (i) The
form is invariant, so $\langle h_0x,y\rangle = -\langle x,h_0y\rangle$ and vectors of different charge $\gamma,\gamma'$ with
$\gamma\ne-\gamma'$ are orthogonal; also different conformal weights are orthogonal. $F_r$ has weight $r$ and charge $-2r$. A
homogeneous generator $u_{(-2)}v$ of $C_2(V)$ of weight $r$ has $\operatorname{wt}u + \operatorname{wt}v = r-1$; writing $u,v$ in
the word basis, each current contributes charge $\le2$ and weight $\ge1$, so the charge is $\le2(r-1)<2r$. Hence
$\langle F_r,u_{(-2)}v\rangle = 0$. (ii) Move $f_{-1}$ across with invariance ($f_{-1}^\dagger\sim e_1$) and use
$e_1E_{r} = $ an explicit multiple of $E_{r-1}$: one gets the recursion $\langle F_{r+1},E_{r+1}\rangle =
-(r+1)(k-r)\langle F_r,E_r\rangle$. (iii) If $\bar E_r\in C_2(L)$, lift: $C_2(L)$ is the image of $C_2(V)$, so $E_r = x+n$ with $x\in
C_2(V)$ and $n$ in the maximal ideal, which is the (right) radical $\{y : \langle x,y\rangle = 0\ \forall x\}$ of the form (Li's theorem: for a VOA of CFT type with
$L_1V_1=0$ the maximal proper ideal equals the radical of the canonical invariant form; Chapter~\ref{ch:L09},
\lean{LiRadical.}\allowbreak\lean{rightRadical_}\allowbreak\lean{eq_maxGradedIdeal}). Pairing
with $F_r$ gives $\langle F_r,E_r\rangle = 0$, contradicting (ii). The classes have distinct weights, so are independent. (iv)
An infinite independent family. (v) Add Theorem~\ref{l16:thm:presL}. \qed

\begin{physics}
$E_r = (J^+_{-1})^r\lsixteenvac$ is the chiral-algebra image of the Higgs-branch operator $\mu^r$, the $r$-th power of the
moment map's highest component, of $SU(2)_R$-spin $r$; and $F_r$ is its conjugate. Their two-point function is nonzero, so these
states are not null; and they are not ``derivatives'' (in $C_2$) because an excited mode costs energy without supplying charge.
The infinite tower of Higgs-branch chiral ring elements is what makes the chiral algebra non-lisse.
\end{physics}

\begin{status}
Proved, unconditional, \lean{pbw}-free: \texttt{\#print axioms} of \lean{VOA.Sl2NonLisse.}\allowbreak\lean{eHarmonicL_}\allowbreak\lean{not_mem_}\allowbreak\lean{O₂},
\lean{VOA.Sl2NonLisse.}\allowbreak\lean{not_C2Cofinite_}\allowbreak\lean{simpleSl2QL}, \lean{VOA.Sl2NonLisse.}\allowbreak\lean{strictEtingofFinite_}\allowbreak\lean{simpleSl2QL}:
\lean{[propext, Classical.}\allowbreak\lean{choice, Quot.}\allowbreak\lean{sound]} (re-run for this chapter).
\end{status}

\begin{warning}[title={Stale ``non-lisseness is not proved'' documentation (documentation only)}]
The module docstrings of \leanpath{AffineSl2QuasiLisse.lean} (``\lean{¬ C2Cofinite (simpleSl2QL)} is still NOT proved''),
\leanpath{AffineSl2Presentation.lean} (``\lean{L₋₄⁄₃(𝔰𝔩₂)} may NOT be described as `proved non-lisse'\,''),
\leanpath{AKSl2Payoff.lean} (``\lean{¬ C2Cofinite (L_{-4/3}(𝔰𝔩₂))} is NOT proved''), \leanpath{AKSl2Instance.lean} (item~1,
``NOT a non-lisse instance''), and the docstrings of \lean{VOA.Sl2Pres.}\allowbreak \lean{not_C2Cofinite_}\allowbreak\lean{simpleSl2QL_}\allowbreak\lean{of_exact},
\lean{VOA.QuasiLisseSl2.}\allowbreak\lean{not_C2Cofinite_}\allowbreak\lean{simpleSl2QL}, \lean{character_}\allowbreak\lean{formalODE_}\allowbreak\lean{of_non_}\allowbreak\lean{lisse_}\allowbreak\lean{simpleSl2QL} and
\lean{character_}\allowbreak\lean{formalODE_}\allowbreak\lean{adjoint_}\allowbreak\lean{of_non_}\allowbreak\lean{lisse_}\allowbreak\lean{simpleSl2QL} all said that non-lisseness is conditional on
\lean{Sl2ExactSymbol}. Already at the snapshot it was the unconditional Theorem~\ref{l16:thm:nonlisse}. \emph{Update
(release audit, 2026-10-27):} these declaration docstrings and module headers have been corrected (by a later revision of
the library); what remains in \leanpath{AffineSl2Presentation.lean} is the correct statement that
\lean{Sl2Presentation.IsExact} (and \lean{Sl2ExactSymbol}) is not proved.
\end{warning}

\begin{theorem}[$R_{L_{-4/3}}\ne0$, two proofs]\label{l16:thm:nontriv}
$R_{L_{-4/3}}$ is a nonzero vector space.
\end{theorem}
\noindent Files \leanpath{VOA/\allowbreak Examples/\allowbreak AffineSl2QuasiLisse.lean},
\leanpath{VOA/\allowbreak Examples/\allowbreak AffineSl2Presentation.lean}.
\begin{leancode}
theorem A₂_simpleSl2QL_nontrivial : Nontrivial (A₂ simpleSl2QL) :=
  A₂_nontrivial simpleSl2QL vacuum_simpleSl2QL_ne_zero

theorem A₂_simpleSl2QL_nontrivial' : Nontrivial (A₂ simpleSl2QL) :=
  A₂_nontrivial simpleSl2QL vacuum_simpleSl2QL_ne_zero'
\end{leancode}

\begin{dictionary}
\lean{Nontrivial X}: $X$ has two distinct elements. The two proofs differ only in how the nonzero vacuum is obtained: the
first (namespace \lean{QuasiLisseSl2}) uses \lean{AffineVOA.}\allowbreak\lean{vacuumA_}\allowbreak\lean{ne_zero}, which goes through
\lean{VOA.pbw} (an axiom at the snapshot, a theorem now); the second (namespace \lean{Sl2Pres}) uses the Wakimoto realisation and
does not use \lean{VOA.pbw}. Since \lean{VOA.pbw} is now a theorem, the difference has no axiom significance: both
report Lean's three standard axioms only. Both are now
implied by Theorem~\ref{l16:thm:nonlisse} (an infinite-dimensional space is nonzero).
\end{dictionary}

\textbf{Proof on paper.} $C_2(V)\subseteq\bigoplus_{n\ne0}V_n$ since $a_{(-2)}b$ has weight $\operatorname{wt}a+\operatorname{wt}b+1\ge1$;
so $[\lsixteenvac]\ne0$ when $\lsixteenvac\ne0$. \qed

\begin{physics}
The identity operator is not a derivative: the classical phase space is non-empty.
\end{physics}

\begin{status}
\lean{VOA.QuasiLisseSl2.}\allowbreak\lean{A₂_simpleSl2QL_}\allowbreak\lean{nontrivial}: \lean{[propext, Classical.}\allowbreak\lean{choice, Quot.}\allowbreak\lean{sound, VOA.}\allowbreak\lean{pbw]} at the snapshot;
now \lean{[propext, Classical.choice, Quot.sound]} (re-probed at a later commit, after \lean{VOA.pbw} became a theorem, not re-run for this edition; confirmed by the release audit). Its
proof passes through the theorem \lean{VOA.pbw}, which is no longer an axiom, so this carries no axiom significance.
\lean{A₂_simpleSl2QL_}\allowbreak\lean{nontrivial'}: standard three only.
\end{status}

%======================================================================
\section{The payoff at $L_{-4/3}(\mathfrak{sl}_2)$}\label{l16:sec:payoff}
%======================================================================

\begin{theorem}[\texttt{VacuumGradeZero} is free]\label{l16:thm:vgz}
For every library VOA, $V_0\subseteq\C\lsixteenvac$.
\end{theorem}
\noindent File \leanpath{VOA/\allowbreak Modularity/\allowbreak AKSl2Payoff.lean}.
\begin{leancode}
theorem vacuumGradeZero_of_vacuum_grade (W : VertexOperatorAlgebra 𝒱) : VacuumGradeZero W :=
  le_of_eq W.vacuum_grade
\end{leancode}

\begin{dictionary}
\lean{W.vacuum_}\allowbreak\lean{grade : W.}\allowbreak\lean{grading 0 = ℂ ∙ W.}\allowbreak\lean{vacuum} is a structure field of \lean{VertexOperatorAlgebra} (CFT type); \lean{le_of_eq}
weakens $=$ to $\le$. So the hypothesis \lean{hzero} carried by every theorem of Sections~\ref{l16:sec:prop52}--\ref{l16:sec:zhu}
is redundant in the strict sense (always satisfied); it is harmless.
\end{dictionary}

\begin{physics}
The vacuum is the unique state of zero energy.
\end{physics}

\begin{status}
Proved, unconditional.
\end{status}

\begin{theorem}[Generic payoff from a presentation]\label{l16:thm:payoffGen}
If $V$ has an $\lsixteensl$ symbol presentation, then the character of every ordinary module with a conformal weight satisfies
a formal equation~\eqref{l16:eq:formalODE}; if the presentation is exact, $V$ is moreover not $C_2$-cofinite.
\end{theorem}
\noindent File \leanpath{VOA/\allowbreak Modularity/\allowbreak AKSl2Payoff.lean}.
\begin{leancode}
theorem character_formalODE_of_sl2Presentation (W : VertexOperatorAlgebra 𝒱)
    (P : Sl2Presentation W) (N : OrdVModule W M) {hw : ℂ} (hcw : HasConformalWeight N hw) :
    ∃ (s : ℕ) (g : ℕ → ℕ → ℂ), 1 ≤ s ∧ (∀ i, Subexponential (fun k => ‖g i k‖)) ∧
      ∀ e : ℂ, (qDeriv^[s] (torusPartition N)) e
        + ∑ i ∈ Finset.range s, qConv (g i) (qDeriv^[i] (torusPartition N)) e = 0 :=

theorem character_formalODE_of_non_lisse (W : VertexOperatorAlgebra 𝒱)
    (P : Sl2Presentation W) (hex : P.IsExact)
    (N : OrdVModule W M) {hw : ℂ} (hcw : HasConformalWeight N hw) :
    ¬ C2Cofinite W ∧
      ∃ (s : ℕ) (g : ℕ → ℕ → ℂ), 1 ≤ s ∧ (∀ i, Subexponential (fun k => ‖g i k‖)) ∧
        ∀ e : ℂ, (qDeriv^[s] (torusPartition N)) e
          + ∑ i ∈ Finset.range s, qConv (g i) (qDeriv^[i] (torusPartition N)) e = 0 :=
\end{leancode}

\begin{dictionary}
Compositions of Theorems~\ref{l16:thm:presCons} and~\ref{l16:thm:formalODE}. For a general $V$ the hypotheses \lean{P},
\lean{hex} are genuine; the conclusion inherits the weakness discussed in Section~\ref{l16:sec:weak}.
\end{dictionary}

\textbf{Proof on paper.} Theorem~\ref{l16:thm:presCons} gives Etingof-finiteness (and non-lisseness from \lean{hex}); then
Theorem~\ref{l16:thm:formalODE}. \qed

\begin{physics}
Any chiral algebra whose classical phase space is a thickened nilpotent cone of $\lsixteensl$ has characters that are
convergent and solve a closed stress-tensor recursion.
\end{physics}

\begin{status}
Proved; \lean{VOA.QuasiLisseSl2.}\allowbreak\lean{character_}\allowbreak\lean{formalODE_}\allowbreak\lean{of_non_}\allowbreak\lean{lisse}: \lean{[propext, Classical.}\allowbreak\lean{choice, Quot.}\allowbreak\lean{sound]} (re-run).
\end{status}

\begin{theorem}[Unconditional consequences at $L_{-4/3}(\lsixteensl)$]\label{l16:thm:payoffL}
For $L = L_{-4/3}(\lsixteensl)$: (i) its character satisfies a formal equation~\eqref{l16:eq:formalODE}; (ii) so does the
character of every ordinary module with a conformal weight, (iii) with an order $s$ independent of the module; (iv) the
character $q$-series converges at every $\tau\in\mathbb H$; (v) the space of analytic genus-one blocks is finite-dimensional.
\end{theorem}
\noindent File \leanpath{VOA/\allowbreak Modularity/\allowbreak AKSl2Instance.lean}.
\begin{leancode}
theorem character_formalODE_simpleSl2QL_uncond :
    ∃ (s : ℕ) (g : ℕ → ℕ → ℂ), 1 ≤ s ∧ (∀ i, Subexponential (fun k => ‖g i k‖)) ∧
      ∀ e : ℂ, (qDeriv^[s] (torusPartition (adjointModule simpleSl2QL))) e
        + ∑ i ∈ Finset.range s,
            qConv (g i) (qDeriv^[i] (torusPartition (adjointModule simpleSl2QL))) e = 0 :=

theorem character_formalODE_module_simpleSl2QL {M : Type} [AddCommGroup M] [Module ℂ M]
    (N : OrdVModule simpleSl2QL M) {hw : ℂ} (hcw : HasConformalWeight N hw) :
    ∃ (s : ℕ) (g : ℕ → ℕ → ℂ), 1 ≤ s ∧ (∀ i, Subexponential (fun k => ‖g i k‖)) ∧
      ∀ e : ℂ, (qDeriv^[s] (torusPartition N)) e
        + ∑ i ∈ Finset.range s, qConv (g i) (qDeriv^[i] (torusPartition N)) e = 0 :=

theorem uniform_order_character_formalODE_simpleSl2QL :
    ∃ s : ℕ, 1 ≤ s ∧
      ∀ {M : Type} [AddCommGroup M] [Module ℂ M] (N : OrdVModule simpleSl2QL M) (hw : ℂ),
        HasConformalWeight N hw →
          ∃ g : ℕ → ℕ → ℂ, (∀ i, Subexponential (fun k => ‖g i k‖)) ∧
            ∀ e : ℂ, (qDeriv^[s] (torusPartition N)) e
              + ∑ i ∈ Finset.range s,
                  qConv (g i) (qDeriv^[i] (torusPartition N)) e = 0 :=

theorem qSummable_character_simpleSl2QL (τ : UpperHalfPlane) :
    QSummableAt (gradedDim (adjointModule simpleSl2QL))
      ((0 : ℂ) - simpleSl2QL.centralCharge / 24) τ :=

theorem finiteDimensional_analyticBlocks_simpleSl2QL :
    FiniteDimensional ℂ (Submodule.map restrictImGT (analyticBlocks simpleSl2QL)) :=
\end{leancode}

\begin{dictionary}
All in \lean{namespace VOA.Sl2Pres}. No hypotheses at all. \lean{QSummableAt χ h τ} is
\lean{Summable (fun n : ℕ => χ (h + n) * (qParam τ) ^ n)} (\leanpath{VOA/\allowbreak Modularity/\allowbreak ConvergenceBridge.lean}),
the convergence of $\sum_n\dim L_n\,q^n$ at $q = e^{2\pi i\tau}$; here $h = -c/24 = 1/4$. Module carriers \lean{M : Type} live in
universe $0$, matching \lean{simpleSl2QL}. (i)--(iii) are formal equations in the weak sense of Section~\ref{l16:sec:weak},
\textbf{not} MLDEs; the file header says so (``NOT AK Theorem 5.1, and NOT a modular linear differential equation''). (iv)
is genuine convergence. (v) is genuine (Theorem~\ref{l16:thm:blocks}). \textbf{What is not here:} an MLDE for the character
of $L_{-4/3}$ packaged as a statement about the character (the genuine MLDE for abstract analytic blocks is
Theorem~\ref{l16:thm:mlde-uncond}(ii), of unspecified order), and Zhu-type modular
invariance (Theorem~\ref{l16:thm:zhuEt} needs \lean{hfib}, \lean{hrank}).
\end{dictionary}

\textbf{Proof on paper.} Theorem~\ref{l16:thm:presL}(iv) gives Etingof-finiteness; Theorem~\ref{l16:thm:vgz} gives
\lean{hzero}; the vacuum of $L$ is nonzero (\lean{pbw}-free); now apply Theorems~\ref{l16:thm:formalODE},
\ref{l16:thm:uniform}, \ref{l16:thm:growth} and~\ref{l16:thm:blocks}. \qed

\begin{physics}
For the $(A_1,A_3)$ Argyres--Douglas theory: the Schur index converges for $|q|<1$ (the protected spectrum has subexponential
growth) and the torus conformal blocks of its chiral algebra form a finite-dimensional space. The MLDE that Beem--Rastelli
observe for this Schur index has order~2; the library proves only that the vacuum insertion of every analytic block satisfies
\emph{some} monic MLDE (Theorem~\ref{l16:thm:mlde-uncond}(ii)), not one of these formal equations, and not the order.
\end{physics}

\begin{status}
Proved, unconditional, \lean{pbw}-free: \lean{VOA.Sl2Pres.}\allowbreak\lean{character_}\allowbreak\lean{formalODE_}\allowbreak\lean{simpleSl2QL_}\allowbreak\lean{uncond},
\lean{uniform_}\allowbreak\lean{order_}\allowbreak\lean{character_}\allowbreak\lean{formalODE_}\allowbreak\lean{simpleSl2QL}, \lean{qSummable_}\allowbreak\lean{character_}\allowbreak\lean{simpleSl2QL},
\lean{finiteDimensional_}\allowbreak\lean{analyticBlocks_}\allowbreak\lean{simpleSl2QL}: \lean{[propext, Classical.}\allowbreak\lean{choice, Quot.}\allowbreak\lean{sound]} (re-run).
\end{status}

\begin{theorem}[The ``non-lisse'' conjunctions, as stated in the library]\label{l16:thm:nonlissePayoff}
$L_{-4/3}(\lsixteensl)$ is not $C_2$-cofinite and its character satisfies a formal equation~\eqref{l16:eq:formalODE}.
\end{theorem}
\noindent Files \leanpath{VOA/\allowbreak Modularity/\allowbreak AKSl2Instance.lean},
\leanpath{VOA/\allowbreak Modularity/\allowbreak AKSl2Payoff.lean}.
\begin{leancode}
theorem character_formalODE_of_non_lisse_simpleSl2QL :
    ¬ C2Cofinite simpleSl2QL ∧
      ∃ (s : ℕ) (g : ℕ → ℕ → ℂ), 1 ≤ s ∧ (∀ i, Subexponential (fun k => ‖g i k‖)) ∧
        ∀ e : ℂ, (qDeriv^[s] (torusPartition (adjointModule simpleSl2QL))) e
          + ∑ i ∈ Finset.range s,
              qConv (g i) (qDeriv^[i] (torusPartition (adjointModule simpleSl2QL))) e = 0 :=

theorem character_formalODE_simpleSl2QL :
    ∃ (s : ℕ) (g : ℕ → ℕ → ℂ), 1 ≤ s ∧ (∀ i, Subexponential (fun k => ‖g i k‖)) ∧
      ∀ e : ℂ, (qDeriv^[s] (torusPartition (adjointModule simpleSl2QL))) e
        + ∑ i ∈ Finset.range s,
            qConv (g i) (qDeriv^[i] (torusPartition (adjointModule simpleSl2QL))) e = 0 :=

theorem character_formalODE_adjoint_of_non_lisse_simpleSl2QL :
    ¬ C2Cofinite simpleSl2QL ∧
      ∃ (s : ℕ) (g : ℕ → ℕ → ℂ), 1 ≤ s ∧ (∀ i, Subexponential (fun k => ‖g i k‖)) ∧
        ∀ e : ℂ, (qDeriv^[s] (torusPartition (adjointModule simpleSl2QL))) e
          + ∑ i ∈ Finset.range s,
              qConv (g i) (qDeriv^[i] (torusPartition (adjointModule simpleSl2QL))) e = 0 :=
\end{leancode}

\begin{dictionary}
The first is \lean{VOA.Sl2Pres.}\allowbreak\lean{character_}\allowbreak\lean{formalODE_}\allowbreak\lean{of_non_}\allowbreak\lean{lisse_}\allowbreak\lean{simpleSl2QL} (\leanpath{AKSl2Instance.lean}); the other two are
in \lean{namespace VOA.}\allowbreak\lean{QuasiLisseSl2} (\leanpath{AKSl2Payoff.lean}). The hypotheses \lean{hex : Sl2ExactSymbol}, \lean{P},
\lean{hex : P.IsExact} were \textbf{all redundant}: the conclusions hold hypothesis-free (checked in Lean; the check is not part of this release). After this
finding the library removed them from these three flagships, whose current, hypothesis-free statements are quoted above.
\end{dictionary}

\textbf{Proof on paper.} Theorem~\ref{l16:thm:nonlisse}(iv) and Theorem~\ref{l16:thm:payoffL}(i). \qed

\begin{physics}
This is the configuration in which Arakawa--Kawasetsu say something Zhu cannot: a chiral algebra with a continuous Higgs
branch whose characters are nonetheless controlled. In the form proved, the control is convergence and a closed recursion,
not modularity.
\end{physics}

\begin{status}
Library: \lean{Sl2Pres.}\allowbreak\lean{character_}\allowbreak\lean{formalODE_}\allowbreak\lean{of_non_}\allowbreak\lean{lisse_}\allowbreak\lean{simpleSl2QL} standard three; the two \lean{QuasiLisseSl2} flagship
statements \lean{character_}\allowbreak\lean{formalODE_}\allowbreak\lean{simpleSl2QL} and \lean{character_}\allowbreak\lean{formalODE_}\allowbreak\lean{adjoint_}\allowbreak\lean{of_non_}\allowbreak\lean{lisse_}\allowbreak\lean{simpleSl2QL}:
\lean{[propext, Classical.}\allowbreak\lean{choice, Quot.}\allowbreak\lean{sound, VOA.}\allowbreak\lean{pbw]} at the snapshot;
now \lean{[propext, Classical.choice, Quot.sound]} (re-probed at a later commit, after \lean{VOA.pbw} became a theorem, not re-run for this edition; confirmed by the release audit,
2026-10-27). At the snapshot they used \lean{VOA.pbw} via \lean{vacuum_}\allowbreak\lean{simpleSl2QL_}\allowbreak\lean{ne_zero}; the
current source routes \lean{QuasiLisseSl2.}\allowbreak\lean{character_}\allowbreak\lean{formalODE_}\allowbreak\lean{simpleSl2QL} through the Wakimoto
\lean{Sl2Pres.}\allowbreak\lean{vacuum_}\allowbreak\lean{simpleSl2QL_}\allowbreak\lean{ne_zero'} instead. Either way, since \lean{VOA.pbw} is now a
theorem, ``\lean{pbw}-free'' is no longer a property that \texttt{\#print axioms} can detect or that carries any weight.
\textbf{Content:} these are the near-vacuous formal equations of Section~\ref{l16:sec:weak}, not MLDEs.
\end{status}

%======================================================================
\section{Traps, dead ends, and Lean defects}\label{l16:sec:defects}
%======================================================================

\begin{warning}[title={Primarity is essential in Lemma 5.5 (recorded dead end)}]
\leanpath{VOA/\allowbreak Modularity/\allowbreak UniCollapse.lean} refutes the unitriangular conversion
(Theorem~\ref{l16:thm:huni}) when quantified over \emph{all} vectors, via $v = L[-1]x$. An earlier note in
\leanpath{QuasiLisseFinite.lean} wrongly identified that refutation with AK Lemmas~5.4--5.6; the corrected note (Status item~1)
explains that only Lemma~5.5 needs primarity and that the refuted instance is not primary.
\end{warning}

\begin{warning}[title={Etingof-finiteness cannot supply Zhu's rank bound}]
\lean{EtingofFinite W} bounds $V/O_2^P(V)$, not Zhu's algebra $A(V)$, since $a_{(0)}b\notin O(V)$ in general. So the rank
hypothesis \lean{hrank} of Theorem~\ref{l16:thm:zhuEt} stays open on the Etingof route (\leanpath{AKZhuModularity.lean},
honesty item~2). Its satisfiability is witnessed only at the trivial VOA, which (the file notes) cannot distinguish it from a
refuted legacy hypothesis.
\end{warning}

\paragraph{Lean defects found while writing this chapter.}
\begin{enumerate}
\item \textbf{Moderate: redundant hypotheses.} \lean{hex : Sl2ExactSymbol} in \lean{VOA.Sl2Pres.}\allowbreak\lean{character_}\allowbreak\lean{formalODE_}\allowbreak\lean{of_non_}\allowbreak\lean{lisse_}\allowbreak\lean{simpleSl2QL},
  \lean{not_C2Cofinite_}\allowbreak\lean{simpleSl2QL_}\allowbreak\lean{of_exact}, \lean{quasiLisseWitness_}\allowbreak\lean{simpleSl2QL_}\allowbreak\lean{of_exact},
  \lean{nilpEquiv_}\allowbreak\lean{simpleSl2QL_}\allowbreak\lean{of_exact} (the first three; the fourth needs it genuinely), and \lean{S}, \lean{hex} in
  \lean{VOA.QuasiLisseSl2.}\allowbreak\lean{not_C2Cofinite_}\allowbreak\lean{simpleSl2QL}, \lean{quasiLisseWitness_}\allowbreak\lean{simpleSl2QL},
  \lean{character_}\allowbreak\lean{formalODE_}\allowbreak\lean{adjoint_}\allowbreak\lean{of_non_}\allowbreak\lean{lisse_}\allowbreak\lean{simpleSl2QL}, and \lean{S} or \lean{P} in
  \lean{QuasiLisseSl2.}\allowbreak\lean{etingofFinite_}\allowbreak\lean{simpleSl2QL}, \lean{finiteDimensional_}\allowbreak\lean{poissonTraces_}\allowbreak\lean{simpleSl2QL},
  \lean{character_}\allowbreak\lean{formalODE_}\allowbreak\lean{simpleSl2QL} (\leanpath{AffineSl2Presentation.lean},
  \leanpath{AffineSl2QuasiLisse.lean}, \leanpath{AKSl2Instance.lean}, \leanpath{AKSl2Payoff.lean}). Evidence: a Lean check against the library (not part of this
  release). Fix: restate with \lean{Sl2NonLisse.}\allowbreak\lean{not_C2Cofinite_}\allowbreak\lean{simpleSl2QL}.
\item \textbf{Moderate: near-vacuous conclusion.} The formal ODE theorems (\lean{character_}\allowbreak\lean{formalODE_}\allowbreak\lean{of_etingofFinite},
  \lean{uniform_}\allowbreak\lean{order_}\allowbreak\lean{character_}\allowbreak\lean{formalODE_}\allowbreak\lean{of_etingofFinite} and their instances) hold for any convergent $q$-series with zeros of
  bounded multiplicity in the disc (Proposition~\ref{l16:prop:weak}, on paper). Fix: state them as growth results, or use the
  $E_4,E_6$ route.
\item \textbf{Closed gap (formerly open): the genuine MLDE's side inputs.} \lean{AKWeightGrading} is now inhabited for every
  VOA (\lean{akWeightGrading}) and \lean{E46SerreClosed} is \lean{e46SerreClosed}, so
  \lean{IsSqPrimary.}\allowbreak\lean{exists_}\allowbreak\lean{monicMLDEOn_}\allowbreak\lean{all_analyticBlocks_}\allowbreak\lean{of_etingofFinite_uncond}
  is conditional only on \lean{EtingofFinite} (Theorem~\ref{l16:thm:mlde-uncond}). \textbf{Documentation only:} the
  suffix \lean{_uncond} is misleading, since the theorem keeps \lean{hE}.
\item \textbf{Documentation only: stale docstrings} saying $L_{-4/3}(\lsixteensl)$ is not proved non-lisse / only the trivial VOA is
  Etingof-finite / the presentation is never constructed: \leanpath{QuasiLisseFinite.lean} (Status item~4, \S7),
  \leanpath{AKModularODE.lean} (header and \lean{trivial_}\allowbreak\lean{character_}\allowbreak\lean{formalODE}), \leanpath{AKUniformODE.lean}
  (\lean{trivial_}\allowbreak\lean{uniform_}\allowbreak\lean{order_}\allowbreak\lean{character_}\allowbreak\lean{formalODE}), \leanpath{AKZhuModularity.lean} (honesty item~1, headline docstring),
  \leanpath{AffineSl2QuasiLisse.lean}, \leanpath{AffineSl2Presentation.lean}, \leanpath{AKSl2Payoff.lean},
  \leanpath{AKSl2Instance.lean}.
\item \textbf{Documentation only: name clash} (\emph{partly resolved}). \lean{VOA.QuasiLisseSl2.}\allowbreak\lean{not_C2Cofinite_}\allowbreak\lean{simpleSl2QL} (needed \lean{S}, \lean{hex} at the snapshot) and
  \lean{VOA.Sl2NonLisse.}\allowbreak\lean{not_C2Cofinite_}\allowbreak\lean{simpleSl2QL} (unconditional); similarly \lean{QuasiLisseSl2.}\allowbreak\lean{etingofFinite_}\allowbreak\lean{simpleSl2QL}
  (needed \lean{S}) versus \lean{Sl2Pres.}\allowbreak\lean{etingofFinite_}\allowbreak\lean{simpleSl2QL'}. At the release audit (2026-10-27) the
  \lean{QuasiLisseSl2} twins in \leanpath{AKSl2Payoff.lean} are hypothesis-free and have the same types as their
  \lean{Sl2NonLisse}/\lean{Sl2Pres} counterparts; the duplicate names remain.
\item \textbf{Documentation only: imprecise citation.} \lean{Sl2ExactSymbol} is described as Arakawa's associated-variety theorem, which only
  gives the radical of the kernel (Definition~\ref{l16:def:exact} and its warning).
\item \textbf{Documentation only: unnecessary \lean{pbw}} (\emph{moot}). At the snapshot \lean{QuasiLisseSl2.}\allowbreak\lean{A₂_simpleSl2QL_}\allowbreak\lean{nontrivial},
  \lean{QuasiLisseSl2.}\allowbreak\lean{character_}\allowbreak\lean{formalODE_}\allowbreak\lean{simpleSl2QL}, \lean{character_}\allowbreak\lean{formalODE_}\allowbreak\lean{adjoint_}\allowbreak\lean{of_non_}\allowbreak\lean{lisse_}\allowbreak\lean{simpleSl2QL} depended on
  the then-axiom \lean{VOA.pbw} although \lean{pbw}-free twins existed. The two formal-ODE statements are now routed
  through the Wakimoto vacuum, and \lean{VOA.pbw} is a theorem (\leanpath{VOA/\allowbreak Foundations/\allowbreak AffineLie.lean}), so
  every declaration of this chapter reports Lean's three standard axioms and ``without \lean{VOA.pbw}'' is not a
  meaningful selling point.
\item \textbf{Documentation only: overclaiming headline docstring} (release audit, finding~3). The docstring of
  \lean{Sl2Pres.}\allowbreak\lean{character_}\allowbreak\lean{formalODE_}\allowbreak\lean{simpleSl2QL_}\allowbreak\lean{uncond} (\leanpath{AKSl2Instance.lean})
  calls the near-vacuous formal equation ``THE HEADLINE'' and ``the Arakawa--Kawasetsu character equation'', contradicting
  the disclaimers elsewhere in the same files. It should say ``formal growth equation''.
\item \textbf{Minor: redundant \lean{hzero : VacuumGradeZero W}} (release audit, finding~6). It holds for every VOA
  (\lean{vacuumGradeZero_holds}, \leanpath{VOA/\allowbreak Modularity/\allowbreak BlockZeroMode.lean}) but is still an argument of
  \lean{IsSqPrimary.}\allowbreak\lean{exists_}\allowbreak\lean{monicMLDEOn_}\allowbreak\lean{all_analyticBlocks_}\allowbreak\lean{of_etingofFinite},
  \lean{character_}\allowbreak\lean{formalODE_}\allowbreak\lean{of_etingofFinite}, \lean{onePointSubexp_}\allowbreak\lean{of_etingofFinite} and
  \lean{finiteDimensional_}\allowbreak\lean{analyticBlocks_}\allowbreak\lean{of_etingofFinite}. Harmless.
\end{enumerate}

%======================================================================
\section{Referee's checklist}\label{l16:sec:checklist}
%======================================================================

\begingroup\small
\begin{longtable}{p{0.30\textwidth}p{0.27\textwidth}p{0.24\textwidth}p{0.11\textwidth}}
\toprule
Lean declaration & Intended meaning & Hypotheses / status & Fidelity check\\
\midrule\endhead
\lean{A₂Algebra}, \lean{QuotVOA.A₂quot} & $R_V$ is a $\C$-algebra; $R_V\twoheadrightarrow R_{V/I}$ Poisson & graded ideal; unconditional & Def.~\ref{l16:def:algebra}, Thm.~\ref{l16:thm:quot}\\
\lean{poissonSpan}, \lean{HP₀}, \lean{poissonTraces} & $\{R,R\}$, $\lsixteenHP(R_V)$, Poisson traces & definitions; standard & Def.~\ref{l16:def:poissonSpan}\\
\lean{EtingofFinite} & $\dim\lsixteenHP(R_V)<\infty$ & definition; \textbf{not} quasi-lisseness (weaker, by Etingof--Schedler) & Def.~\ref{l16:def:etingof}\\
\lean{EtingofFinite_}\allowbreak\lean{of_C2Cofinite}, \lean{finiteDimensional_}\allowbreak\lean{poissonTraces} & lisse $\Rightarrow$ Etingof-finite; finite trace space & unconditional & Thm.~\ref{l16:thm:c2etingof}\\
\lean{etingofFinite_}\allowbreak\lean{iff_quot_}\allowbreak\lean{O₂P} & AK Thm.~2.2 upstairs & unconditional & Thm.~\ref{l16:thm:O2P}\\
\lean{vSqZ_zero_eq}, \lean{sqDefect0_mem_of_lower} & $a[0]b = a_{(0)}b+$lower & unconditional & Thm.~\ref{l16:thm:triang}\\
\lean{exists_finite_gen_sup_}\allowbreak\lean{blockRelSpanG_}\allowbreak\lean{of_etingofFinite} & AK Prop.~5.2 & \lean{hE}; \lean{hzero} free & Thm.~\ref{l16:thm:prop52}\\
\lean{onePointSubexp_}\allowbreak\lean{of_etingofFinite} & one-point functions converge & \lean{hcw}, \lean{hE} & Thm.~\ref{l16:thm:growth}\\
\lean{IsSqPrimary}, \lean{Lsq_tower} & primaries; Virasoro tower & unconditional & Def.~\ref{l16:def:sqprimary}, Thm.~\ref{l16:thm:tower}\\
\lean{huni_of_sqPrimary} & AK Lemma 5.5 ($\theta$-basis, $\mathcal S$-coefficients, block-dependent) & primary $a$ & Thm.~\ref{l16:thm:huni}\\
\lean{exists_}\allowbreak\lean{monic_}\allowbreak\lean{Lsq_relation_}\allowbreak\lean{of_etingofFinite} & AK Lemma 5.3 (block form) & \lean{hE} & Thm.~\ref{l16:thm:lemma53}\\
\lean{character_}\allowbreak\lean{formalODE_}\allowbreak\lean{of_etingofFinite} & ``AK eq.~(2)'' formal; \textbf{not an MLDE, not AK Thm.~5.1; near-vacuous} (Prop.~\ref{l16:prop:weak}) & \lean{hcw}, \lean{hE} & Thm.~\ref{l16:thm:formalODE}, Prop.~\ref{l16:prop:weak}\\
\lean{exists_}\allowbreak\lean{uniform_}\allowbreak\lean{monic_}\allowbreak\lean{Lsq_relation_}\allowbreak\lean{of_etingofFinite}, \lean{uniform_}\allowbreak\lean{order_}\allowbreak\lean{character_}\allowbreak\lean{formalODE_}\allowbreak\lean{of_etingofFinite} & uniform Lemma 5.3; uniform order only & \lean{hE} & Thm.~\ref{l16:thm:uniform}\\
\lean{MonicMLDEOn}, \lean{AKWeightGrading}, \lean{E46SerreClosed} & genuine MLDE; grading data; Ramanujan closure & definitions; \lean{AKWeightGrading} inhabited for every VOA (\lean{akWeightGrading}); \lean{E46SerreClosed} proved (\lean{e46SerreClosed}) & Defs.~\ref{l16:def:mlde}, \ref{l16:def:akgrading}, Thm.~\ref{l16:thm:akgrading}\\
\lean{IsSqPrimary.}\allowbreak\lean{exists_}\allowbreak\lean{monicMLDEOn_}\allowbreak\lean{all_analyticBlocks_}\allowbreak\lean{of_etingofFinite} & AK Thm.~5.1 (one MLDE for all blocks) & \lean{hE}, \lean{hgrade}, \lean{hclose}; \lean{hclose} discharged in the primed form & Thm.~\ref{l16:thm:mlde}\\
\lean{IsSqPrimary.}\allowbreak\lean{exists_}\allowbreak\lean{monicMLDEOn_}\allowbreak\lean{all_analyticBlocks_}\allowbreak\lean{of_etingofFinite_uncond} & AK Thm.~5.1, side inputs discharged & \lean{hE} only (despite the name); order unbounded & Thm.~\ref{l16:thm:mlde-uncond}\\
\lean{exists_}\allowbreak\lean{monicMLDEOn_}\allowbreak\lean{all_analyticBlocks_}\allowbreak\lean{simpleSl2QL} & one monic MLDE for all blocks of $L_{-4/3}$ at the vacuum & unconditional; order unspecified (AK's order~2 not proved) & Thm.~\ref{l16:thm:mlde-uncond}\\
\lean{finiteDimensional_}\allowbreak\lean{analyticBlocks_}\allowbreak\lean{of_etingofFinite} & Zhu 5.3.2 under \lean{hE} & \lean{hE} & Thm.~\ref{l16:thm:blocks}\\
\lean{exists_}\allowbreak\lean{zhuModularity_}\allowbreak\lean{of_}\allowbreak\lean{sqAnalyticRank_}\allowbreak\lean{le_irr_}\allowbreak\lean{etingofFinite} & Zhu modularity under \lean{hE} & \lean{hE}, \lean{hfib}, \lean{hrank} (discharged only for the trivial VOA; \lean{hrank} expected \textbf{false} at $L_{-4/3}$, paper argument) & Thm.~\ref{l16:thm:zhuEt}\\
\lean{QuasiLisseSl2.}\allowbreak\lean{bigSub_}\allowbreak\lean{eq_top}, \lean{not_finiteDimensional_}\allowbreak\lean{NilpAlg}, \lean{finiteDimensional_}\allowbreak\lean{quot_adRange} & classical model: infinite $\lsixteenR$, finite $\lsixteenHP$ & unconditional & Thm.~\ref{l16:thm:model}\\
\lean{Sl2Presentation}, \lean{.etingofFinite}, \lean{.not_C2Cofinite} & presentations and consequences & presentation (+ exactness) & Def.~\ref{l16:def:pres}, Thm.~\ref{l16:thm:presCons}\\
\lean{affineSl2QL}, \lean{simpleSl2QL} & $V^{-4/3}(\lsixteensl)$, $L_{-4/3}(\lsixteensl)$ & definitions; $=$ \lean{L43} of L09 & Def.~\ref{l16:def:flagship}\\
\lean{Sl2Pres.}\allowbreak\lean{sl2PresentationSimpleSl2QL}, \lean{etingofFinite_}\allowbreak\lean{simpleSl2QL'} & presentation constructed; $L_{-4/3}$ Etingof-finite & \textbf{unconditional}, \lean{pbw}-free & Thm.~\ref{l16:thm:presL}\\
\lean{Sl2ExactSymbol} & exactness of the presentation & \textbf{unproved}; now redundant; citation imprecise & Def.~\ref{l16:def:exact}\\
\lean{Sl2NonLisse.}\allowbreak\lean{not_C2Cofinite_}\allowbreak\lean{simpleSl2QL}, \lean{strictEtingofFinite_}\allowbreak\lean{simpleSl2QL} & $L_{-4/3}$ not lisse, and Etingof-finite & \textbf{unconditional}, \lean{pbw}-free & Thm.~\ref{l16:thm:nonlisse}\\
\lean{A₂_simpleSl2QL_}\allowbreak\lean{nontrivial} (and \lean{'}) & $R_{L}\neq0$ & unconditional; standard axioms (unprimed passes through the theorem \lean{VOA.pbw}) & Thm.~\ref{l16:thm:nontriv}\\
\lean{Sl2Pres.}\allowbreak\lean{character_}\allowbreak\lean{formalODE_}\allowbreak\lean{simpleSl2QL_}\allowbreak\lean{uncond} and siblings & formal ODE, convergence, finite blocks at $L_{-4/3}$ & unconditional; formal ODE \textbf{near-vacuous}, not an MLDE & Thm.~\ref{l16:thm:payoffL}\\
\lean{character_}\allowbreak\lean{formalODE_}\allowbreak\lean{of_non_}\allowbreak\lean{lisse_}\allowbreak\lean{simpleSl2QL} and \lean{QuasiLisseSl2} flagships & not lisse $\wedge$ formal ODE (\textbf{not an MLDE}) & none (\lean{hex}/\lean{P} removed); standard axioms & Thm.~\ref{l16:thm:nonlissePayoff}\\
\bottomrule
\end{longtable}
\endgroup

%======================================================================
\section{Notes and references}\label{l16:sec:notes}
%======================================================================

\textbf{The main theorem.} T.~Arakawa and K.~Kawasetsu, ``Quasi-lisse vertex algebras and modular linear differential
equations'', in \emph{Lie Groups, Geometry, and Representation Theory} (V.~G.~Kac and V.~L.~Popov, eds.), Progr.\ Math.\ 326,
Birkh\"auser (2018), 41--57, arXiv:1610.05865. The numbering used in this chapter (Theorem~2.2, Proposition~5.2,
Lemmas~5.3--5.6, Theorems~5.1 and~5.7, equation~(2)) follows arXiv version~4, which is the version cited in the source.

\textbf{Associated varieties.} T.~Arakawa, ``A remark on the $C_2$-cofiniteness condition on vertex algebras'', \emph{Math.\
Z.} 270 (2012), 559--575, arXiv:1004.1492 (definition of the associated variety; lisse iff $X_V$ is a point); T.~Arakawa,
``Associated varieties of modules over Kac--Moody algebras and $C_2$-cofiniteness of W-algebras'', \emph{Int.\ Math.\ Res.\
Not.} 2015, 11605--11666, arXiv:1004.1554 ($X_{L_k(\mathfrak g)}$ at admissible level is a nilpotent orbit closure; for
$\lsixteensl$, the nilpotent cone).

\textbf{Poisson traces.} P.~Etingof and T.~Schedler, ``Poisson traces and $D$-modules on Poisson varieties'',
\emph{Geom.\ Funct.\ Anal.} 20 (2010), 958--987, arXiv:0908.3868 (finitely many symplectic leaves $\Rightarrow$
$\dim\lsixteenHP<\infty$).

\textbf{Zhu's theorem.} Y.~Zhu, ``Modular invariance of characters of vertex operator algebras'', \emph{J.\ Amer.\ Math.\
Soc.} 9 (1996), 237--302 (relation module, Lemmas~4.4.1--4.4.2, Theorem~5.3.2). C.~Dong, H.~Li and G.~Mason,
``Modular-invariance of trace functions in orbifold theory and generalized Moonshine'', \emph{Comm.\ Math.\ Phys.} 214
(2000), 1--56 (Lemma~5.2, the $\mathcal G$-module form).

\textbf{Admissible $\widehat{\lsixteensl}$.} V.~G.~Kac and M.~Wakimoto, ``Modular invariant representations of
infinite-dimensional Lie algebras and superalgebras'', \emph{Proc.\ Nat.\ Acad.\ Sci.\ USA} 85 (1988), 4956--4960;
F.~G.~Malikov, B.~L.~Feigin and D.~B.~Fuchs, ``Singular vectors in Verma modules over Kac--Moody algebras'', \emph{Funct.\
Anal.\ Appl.} 20 (1986), 103--113; D.~Adamovi\'c and A.~Milas, ``Vertex operator algebras associated to modular invariant
representations for $A_1^{(1)}$'', \emph{Math.\ Res.\ Lett.} 2 (1995), 563--575 (classification of $L_k(\lsixteensl)$-modules
at admissible level).

\textbf{Ideals of holomorphic functions} (used in Proposition~\ref{l16:prop:weak}). J.~H.~M.~Wedderburn, ``On matrices whose
coefficients are functions of a single variable'', \emph{Trans.\ Amer.\ Math.\ Soc.} 16 (1915), 328--332; O.~Helmer,
``Divisibility properties of integral functions'', \emph{Duke Math.\ J.} 6 (1940), 345--356.

\textbf{Physics.} C.~Beem, M.~Lemos, P.~Liendo, W.~Peelaers, L.~Rastelli and B.~C.~van~Rees, ``Infinite chiral symmetry in
four dimensions'', \emph{Comm.\ Math.\ Phys.} 336 (2015), 1359--1433, arXiv:1312.5344 (the chiral algebra of a 4d
$\mathcal N=2$ SCFT, $c_{2d} = -12c_{4d}$, Schur index as vacuum character). C.~Beem and L.~Rastelli, ``Vertex operator
algebras, Higgs branches, and modular differential equations'', \emph{J.\ High Energy Phys.} 2018(8), 114,
arXiv:1707.07679 (Higgs branch $=$ associated variety, quasi-lisseness, MLDEs for Schur indices). T.~Creutzig, ``W-algebras for
Argyres--Douglas theories'', \emph{Eur.\ J.\ Math.} 3 (2017), 659--690, arXiv:1701.05926, and C.~Cordova and S.-H.~Shao,
``Schur indices, BPS particles, and Argyres--Douglas theories'', \emph{J.\ High Energy Phys.} 2016(1), 040, arXiv:1506.00265
(the $(A_1,D_{2n+1})$ theories and $L_{-4n/(2n+1)}(\lsixteensl)$). S.~D.~Mathur, S.~Mukhi and A.~Sen, ``On the classification of
rational conformal field theories'', \emph{Phys.\ Lett.\ B} 213 (1988), 303--308 (MLDEs in CFT).

\paragraph{Verification record.} Statements re-checked against the source for this chapter at the earlier commits indicated in its
status boxes (the snapshot and a later re-probe); they were not re-run for this edition.
Release audit 2026-10-27: see Notes/Textbook/referee/RELEASE\_AUDIT\_L16.md.

\textbf{Project records.} The honesty blocks in the module docstrings of the thirteen files of this chapter; the axiom audit
\texttt{scripts/\allowbreak audit/\allowbreak AKZhuModularityAxiomAudit.lean}.

\endgroup
